% !TEX program = lualatex
% ---------------------------------------------------------------------------------------------
%  Deformalization 1: twisted cokernels in a cyclotomic tower   (riemann-rev)
%
%  Build:  python3 -B scripts/build_digest.py    (writes data.tex and figure.tex, runs LuaLaTeX, checks)
%
%  One cluster of the Lean archive behind L. Alpöge's announced proof of RH, rewritten as ordinary
%  mathematics. The prose is hand-written. Lean names, line ranges, counts and the numbers of the
%  examples come from data.tex; the quoted Lean statements are copied from the archive by the script.
%  Each statement was compared with the Lean source at the pinned commit (Appendix A).
%
%  Rules of the text: standard terminology; every symbol is explained where it first appears and
%  listed again in the notation at the end; no letter has two meanings. Letters differ from the
%  archive's where that helps: its twist e is h here (as in the guide), its centres s_i are c_i,
%  its level c is a, its reference layer jZero l R is m_0(l,R).
%
%  The signature rule: a function that the archive defines is written with every explicit argument of its
%  Lean definition, in every place, through a macro of data.tex (\tw, \mz, \vl, …). No argument is fixed and dropped.
%  Every hypothesis of a theorem stands in its statement or in a numbered Setting that the statement cites.
%
%  Layout: the reading edition of the guide (11 pt Libertinus on a 140 mm measure).
% ---------------------------------------------------------------------------------------------
\documentclass[11pt,a4paper]{article}
\ifdefined\pdfvariable\pdfvariable suppressoptionalinfo \numexpr 1+2+4+8\relax\fi
\usepackage[a4paper,hmargin=35mm,top=22mm,bottom=27mm,footskip=12mm,nohead]{geometry}

\usepackage{fontspec}
\directlua{
  luaotfload.add_fallback("gfallback", {"STIXTwoMath-Regular.otf:mode=node;", "DejaVuSans.ttf:mode=node;"})
  luaotfload.add_fallback("gmonofallback", {"DejaVuSansMono.ttf:mode=node;", "STIXTwoMath-Regular.otf:mode=node;"})
}
\setmainfont{LibertinusSerif}[Extension=.otf,UprightFont=*-Regular,ItalicFont=*-Italic,
  BoldFont=*-Semibold,BoldItalicFont=*-SemiboldItalic,RawFeature={fallback=gfallback}]
\setsansfont{Inter}[Extension=.otf,UprightFont=*-Regular,ItalicFont=*-Italic,
  BoldFont=*-SemiBold,BoldItalicFont=*-SemiBoldItalic,RawFeature={fallback=gfallback}]
\setmonofont{JetBrainsMono}[Extension=.otf,UprightFont=*-Regular,ItalicFont=*-Italic,
  BoldFont=*-SemiBold,BoldItalicFont=*-SemiBoldItalic,Scale=MatchLowercase,
  RawFeature={-calt;-liga;fallback=gmonofallback}]
\usepackage{amsmath}
\usepackage{amsthm}
\usepackage[math-style=TeX]{unicode-math}
\setmathfont{LibertinusMath-Regular.otf}

\usepackage{xcolor}
\usepackage{tikz}
\usetikzlibrary{calc,arrows.meta,positioning,backgrounds}
\usepackage[protrusion=true,expansion=true]{microtype}
\usepackage{enumitem}
\usepackage{needspace}
\usepackage{booktabs}
\usepackage{array}
\usepackage{longtable}
\usepackage{fancyhdr}
\usepackage{lastpage}
\usepackage{multicol}
\usepackage[hidelinks,pdfusetitle]{hyperref}

\definecolor{ink}{HTML}{0F172A}
\definecolor{inksoft}{HTML}{475569}
\definecolor{inkfaint}{HTML}{94A3B8}
\definecolor{rulegrey}{HTML}{CBD5E1}
\definecolor{arith}{HTML}{C2410C}
\definecolor{est}{HTML}{7C3AED}
\definecolor{ana}{HTML}{0891B2}
\definecolor{ok}{HTML}{047857}
\definecolor{leancol}{HTML}{3730A3}

\title{Twisted cokernels in a cyclotomic tower}
\author{riemann-rev}
\date{}

% ------------------------------------------------------------------------------------- macros
\newcommand{\leanfont}{\ttfamily\fontsize{8.6}{10.5}\selectfont\color{leancol}}
\newcommand{\lean}[1]{\csname LN@#1\endcsname}
\newcommand{\leanquote}[1]{\csname LQ@#1\endcsname}
% generated by scripts/build_digest.py -- do not edit
\expandafter\newcommand\csname LN@GenBy\endcsname{\href{https://github.com/Deicyde/riemann-hypothesis/blob/40e56bab026552e089d7afcfeff08d2defb9c005/src/Development/AutomorphicArithmetic/AutomorphicVarieties/AdditiveValuationNorms/CokLenLeOfRangeLe.lean\#L9-L10}{\leanfont GenBy}}
\expandafter\newcommand\csname LN@jZero\endcsname{\href{https://github.com/Deicyde/riemann-hypothesis/blob/40e56bab026552e089d7afcfeff08d2defb9c005/src/Development/AutomorphicArithmetic/AutomorphicVarieties/AdditiveValuationNorms/CokLenLeOfRangeLe.lean\#L12-L12}{\leanfont jZero}}
\expandafter\newcommand\csname LN@IsLevel\endcsname{\href{https://github.com/Deicyde/riemann-hypothesis/blob/40e56bab026552e089d7afcfeff08d2defb9c005/src/Development/AutomorphicArithmetic/AutomorphicVarieties/AdditiveValuationNorms/CokLenLeOfRangeLe.lean\#L16-L16}{\leanfont IsLevel}}
\expandafter\newcommand\csname LN@cokLen\endcsname{\href{https://github.com/Deicyde/riemann-hypothesis/blob/40e56bab026552e089d7afcfeff08d2defb9c005/src/Development/AutomorphicArithmetic/AutomorphicVarieties/AdditiveValuationNorms/CokLenLeOfRangeLe.lean\#L20-L20}{\leanfont cokLen}}
\expandafter\newcommand\csname LN@InPow\endcsname{\href{https://github.com/Deicyde/riemann-hypothesis/blob/40e56bab026552e089d7afcfeff08d2defb9c005/src/Development/AutomorphicArithmetic/AutomorphicVarieties/AdditiveValuationNorms/CokLenLeOfRangeLe.lean\#L22-L22}{\leanfont InPow}}
\expandafter\newcommand\csname LN@Unip\endcsname{\href{https://github.com/Deicyde/riemann-hypothesis/blob/40e56bab026552e089d7afcfeff08d2defb9c005/src/Development/AutomorphicArithmetic/AutomorphicVarieties/AdditiveValuationNorms/CokLenLeOfRangeLe.lean\#L24-L24}{\leanfont Unip}}
\expandafter\newcommand\csname LN@twLen\endcsname{\href{https://github.com/Deicyde/riemann-hypothesis/blob/40e56bab026552e089d7afcfeff08d2defb9c005/src/Development/AutomorphicArithmetic/AutomorphicVarieties/AdditiveValuationNorms/CokLenLeOfRangeLe.lean\#L26-L26}{\leanfont twLen}}
\expandafter\newcommand\csname LN@range\endcsname{\href{https://github.com/Deicyde/riemann-hypothesis/blob/40e56bab026552e089d7afcfeff08d2defb9c005/src/Development/AutomorphicArithmetic/AutomorphicVarieties/AdditiveValuationNorms/CokLenLeOfRangeLe.lean\#L28-L30}{\leanfont cokLen\_le\_of\_range\_le}}
\expandafter\newcommand\csname LN@lenComp\endcsname{\href{https://github.com/Deicyde/riemann-hypothesis/blob/40e56bab026552e089d7afcfeff08d2defb9c005/src/Development/AutomorphicArithmetic/AutomorphicVarieties/GaloisLibrary/SelmerTransport/LenComp.lean\#L17-L41}{\leanfont Base.lenComp}}
\expandafter\newcommand\csname LN@smith\endcsname{\href{https://github.com/Deicyde/riemann-hypothesis/blob/40e56bab026552e089d7afcfeff08d2defb9c005/src/Development/AutomorphicArithmetic/CyclotomicBounds/ClassNumberBounds/SelmerDuality/LengthDet.lean\#L52-L68}{\leanfont length\_cokernel\_eq\_addVal\_det}}
\expandafter\newcommand\csname LN@lift\endcsname{\href{https://github.com/Deicyde/riemann-hypothesis/blob/40e56bab026552e089d7afcfeff08d2defb9c005/src/Development/AutomorphicArithmetic/AutomorphicVarieties/AdditiveValuationNorms/ImportsAlt.lean\#L7-L44}{\leanfont lift}}
\expandafter\newcommand\csname LN@lte\endcsname{\href{https://github.com/Deicyde/riemann-hypothesis/blob/40e56bab026552e089d7afcfeff08d2defb9c005/src/Development/AutomorphicArithmetic/AutomorphicVarieties/GaloisLibrary/LocalRingsCompletion/IsUnitAddMem.lean\#L109-L126}{\leanfont abelian₂.lte}}
\expandafter\newcommand\csname LN@lteNorm\endcsname{\href{https://github.com/Deicyde/riemann-hypothesis/blob/40e56bab026552e089d7afcfeff08d2defb9c005/src/Development/AutomorphicArithmetic/AutomorphicVarieties/AdditiveValuationNorms/LteNorm.lean\#L19-L94}{\leanfont lteNorm}}
\expandafter\newcommand\csname LN@units\endcsname{\href{https://github.com/Deicyde/riemann-hypothesis/blob/40e56bab026552e089d7afcfeff08d2defb9c005/src/Development/AutomorphicArithmetic/AutomorphicVarieties/AdditiveValuationNorms/PadicUnitsDistanceLeMax.lean\#L8-L35}{\leanfont padic\_units\_distance\_le\_max}}
\expandafter\newcommand\csname LN@helly\endcsname{\href{https://github.com/Deicyde/riemann-hypothesis/blob/40e56bab026552e089d7afcfeff08d2defb9c005/src/Development/AutomorphicArithmetic/AutomorphicVarieties/AdditiveValuationNorms/Helly.lean\#L6-L25}{\leanfont helly}}
\expandafter\newcommand\csname LN@ball\endcsname{\href{https://github.com/Deicyde/riemann-hypothesis/blob/40e56bab026552e089d7afcfeff08d2defb9c005/src/Development/AutomorphicArithmetic/AutomorphicVarieties/AdditiveValuationNorms/PadicUnitsDistanceLeMax.lean\#L37-L63}{\leanfont ball}}
\expandafter\newcommand\csname LN@detFactor\endcsname{\href{https://github.com/Deicyde/riemann-hypothesis/blob/40e56bab026552e089d7afcfeff08d2defb9c005/src/Development/AutomorphicArithmetic/AutomorphicVarieties/AdditiveValuationNorms/DetFactor.lean\#L6-L49}{\leanfont detFactor}}
\expandafter\newcommand\csname LN@normVal\endcsname{\href{https://github.com/Deicyde/riemann-hypothesis/blob/40e56bab026552e089d7afcfeff08d2defb9c005/src/Development/AutomorphicArithmetic/AutomorphicVarieties/AdditiveValuationNorms/NormVal.lean\#L30-L46}{\leanfont normVal}}
\expandafter\newcommand\csname LN@detBound\endcsname{\href{https://github.com/Deicyde/riemann-hypothesis/blob/40e56bab026552e089d7afcfeff08d2defb9c005/src/Development/AutomorphicArithmetic/AutomorphicVarieties/AdditiveValuationNorms/DetBound.lean\#L9-L34}{\leanfont detBound}}
\expandafter\newcommand\csname LN@level\endcsname{\href{https://github.com/Deicyde/riemann-hypothesis/blob/40e56bab026552e089d7afcfeff08d2defb9c005/src/Development/AutomorphicArithmetic/AutomorphicVarieties/AdditiveValuationNorms/ExistsCentresLevel.lean\#L23-L32}{\leanfont exists\_centres\_level}}
\expandafter\newcommand\csname LN@cokL\endcsname{\href{https://github.com/Deicyde/riemann-hypothesis/blob/40e56bab026552e089d7afcfeff08d2defb9c005/src/Development/AutomorphicArithmetic/AutomorphicVarieties/AdditiveValuationNorms/CokLen.lean\#L6-L41}{\leanfont cokL}}
\expandafter\newcommand\csname LN@nilLen\endcsname{\href{https://github.com/Deicyde/riemann-hypothesis/blob/40e56bab026552e089d7afcfeff08d2defb9c005/src/Development/AutomorphicArithmetic/AutomorphicVarieties/AdditiveValuationNorms/NilLen.lean\#L29-L59}{\leanfont nilLen}}
\expandafter\newcommand\csname LN@unipPow\endcsname{\href{https://github.com/Deicyde/riemann-hypothesis/blob/40e56bab026552e089d7afcfeff08d2defb9c005/src/Development/AutomorphicArithmetic/AutomorphicVarieties/AdditiveValuationNorms/UnipPow.lean\#L8-L46}{\leanfont unipPow}}
\expandafter\newcommand\csname LN@frob\endcsname{\href{https://github.com/Deicyde/riemann-hypothesis/blob/40e56bab026552e089d7afcfeff08d2defb9c005/src/Development/AutomorphicArithmetic/AutomorphicVarieties/AdditiveValuationNorms/Imports.lean\#L6-L36}{\leanfont frob}}
\expandafter\newcommand\csname LN@powIn\endcsname{\href{https://github.com/Deicyde/riemann-hypothesis/blob/40e56bab026552e089d7afcfeff08d2defb9c005/src/Development/AutomorphicArithmetic/AutomorphicVarieties/AdditiveValuationNorms/PowIn.lean\#L6-L29}{\leanfont powIn}}
\expandafter\newcommand\csname LN@phi\endcsname{\href{https://github.com/Deicyde/riemann-hypothesis/blob/40e56bab026552e089d7afcfeff08d2defb9c005/src/Development/AutomorphicArithmetic/AutomorphicVarieties/AdditiveValuationNorms/ImportsTer.lean\#L7-L46}{\leanfont phi}}
\expandafter\newcommand\csname LN@step\endcsname{\href{https://github.com/Deicyde/riemann-hypothesis/blob/40e56bab026552e089d7afcfeff08d2defb9c005/src/Development/AutomorphicArithmetic/AutomorphicVarieties/AdditiveValuationNorms/ImportsMore.lean\#L9-L22}{\leanfont step}}
\expandafter\newcommand\csname LN@mono\endcsname{\href{https://github.com/Deicyde/riemann-hypothesis/blob/40e56bab026552e089d7afcfeff08d2defb9c005/src/Development/AutomorphicArithmetic/AutomorphicVarieties/AdditiveValuationNorms/ImportsBis.lean\#L6-L13}{\leanfont mono}}
\expandafter\newcommand\csname LN@twist\endcsname{\href{https://github.com/Deicyde/riemann-hypothesis/blob/40e56bab026552e089d7afcfeff08d2defb9c005/src/Development/AutomorphicArithmetic/AutomorphicVarieties/AdditiveValuationNorms/ImportsFurther.lean\#L7-L50}{\leanfont twist}}
\expandafter\newcommand\csname LN@growth\endcsname{\href{https://github.com/Deicyde/riemann-hypothesis/blob/40e56bab026552e089d7afcfeff08d2defb9c005/src/Development/AutomorphicArithmetic/AutomorphicVarieties/AdditiveValuationNorms/Growth.lean\#L11-L56}{\leanfont growth}}
\expandafter\newcommand\csname LN@twgrowth\endcsname{\href{https://github.com/Deicyde/riemann-hypothesis/blob/40e56bab026552e089d7afcfeff08d2defb9c005/src/Development/AutomorphicArithmetic/AutomorphicVarieties/AdditiveValuationNorms/ExistsCentresLevel.lean\#L17-L21}{\leanfont twLen\_le\_growth}}
\expandafter\newcommand\csname LN@centres\endcsname{\href{https://github.com/Deicyde/riemann-hypothesis/blob/40e56bab026552e089d7afcfeff08d2defb9c005/src/Development/AutomorphicArithmetic/AutomorphicVarieties/AdditiveValuationNorms/ExistsCentresLevel.lean\#L34-L49}{\leanfont exists\_centres}}
\expandafter\newcommand\csname LN@centresnorm\endcsname{\href{https://github.com/Deicyde/riemann-hypothesis/blob/40e56bab026552e089d7afcfeff08d2defb9c005/src/Development/AutomorphicArithmetic/AutomorphicVarieties/AdditiveValuationNorms/ExistsCentresLevel.lean\#L57-L69}{\leanfont exists\_centres\_norm}}
\expandafter\newcommand\csname LN@near1\endcsname{\href{https://github.com/Deicyde/riemann-hypothesis/blob/40e56bab026552e089d7afcfeff08d2defb9c005/src/Development/AutomorphicArithmetic/AutomorphicVarieties/LayeredIsomorphisms/Near1.lean\#L6-L43}{\leanfont Win.near1}}
\expandafter\newcommand\csname LN@exc\endcsname{\href{https://github.com/Deicyde/riemann-hypothesis/blob/40e56bab026552e089d7afcfeff08d2defb9c005/src/Development/AutomorphicArithmetic/AutomorphicVarieties/LayeredIsomorphisms/Exc.lean\#L7-L42}{\leanfont Win.exc}}
\expandafter\newcommand\csname LN@qzInj\endcsname{\href{https://github.com/Deicyde/riemann-hypothesis/blob/40e56bab026552e089d7afcfeff08d2defb9c005/src/Development/AutomorphicArithmetic/AutomorphicVarieties/AdditiveValuationNorms/QzInj.lean\#L8-L27}{\leanfont qzInj}}
\expandafter\newcommand\csname LN@dualDiv\endcsname{\href{https://github.com/Deicyde/riemann-hypothesis/blob/40e56bab026552e089d7afcfeff08d2defb9c005/src/Development/AutomorphicArithmetic/AutomorphicVarieties/AdditiveValuationNorms/DualDiv.lean\#L8-L31}{\leanfont dualDiv}}
\expandafter\newcommand\csname LN@torsFin\endcsname{\href{https://github.com/Deicyde/riemann-hypothesis/blob/40e56bab026552e089d7afcfeff08d2defb9c005/src/Development/AutomorphicArithmetic/AutomorphicVarieties/AdditiveValuationNorms/TorsFin.lean\#L9-L84}{\leanfont torsFin}}
\expandafter\newcommand\csname LN@dualUnip\endcsname{\href{https://github.com/Deicyde/riemann-hypothesis/blob/40e56bab026552e089d7afcfeff08d2defb9c005/src/Development/AutomorphicArithmetic/AutomorphicVarieties/AdditiveValuationNorms/DualUnip.lean\#L9-L34}{\leanfont dualUnip}}
\expandafter\newcommand\csname LN@XDual\endcsname{\href{https://github.com/Deicyde/riemann-hypothesis/blob/40e56bab026552e089d7afcfeff08d2defb9c005/src/Development/AutomorphicArithmetic/AutomorphicVarieties/GaloisLibrary/KummerTheory/Mk0IdMapCongr.lean\#L90-L99}{\leanfont XDual}}
\expandafter\newcommand\csname LN@IsTower\endcsname{\href{https://github.com/Deicyde/riemann-hypothesis/blob/40e56bab026552e089d7afcfeff08d2defb9c005/src/Development/AutomorphicArithmetic/AutomorphicVarieties/NumberFieldTowers/XTwApplyApply.lean\#L120-L124}{\leanfont Tower.IsTower}}
\expandafter\newcommand\csname LN@IsLayer\endcsname{\href{https://github.com/Deicyde/riemann-hypothesis/blob/40e56bab026552e089d7afcfeff08d2defb9c005/src/Development/AutomorphicArithmetic/AutomorphicVarieties/StmtFields.lean\#L47-L48}{\leanfont IsLayer}}
\expandafter\newcommand\csname LN@shaInv\endcsname{\href{https://github.com/Deicyde/riemann-hypothesis/blob/40e56bab026552e089d7afcfeff08d2defb9c005/src/Development/AutomorphicArithmetic/AutomorphicVarieties/GaloisLibrary/GaloisCohomologyLength/LevelInertiaLeLevelDec.lean\#L70-L71}{\leanfont shaInv}}
\expandafter\newcommand\csname LN@VLen\endcsname{\href{https://github.com/Deicyde/riemann-hypothesis/blob/40e56bab026552e089d7afcfeff08d2defb9c005/src/Formalism/GaloisCohomology/Statement.lean\#L80-L80}{\leanfont VLen}}
\expandafter\newcommand\csname LN@ABound\endcsname{\href{https://github.com/Deicyde/riemann-hypothesis/blob/40e56bab026552e089d7afcfeff08d2defb9c005/src/Development/AutomorphicArithmetic/AutomorphicVarieties/LayeredIsomorphisms/SecondBound.lean\#L19-L25}{\leanfont Win.ABound}}
\expandafter\newcommand\csname LN@SecondBound\endcsname{\href{https://github.com/Deicyde/riemann-hypothesis/blob/40e56bab026552e089d7afcfeff08d2defb9c005/src/Development/AutomorphicArithmetic/AutomorphicVarieties/LayeredIsomorphisms/SecondBound.lean\#L13-L17}{\leanfont Win.SecondBound}}
\expandafter\newcommand\csname LN@FirstBound\endcsname{\href{https://github.com/Deicyde/riemann-hypothesis/blob/40e56bab026552e089d7afcfeff08d2defb9c005/src/Development/AutomorphicArithmetic/AutomorphicVarieties/LayeredIsomorphisms/SecondBound.lean\#L27-L30}{\leanfont Win.FirstBound}}
\expandafter\newcommand\csname LN@dinfNil\endcsname{\href{https://github.com/Deicyde/riemann-hypothesis/blob/40e56bab026552e089d7afcfeff08d2defb9c005/src/Development/AutomorphicArithmetic/AutomorphicVarieties/GaloisLibrary/KummerTheory/DinfNil.lean\#L62-L235}{\leanfont Kum.dinfNil}}
\expandafter\newcommand\csname LN@xUnip\endcsname{\href{https://github.com/Deicyde/riemann-hypothesis/blob/40e56bab026552e089d7afcfeff08d2defb9c005/src/Development/AutomorphicArithmetic/AutomorphicVarieties/AdditiveValuationNorms/XUnip.lean\#L9-L27}{\leanfont xUnip}}
\expandafter\newcommand\csname LN@step0\endcsname{\href{https://github.com/Deicyde/riemann-hypothesis/blob/40e56bab026552e089d7afcfeff08d2defb9c005/src/Development/AutomorphicArithmetic/AutomorphicVarieties/NumberFieldTowers/ExistsDescends.lean\#L26-L46}{\leanfont Tower.step0}}
\expandafter\newcommand\csname LN@isTowerIff\endcsname{\href{https://github.com/Deicyde/riemann-hypothesis/blob/40e56bab026552e089d7afcfeff08d2defb9c005/src/Development/AutomorphicArithmetic/AutomorphicVarieties/NumberFieldTowers/ExistsDescends.lean\#L60-L62}{\leanfont Tower.isTower\_iff}}
\expandafter\newcommand\csname LN@Descends\endcsname{\href{https://github.com/Deicyde/riemann-hypothesis/blob/40e56bab026552e089d7afcfeff08d2defb9c005/src/Development/AutomorphicArithmetic/AutomorphicVarieties/AdditiveValuationNorms/ABoundOfDescends.lean\#L16-L18}{\leanfont Descends}}
\expandafter\newcommand\csname LN@descends\endcsname{\href{https://github.com/Deicyde/riemann-hypothesis/blob/40e56bab026552e089d7afcfeff08d2defb9c005/src/Development/AutomorphicArithmetic/AutomorphicVarieties/AdditiveValuationNorms/Descends.lean\#L8-L15}{\leanfont descends}}
\expandafter\newcommand\csname LN@cycRange\endcsname{\href{https://github.com/Deicyde/riemann-hypothesis/blob/40e56bab026552e089d7afcfeff08d2defb9c005/src/Development/AutomorphicArithmetic/AutomorphicVarieties/NumberFieldTowers/CycRange.lean\#L9-L62}{\leanfont Tower.cycRange}}
\expandafter\newcommand\csname LN@genLevel\endcsname{\href{https://github.com/Deicyde/riemann-hypothesis/blob/40e56bab026552e089d7afcfeff08d2defb9c005/src/Development/AutomorphicArithmetic/AutomorphicVarieties/AdditiveValuationNorms/GenLevel.lean\#L11-L31}{\leanfont genLevel}}
\expandafter\newcommand\csname LN@shaLeV\endcsname{\href{https://github.com/Deicyde/riemann-hypothesis/blob/40e56bab026552e089d7afcfeff08d2defb9c005/src/Development/AutomorphicArithmetic/AutomorphicVarieties/AdditiveValuationNorms/ShaLeV.lean\#L7-L16}{\leanfont shaLeV}}
\expandafter\newcommand\csname LN@aDesc\endcsname{\href{https://github.com/Deicyde/riemann-hypothesis/blob/40e56bab026552e089d7afcfeff08d2defb9c005/src/Development/AutomorphicArithmetic/AutomorphicVarieties/AdditiveValuationNorms/ABoundOfDescends.lean\#L36-L55}{\leanfont aBound\_of\_descends}}
\expandafter\newcommand\csname LN@aTower\endcsname{\href{https://github.com/Deicyde/riemann-hypothesis/blob/40e56bab026552e089d7afcfeff08d2defb9c005/src/Development/AutomorphicArithmetic/AutomorphicVarieties/AdditiveValuationNorms/ABoundOfIsTower.lean\#L7-L11}{\leanfont aBound\_of\_isTower}}
\expandafter\newcommand\csname LN@FW\endcsname{\href{https://github.com/Deicyde/riemann-hypothesis/blob/40e56bab026552e089d7afcfeff08d2defb9c005/src/Development/AutomorphicArithmetic/AutomorphicVarieties/PublishedTheorems/FerreroWashington.lean\#L7-L11}{\leanfont Published.FerreroWashington}}
\expandafter\newcommand\csname LN@FWthm\endcsname{\href{https://github.com/Deicyde/riemann-hypothesis/blob/40e56bab026552e089d7afcfeff08d2defb9c005/src/Development/AutomorphicArithmetic/AutomorphicVarieties/GaloisLibrary/FerreroWashingtonIndex/FerreroWashington.lean\#L104-L119}{\leanfont FW.ferreroWashington}}
\expandafter\newcommand\csname LN@towerFixed\endcsname{\href{https://github.com/Deicyde/riemann-hypothesis/blob/40e56bab026552e089d7afcfeff08d2defb9c005/src/Development/AutomorphicArithmetic/AutomorphicVarieties/QuadraticExtensions/IsAlgClosureRat.lean\#L98-L110}{\leanfont Two.towerFixed}}
\expandafter\newcommand\csname LN@fixedIn\endcsname{\href{https://github.com/Deicyde/riemann-hypothesis/blob/40e56bab026552e089d7afcfeff08d2defb9c005/src/Development/AutomorphicArithmetic/AutomorphicVarieties/QuadraticExtensions/IsAlgClosureRat.lean\#L139-L141}{\leanfont Two.fixedIn}}
\expandafter\newcommand\csname LN@twoIn\endcsname{\href{https://github.com/Deicyde/riemann-hypothesis/blob/40e56bab026552e089d7afcfeff08d2defb9c005/src/Development/AutomorphicArithmetic/AutomorphicVarieties/QuadraticExtensions/IsAlgClosureRat.lean\#L133-L137}{\leanfont Two.twoIn}}
\expandafter\newcommand\csname LN@LemmaK\endcsname{\href{https://github.com/Deicyde/riemann-hypothesis/blob/40e56bab026552e089d7afcfeff08d2defb9c005/src/Development/AutomorphicArithmetic/AutomorphicVarieties/GeneratorCountTuples/LemmaK.lean\#L13-L17}{\leanfont LemmaK}}
\expandafter\newcommand\csname LN@window\endcsname{\href{https://github.com/Deicyde/riemann-hypothesis/blob/40e56bab026552e089d7afcfeff08d2defb9c005/src/Development/AutomorphicArithmetic/AutomorphicVarieties/LayeredIsomorphisms/Window.lean\#L7-L43}{\leanfont Win.window}}
\expandafter\newcommand\csname LN@thmBab\endcsname{\href{https://github.com/Deicyde/riemann-hypothesis/blob/40e56bab026552e089d7afcfeff08d2defb9c005/src/Development/AutomorphicArithmetic/AutomorphicVarieties/LayeredIsomorphisms/ThmBabOf.lean\#L11-L103}{\leanfont Win.thmBab\_of}}
\expandafter\newcommand\csname LN@thmBabK\endcsname{\href{https://github.com/Deicyde/riemann-hypothesis/blob/40e56bab026552e089d7afcfeff08d2defb9c005/src/Development/AutomorphicArithmetic/AutomorphicVarieties/LayeredIsomorphisms/ThmBabOfLemmaK.lean\#L9-L12}{\leanfont Win.thmBab\_of\_lemmaK}}
\expandafter\newcommand\csname LN@vl\endcsname{\href{https://github.com/Deicyde/riemann-hypothesis/blob/40e56bab026552e089d7afcfeff08d2defb9c005/src/Development/AutomorphicArithmetic/AutomorphicVarieties/AdditiveValuationNorms/CokLenLeOfRangeLe.lean\#L14-L14}{\leanfont vl}}
\expandafter\newcommand\csname LN@normDvd\endcsname{\href{https://github.com/Deicyde/riemann-hypothesis/blob/40e56bab026552e089d7afcfeff08d2defb9c005/src/Development/AutomorphicArithmetic/AutomorphicVarieties/AdditiveValuationNorms/LteNorm.lean\#L7-L10}{\leanfont norm\_le\_of\_dvd'}}
\expandafter\newcommand\csname LN@normAddVal\endcsname{\href{https://github.com/Deicyde/riemann-hypothesis/blob/40e56bab026552e089d7afcfeff08d2defb9c005/src/Development/AutomorphicArithmetic/AutomorphicVarieties/AdditiveValuationNorms/LteNorm.lean\#L12-L17}{\leanfont norm\_eq\_of\_addVal}}
\expandafter\newcommand\csname LN@dvdNorm\endcsname{\href{https://github.com/Deicyde/riemann-hypothesis/blob/40e56bab026552e089d7afcfeff08d2defb9c005/src/Development/AutomorphicArithmetic/AutomorphicVarieties/AdditiveValuationNorms/NormVal.lean\#L6-L19}{\leanfont dvd\_of\_norm\_le'}}
\expandafter\newcommand\csname LN@vlProd\endcsname{\href{https://github.com/Deicyde/riemann-hypothesis/blob/40e56bab026552e089d7afcfeff08d2defb9c005/src/Development/AutomorphicArithmetic/AutomorphicVarieties/AdditiveValuationNorms/NormVal.lean\#L21-L28}{\leanfont vl\_prod}}
\expandafter\newcommand\csname LN@nilSucc\endcsname{\href{https://github.com/Deicyde/riemann-hypothesis/blob/40e56bab026552e089d7afcfeff08d2defb9c005/src/Development/AutomorphicArithmetic/AutomorphicVarieties/AdditiveValuationNorms/NilLen.lean\#L6-L9}{\leanfont nilLen\_range\_succ}}
\expandafter\newcommand\csname LN@nilLt\endcsname{\href{https://github.com/Deicyde/riemann-hypothesis/blob/40e56bab026552e089d7afcfeff08d2defb9c005/src/Development/AutomorphicArithmetic/AutomorphicVarieties/AdditiveValuationNorms/NilLen.lean\#L11-L27}{\leanfont nilLen\_lt}}
\expandafter\newcommand\csname LN@vlNorm\endcsname{\href{https://github.com/Deicyde/riemann-hypothesis/blob/40e56bab026552e089d7afcfeff08d2defb9c005/src/Development/AutomorphicArithmetic/AutomorphicVarieties/AdditiveValuationNorms/ExistsCentresLevel.lean\#L51-L55}{\leanfont vl\_le\_of\_norm}}
\expandafter\newcommand\csname LN@memLev\endcsname{\href{https://github.com/Deicyde/riemann-hypothesis/blob/40e56bab026552e089d7afcfeff08d2defb9c005/src/Development/AutomorphicArithmetic/AutomorphicVarieties/AdditiveValuationNorms/ABoundOfDescends.lean\#L22-L28}{\leanfont mem\_levU\_one}}
\expandafter\newcommand\csname LN@eqTwLen\endcsname{\href{https://github.com/Deicyde/riemann-hypothesis/blob/40e56bab026552e089d7afcfeff08d2defb9c005/src/Development/AutomorphicArithmetic/AutomorphicVarieties/AdditiveValuationNorms/ABoundOfDescends.lean\#L30-L34}{\leanfont Descends.eq\_twLen}}
\expandafter\newcommand\csname LN@qzSep\endcsname{\href{https://github.com/Deicyde/riemann-hypothesis/blob/40e56bab026552e089d7afcfeff08d2defb9c005/src/Development/AutomorphicArithmetic/AutomorphicVarieties/NumberFieldTowers/QzSep.lean\#L47-L66}{\leanfont Tower.qzSep}}
\expandafter\newcommand\csname LN@qzAnn\endcsname{\href{https://github.com/Deicyde/riemann-hypothesis/blob/40e56bab026552e089d7afcfeff08d2defb9c005/src/Development/AutomorphicArithmetic/AutomorphicVarieties/NumberFieldTowers/QzSep.lean\#L8-L45}{\leanfont exists\_qz\_ne\_zero\_annihilators\_of\_ne\_zero}}
\expandafter\newcommand\csname LN@exactDual\endcsname{\href{https://github.com/Deicyde/riemann-hypothesis/blob/40e56bab026552e089d7afcfeff08d2defb9c005/src/Development/AutomorphicArithmetic/AutomorphicVarieties/NumberFieldTowers/LenDual.lean\#L39-L47}{\leanfont exact\_dual}}
\expandafter\newcommand\csname LN@simpleDual\endcsname{\href{https://github.com/Deicyde/riemann-hypothesis/blob/40e56bab026552e089d7afcfeff08d2defb9c005/src/Development/AutomorphicArithmetic/AutomorphicVarieties/NumberFieldTowers/LenDual.lean\#L49-L109}{\leanfont simple\_dual}}
\expandafter\newcommand\csname LN@leDual\endcsname{\href{https://github.com/Deicyde/riemann-hypothesis/blob/40e56bab026552e089d7afcfeff08d2defb9c005/src/Development/AutomorphicArithmetic/AutomorphicVarieties/NumberFieldTowers/LenDual.lean\#L10-L37}{\leanfont le\_dual}}
\expandafter\newcommand\csname LN@lenDual\endcsname{\href{https://github.com/Deicyde/riemann-hypothesis/blob/40e56bab026552e089d7afcfeff08d2defb9c005/src/Development/AutomorphicArithmetic/AutomorphicVarieties/NumberFieldTowers/LenDual.lean\#L111-L137}{\leanfont Tower.lenDual}}
\expandafter\newcommand\csname LN@dualKer\endcsname{\href{https://github.com/Deicyde/riemann-hypothesis/blob/40e56bab026552e089d7afcfeff08d2defb9c005/src/Development/AutomorphicArithmetic/AutomorphicVarieties/NumberFieldTowers/DualKer.lean\#L7-L40}{\leanfont Tower.dualKer}}
\expandafter\newcommand\csname LN@cokerDual\endcsname{\href{https://github.com/Deicyde/riemann-hypothesis/blob/40e56bab026552e089d7afcfeff08d2defb9c005/src/Development/AutomorphicArithmetic/AutomorphicVarieties/NumberFieldTowers/CokerDual.lean\#L8-L19}{\leanfont Tower.cokerDual}}
\expandafter\newcommand\csname LN@pgIndex\endcsname{\href{https://github.com/Deicyde/riemann-hypothesis/blob/40e56bab026552e089d7afcfeff08d2defb9c005/src/Development/AutomorphicArithmetic/AutomorphicVarieties/NumberFieldTowers/PrincGen.lean\#L24-L30}{\leanfont pg\_index}}
\expandafter\newcommand\csname LN@pgExists\endcsname{\href{https://github.com/Deicyde/riemann-hypothesis/blob/40e56bab026552e089d7afcfeff08d2defb9c005/src/Development/AutomorphicArithmetic/AutomorphicVarieties/NumberFieldTowers/PrincGen.lean\#L32-L47}{\leanfont pg\_exists}}
\expandafter\newcommand\csname LN@pgEmult\endcsname{\href{https://github.com/Deicyde/riemann-hypothesis/blob/40e56bab026552e089d7afcfeff08d2defb9c005/src/Development/AutomorphicArithmetic/AutomorphicVarieties/NumberFieldTowers/PrincGen.lean\#L13-L22}{\leanfont pg\_emult}}
\expandafter\newcommand\csname LN@princGen\endcsname{\href{https://github.com/Deicyde/riemann-hypothesis/blob/40e56bab026552e089d7afcfeff08d2defb9c005/src/Development/AutomorphicArithmetic/AutomorphicVarieties/NumberFieldTowers/PrincGen.lean\#L49-L117}{\leanfont Tower.princGen}}
\expandafter\newcommand\csname LN@princClosed\endcsname{\href{https://github.com/Deicyde/riemann-hypothesis/blob/40e56bab026552e089d7afcfeff08d2defb9c005/src/Development/AutomorphicArithmetic/CyclotomicBounds/ClassNumberBounds/TotallyRealRoots/Generate.lean\#L103-L109}{\leanfont princ\_le\_of\_isClosed}}
\expandafter\newcommand\csname LN@lteOdd\endcsname{\href{https://github.com/Deicyde/riemann-hypothesis/blob/40e56bab026552e089d7afcfeff08d2defb9c005/src/Development/AutomorphicArithmetic/CyclotomicBounds/ClassNumberBounds/TotallyRealRoots/EmultiplicityPowSubOneOdd.lean\#L46-L52}{\leanfont emultiplicity\_pow\_sub\_one\_odd}}
\expandafter\newcommand\csname LN@lteSq\endcsname{\href{https://github.com/Deicyde/riemann-hypothesis/blob/40e56bab026552e089d7afcfeff08d2defb9c005/src/Development/AutomorphicArithmetic/CyclotomicBounds/ClassNumberBounds/TotallyRealRoots/EmultiplicityPowSubOneOdd.lean\#L72-L88}{\leanfont emultiplicity\_pow\_sub\_one\_sq}}
\expandafter\newcommand\csname LN@cZero\endcsname{\href{https://github.com/Deicyde/riemann-hypothesis/blob/40e56bab026552e089d7afcfeff08d2defb9c005/src/Development/AutomorphicArithmetic/AutomorphicVarieties/NumberFieldTowers/XTwApplyApply.lean\#L104-L104}{\leanfont Tower.cZero}}
\expandafter\newcommand\csname LN@layU\endcsname{\href{https://github.com/Deicyde/riemann-hypothesis/blob/40e56bab026552e089d7afcfeff08d2defb9c005/src/Development/AutomorphicArithmetic/AutomorphicVarieties/NumberFieldTowers/XTwApplyApply.lean\#L108-L108}{\leanfont Tower.layU}}
\expandafter\newcommand\csname LN@muQ\endcsname{\href{https://github.com/Deicyde/riemann-hypothesis/blob/40e56bab026552e089d7afcfeff08d2defb9c005/src/Development/AutomorphicArithmetic/AutomorphicVarieties/NumberFieldTowers/CycMemPrincIff.lean\#L81-L104}{\leanfont Tower.muQ}}
\expandafter\newcommand\csname LN@galOfCyc\endcsname{\href{https://github.com/Deicyde/riemann-hypothesis/blob/40e56bab026552e089d7afcfeff08d2defb9c005/src/Development/AutomorphicArithmetic/AutomorphicVarieties/GaloisLibrary/KummerTheory/GalOfCyc.lean\#L5-L65}{\leanfont Kum.galOfCyc}}
\expandafter\newcommand\csname LN@galEmbTop\endcsname{\href{https://github.com/Deicyde/riemann-hypothesis/blob/40e56bab026552e089d7afcfeff08d2defb9c005/src/Development/AutomorphicArithmetic/AutomorphicVarieties/GaloisLibrary/SelmerTransport/GalEmbTop.lean\#L32-L69}{\leanfont Base.galEmbTop}}
\expandafter\newcommand\csname LN@luKey\endcsname{\href{https://github.com/Deicyde/riemann-hypothesis/blob/40e56bab026552e089d7afcfeff08d2defb9c005/src/Development/AutomorphicArithmetic/AutomorphicVarieties/NumberFieldTowers/LayerUniq.lean\#L11-L23}{\leanfont lu\_key}}
\expandafter\newcommand\csname LN@layerUniq\endcsname{\href{https://github.com/Deicyde/riemann-hypothesis/blob/40e56bab026552e089d7afcfeff08d2defb9c005/src/Development/AutomorphicArithmetic/AutomorphicVarieties/NumberFieldTowers/LayerUniq.lean\#L25-L43}{\leanfont Tower.layerUniq}}
\expandafter\newcommand\csname LN@layerGen\endcsname{\href{https://github.com/Deicyde/riemann-hypothesis/blob/40e56bab026552e089d7afcfeff08d2defb9c005/src/Development/AutomorphicArithmetic/AutomorphicVarieties/NumberFieldTowers/LayerGen.lean\#L10-L39}{\leanfont Tower.layerGen}}
\expandafter\newcommand\csname LN@IsGen\endcsname{\href{https://github.com/Deicyde/riemann-hypothesis/blob/40e56bab026552e089d7afcfeff08d2defb9c005/src/Development/AutomorphicArithmetic/AutomorphicVarieties/NumberFieldTowers/XTwApplyApply.lean\#L36-L37}{\leanfont Tower.IsGen}}
\expandafter\newcommand\csname LN@dTw\endcsname{\href{https://github.com/Deicyde/riemann-hypothesis/blob/40e56bab026552e089d7afcfeff08d2defb9c005/src/Development/AutomorphicArithmetic/AutomorphicVarieties/NumberFieldTowers/XTwApplyApply.lean\#L56-L56}{\leanfont Tower.dTw}}
\expandafter\newcommand\csname LN@xTw\endcsname{\href{https://github.com/Deicyde/riemann-hypothesis/blob/40e56bab026552e089d7afcfeff08d2defb9c005/src/Development/AutomorphicArithmetic/AutomorphicVarieties/NumberFieldTowers/XTwApplyApply.lean\#L58-L58}{\leanfont Tower.xTw}}
\expandafter\newcommand\csname LN@xTwApply\endcsname{\href{https://github.com/Deicyde/riemann-hypothesis/blob/40e56bab026552e089d7afcfeff08d2defb9c005/src/Development/AutomorphicArithmetic/AutomorphicVarieties/NumberFieldTowers/XTwApplyApply.lean\#L66-L69}{\leanfont Tower.xTw\_apply\_apply}}
\expandafter\newcommand\csname LN@ShaD\endcsname{\href{https://github.com/Deicyde/riemann-hypothesis/blob/40e56bab026552e089d7afcfeff08d2defb9c005/src/Development/AutomorphicArithmetic/AutomorphicVarieties/NumberFieldTowers/XTwApplyApply.lean\#L73-L74}{\leanfont Tower.ShaD}}
\expandafter\newcommand\csname LN@desc\endcsname{\href{https://github.com/Deicyde/riemann-hypothesis/blob/40e56bab026552e089d7afcfeff08d2defb9c005/src/Development/AutomorphicArithmetic/AutomorphicVarieties/NumberFieldTowers/Imports.lean\#L10-L147}{\leanfont Tower.desc}}
\expandafter\newcommand\csname LN@shaLenInv\endcsname{\href{https://github.com/Deicyde/riemann-hypothesis/blob/40e56bab026552e089d7afcfeff08d2defb9c005/src/Development/AutomorphicArithmetic/AutomorphicVarieties/GaloisLibrary/GaloisCohomologyLength/ShaLenEqShaInv.lean\#L58-L60}{\leanfont shaLen\_eq\_shaInv}}
\expandafter\newcommand\csname LN@shaDGen\endcsname{\href{https://github.com/Deicyde/riemann-hypothesis/blob/40e56bab026552e089d7afcfeff08d2defb9c005/src/Development/AutomorphicArithmetic/AutomorphicVarieties/NumberFieldTowers/ShaDGen.lean\#L9-L59}{\leanfont Tower.shaDGen}}
\newcommand{\Pin}{40e56bab}
\newcommand{\PinUrl}{https://github.com/Deicyde/riemann-hypothesis/tree/40e56bab026552e089d7afcfeff08d2defb9c005/src}
\newcommand{\AtlasUrl}{https://nasqret.github.io/riemann-proof-atlas/}
\newcommand{\AtlasClusterUrl}{https://nasqret.github.io/riemann-proof-atlas/\#theorem=385692}
\newcommand{\AtlasCommit}{d05bd171}
\newcommand{\AtlasDecls}{818{,}755}
\newcommand{\AtlasBundles}{4{,}983}
\newcommand{\ClusterSize}{4}
\newcommand{\DirDecls}{42}
\newcommand{\NChecked}{3}
\newcommand{\NCovered}{1}
\newcommand{\SubjectCount}{9}
\newcommand{\SubjectRows}{Divisibility, congruences and arithmetic functions & 13\\Matrices, determinants and linear algebra & 13\\Automorphic representations, trace formulae and theta lifting & 4\\Galois and profinite cohomology & 4\\Global fields, idèles and class groups & 2\\Iwasawa theory and cyclotomic class-group towers & 2\\p-adic arithmetic and power series & 2\\Commutative algebra, ideals and modules & 1\\Local fields, valuation and ramification & 1\\}
\newcommand{\DirFiles}{29}
\newcommand{\DirLines}{1{,}193}
\newcommand{\ConcDecls}{108}
\newcommand{\ConcFiles}{65}
\newcommand{\ExZ}{150}
\newcommand{\ExK}{4}
\newcommand{\ExMzero}{3}
\newcommand{\ExLevels}{5}
\newcommand{\ExExc}{5 and 121}
\newcommand{\ExWorst}{4}
\newcommand{\ExTwoMzero}{4}
\newcommand{\ExTwoRowOne}{2 & 6 & 9 & 12 & 15}
\newcommand{\ExTwoRowNine}{4 & 12 & 15 & 18 & 21}
\newcommand{\ExTwoJump}{8}
\newcommand{\RandCases}{120}
\newcommand{\RandSteps}{360}
\newcommand{\RandSharp}{0}
\newcommand{\UnitMembers}{67}
\newcommand{\UnitInDir}{51}
\newcommand{\UnitOutside}{16}
\newcommand{\UnitInNote}{67}
\newcommand{\UnitElsewhere}{0}
\newcommand{\ElsewhereRows}{}
\newcommand{\UnitProofs}{58}
\newcommand{\UnitSteps}{452}
\newcommand{\ConcordanceRows}{\crow{GenBy}{\lean{GenBy}}{CokLenLeOfRangeLe.lean}{9--10}
\crow{jZero}{\lean{jZero}}{CokLenLeOfRangeLe.lean}{12--12}
\crow{IsLevel}{\lean{IsLevel}}{CokLenLeOfRangeLe.lean}{16--16}
\crow{cokLen}{\lean{cokLen}}{CokLenLeOfRangeLe.lean}{20--20}
\crow{InPow}{\lean{InPow}}{CokLenLeOfRangeLe.lean}{22--22}
\crow{Unip}{\lean{Unip}}{CokLenLeOfRangeLe.lean}{24--24}
\crow{twLen}{\lean{twLen}}{CokLenLeOfRangeLe.lean}{26--26}
\crow{range}{\lean{range}}{CokLenLeOfRangeLe.lean}{28--30}
\crow{lenComp}{\lean{lenComp}}{LenComp.lean}{17--41}
\crow{smith}{\lean{smith}}{LengthDet.lean}{52--68}
\crow{lift}{\lean{lift}}{ImportsAlt.lean}{7--44}
\crow{lte}{\lean{lte}}{IsUnitAddMem.lean}{109--126}
\crow{lteNorm}{\lean{lteNorm}}{LteNorm.lean}{19--94}
\crow{units}{\lean{units}}{PadicUnitsDistanceLeMax.lean}{8--35}
\crow{helly}{\lean{helly}}{Helly.lean}{6--25}
\crow{ball}{\lean{ball}}{PadicUnitsDistanceLeMax.lean}{37--63}
\crow{detFactor}{\lean{detFactor}}{DetFactor.lean}{6--49}
\crow{normVal}{\lean{normVal}}{NormVal.lean}{30--46}
\crow{detBound}{\lean{detBound}}{DetBound.lean}{9--34}
\crow{level}{\lean{level}}{ExistsCentresLevel.lean}{23--32}
\crow{cokL}{\lean{cokL}}{CokLen.lean}{6--41}
\crow{nilLen}{\lean{nilLen}}{NilLen.lean}{29--59}
\crow{unipPow}{\lean{unipPow}}{UnipPow.lean}{8--46}
\crow{frob}{\lean{frob}}{Imports.lean}{6--36}
\crow{powIn}{\lean{powIn}}{PowIn.lean}{6--29}
\crow{phi}{\lean{phi}}{ImportsTer.lean}{7--46}
\crow{step}{\lean{step}}{ImportsMore.lean}{9--22}
\crow{mono}{\lean{mono}}{ImportsBis.lean}{6--13}
\crow{twist}{\lean{twist}}{ImportsFurther.lean}{7--50}
\crow{growth}{\lean{growth}}{Growth.lean}{11--56}
\crow{twgrowth}{\lean{twgrowth}}{ExistsCentresLevel.lean}{17--21}
\crow{centres}{\lean{centres}}{ExistsCentresLevel.lean}{34--49}
\crow{centresnorm}{\lean{centresnorm}}{ExistsCentresLevel.lean}{57--69}
\crow{near1}{\lean{near1}}{Near1.lean}{6--43}
\crow{exc}{\lean{exc}}{Exc.lean}{7--42}
\crow{qzInj}{\lean{qzInj}}{QzInj.lean}{8--27}
\crow{dualDiv}{\lean{dualDiv}}{DualDiv.lean}{8--31}
\crow{torsFin}{\lean{torsFin}}{TorsFin.lean}{9--84}
\crow{dualUnip}{\lean{dualUnip}}{DualUnip.lean}{9--34}
\crow{XDual}{\lean{XDual}}{Mk0IdMapCongr.lean}{90--99}
\crow{IsTower}{\lean{IsTower}}{XTwApplyApply.lean}{120--124}
\crow{IsLayer}{\lean{IsLayer}}{StmtFields.lean}{47--48}
\crow{shaInv}{\lean{shaInv}}{LevelInertiaLeLevelDec.lean}{70--71}
\crow{VLen}{\lean{VLen}}{Statement.lean}{80--80}
\crow{ABound}{\lean{ABound}}{SecondBound.lean}{19--25}
\crow{SecondBound}{\lean{SecondBound}}{SecondBound.lean}{13--17}
\crow{FirstBound}{\lean{FirstBound}}{SecondBound.lean}{27--30}
\crow{dinfNil}{\lean{dinfNil}}{DinfNil.lean}{62--235}
\crow{xUnip}{\lean{xUnip}}{XUnip.lean}{9--27}
\crow{step0}{\lean{step0}}{ExistsDescends.lean}{26--46}
\crow{isTowerIff}{\lean{isTowerIff}}{ExistsDescends.lean}{60--62}
\crow{Descends}{\lean{Descends}}{ABoundOfDescends.lean}{16--18}
\crow{descends}{\lean{descends}}{Descends.lean}{8--15}
\crow{cycRange}{\lean{cycRange}}{CycRange.lean}{9--62}
\crow{genLevel}{\lean{genLevel}}{GenLevel.lean}{11--31}
\crow{shaLeV}{\lean{shaLeV}}{ShaLeV.lean}{7--16}
\crow{aDesc}{\lean{aDesc}}{ABoundOfDescends.lean}{36--55}
\crow{aTower}{\lean{aTower}}{ABoundOfIsTower.lean}{7--11}
\crow{FW}{\lean{FW}}{FerreroWashington.lean}{7--11}
\crow{FWthm}{\lean{FWthm}}{FerreroWashington.lean}{104--119}
\crow{towerFixed}{\lean{towerFixed}}{IsAlgClosureRat.lean}{98--110}
\crow{fixedIn}{\lean{fixedIn}}{IsAlgClosureRat.lean}{139--141}
\crow{twoIn}{\lean{twoIn}}{IsAlgClosureRat.lean}{133--137}
\crow{LemmaK}{\lean{LemmaK}}{LemmaK.lean}{13--17}
\crow{window}{\lean{window}}{Window.lean}{7--43}
\crow{thmBab}{\lean{thmBab}}{ThmBabOf.lean}{11--103}
\crow{thmBabK}{\lean{thmBabK}}{ThmBabOfLemmaK.lean}{9--12}
\crow{vl}{\lean{vl}}{CokLenLeOfRangeLe.lean}{14--14}
\crow{normDvd}{\lean{normDvd}}{LteNorm.lean}{7--10}
\crow{normAddVal}{\lean{normAddVal}}{LteNorm.lean}{12--17}
\crow{dvdNorm}{\lean{dvdNorm}}{NormVal.lean}{6--19}
\crow{vlProd}{\lean{vlProd}}{NormVal.lean}{21--28}
\crow{nilSucc}{\lean{nilSucc}}{NilLen.lean}{6--9}
\crow{nilLt}{\lean{nilLt}}{NilLen.lean}{11--27}
\crow{vlNorm}{\lean{vlNorm}}{ExistsCentresLevel.lean}{51--55}
\crow{memLev}{\lean{memLev}}{ABoundOfDescends.lean}{22--28}
\crow{eqTwLen}{\lean{eqTwLen}}{ABoundOfDescends.lean}{30--34}
\crow{qzSep}{\lean{qzSep}}{QzSep.lean}{47--66}
\crow{qzAnn}{\lean{qzAnn}}{QzSep.lean}{8--45}
\crow{exactDual}{\lean{exactDual}}{LenDual.lean}{39--47}
\crow{simpleDual}{\lean{simpleDual}}{LenDual.lean}{49--109}
\crow{leDual}{\lean{leDual}}{LenDual.lean}{10--37}
\crow{lenDual}{\lean{lenDual}}{LenDual.lean}{111--137}
\crow{dualKer}{\lean{dualKer}}{DualKer.lean}{7--40}
\crow{cokerDual}{\lean{cokerDual}}{CokerDual.lean}{8--19}
\crow{pgIndex}{\lean{pgIndex}}{PrincGen.lean}{24--30}
\crow{pgExists}{\lean{pgExists}}{PrincGen.lean}{32--47}
\crow{pgEmult}{\lean{pgEmult}}{PrincGen.lean}{13--22}
\crow{princGen}{\lean{princGen}}{PrincGen.lean}{49--117}
\crow{princClosed}{\lean{princClosed}}{Generate.lean}{103--109}
\crow{lteOdd}{\lean{lteOdd}}{EmultiplicityPowSubOneOdd.lean}{46--52}
\crow{lteSq}{\lean{lteSq}}{EmultiplicityPowSubOneOdd.lean}{72--88}
\crow{cZero}{\lean{cZero}}{XTwApplyApply.lean}{104--104}
\crow{layU}{\lean{layU}}{XTwApplyApply.lean}{108--108}
\crow{muQ}{\lean{muQ}}{CycMemPrincIff.lean}{81--104}
\crow{galOfCyc}{\lean{galOfCyc}}{GalOfCyc.lean}{5--65}
\crow{galEmbTop}{\lean{galEmbTop}}{GalEmbTop.lean}{32--69}
\crow{luKey}{\lean{luKey}}{LayerUniq.lean}{11--23}
\crow{layerUniq}{\lean{layerUniq}}{LayerUniq.lean}{25--43}
\crow{layerGen}{\lean{layerGen}}{LayerGen.lean}{10--39}
\crow{IsGen}{\lean{IsGen}}{XTwApplyApply.lean}{36--37}
\crow{dTw}{\lean{dTw}}{XTwApplyApply.lean}{56--56}
\crow{xTw}{\lean{xTw}}{XTwApplyApply.lean}{58--58}
\crow{xTwApply}{\lean{xTwApply}}{XTwApplyApply.lean}{66--69}
\crow{ShaD}{\lean{ShaD}}{XTwApplyApply.lean}{73--74}
\crow{desc}{\lean{desc}}{Imports.lean}{10--147}
\crow{shaLenInv}{\lean{shaLenInv}}{ShaLenEqShaInv.lean}{58--60}
\crow{shaDGen}{\lean{shaDGen}}{ShaDGen.lean}{9--59}}
\expandafter\newcommand\csname LQ@aTower\endcsname{\leanquoteline{0}{theorem UpperTheory.V.PropA.aBound\_of\_isTower (hP : GlobalTheory.Published.PoitouTateDuality) (hF : LocalFinitenessK) (hT : LocalTateDualityK)}\leanquoteline{4}{(l : ℕ) [Fact l.Prime] (L₀ : Type) [Field L₀] [NumberField L₀] (R : ℕ) (ht : UpperTheory.V.Tower.IsTower l L₀ R) :}\leanquoteline{4}{UpperTheory.V.Win.ABound l R L₀}}
\expandafter\newcommand\csname LQ@ABound\endcsname{\leanquoteline{0}{def ABound (l R : ℕ) [Fact l.Prime] (L₀ : Type) [Field L₀] : Prop :=}\leanquoteline{2}{(∃ (K : Type) (\_ : NormedField K) (\_ : NormedAlgebra ℚ\_[l] K) (\_ : IsUltrametricDist K) (c R' : ℕ) (s : Fin R' → K),}\leanquoteline{4}{R' ≤ R ∧ (∃ ζ : L₀, IsPrimitiveRoot ζ (l \textasciicircum{} c)) ∧}\leanquoteline{4}{∀ e : ℕ, 1 ≤ e → ∀ (m : ℕ) (L₁ : Type) [Field L₁] [NumberField L₁] [Algebra L₀ L₁], IsLayer l m L₀ L₁ →}\leanquoteline{6}{∀ B : ℕ, (∀ i, (l : ℝ) \textasciicircum{} (-(B : ℤ)) ≤ ‖s i - (e : K)‖) →}\leanquoteline{8}{shaInv L₁ l e ≤ ((R' * (max m (jZero l R) + c + B) : ℕ) : ℕ∞)) ∧}\leanquoteline{2}{SecondBound l R L₀}}
\expandafter\newcommand\csname LQ@SecondBound\endcsname{\leanquoteline{0}{def SecondBound (l R : ℕ) [Fact l.Prime] (L₀ : Type) [Field L₀] : Prop :=}\leanquoteline{2}{∀ e : ℕ, 1 ≤ e → ∀ (m : ℕ) (L₁ : Type) [Field L₁] [NumberField L₁] [Algebra L₀ L₁], IsLayer l m L₀ L₁ →}\leanquoteline{4}{∀ (L₂ : Type) [Field L₂] [NumberField L₂] [Algebra L₀ L₂], IsLayer l (jZero l R) L₀ L₂ →}\leanquoteline{6}{∀ (Ω : Type) [Field Ω] [Algebra L₂ Ω] [IsAlgClosure L₂ Ω] (D : LData L₂ Ω l),}\leanquoteline{8}{shaInv L₁ l e ≤ VLen D e + ((R * (m - jZero l R) : ℕ) : ℕ∞)}}
\expandafter\newcommand\csname LQ@centres\endcsname{\leanquoteline{0}{theorem exists\_centres \{R c : ℕ\} \{γ : Module.End ℤ\_[l] X\} \{u : ℤ\_[l]\} (hX : GenBy l R X) (hγ : Unip γ) (hu : IsLevel l u c) :}\leanquoteline{4}{∃ s : Fin R → ℤ\_[l], ∀ e m : ℕ,}\leanquoteline{6}{twLen γ u e m ≤ ∑ i, (((max m (jZero l R) + c : ℕ) : ℕ∞) + vl l (s i - e))}}
\expandafter\newcommand\csname LQ@growth\endcsname{\leanquoteline{0}{theorem UpperTheory.V.PropA.growth :}\leanquoteline{4}{∀ (l : ℕ) [Fact l.Prime] (X : Type) [AddCommGroup X] [Module ℤ\_[l] X] (R : ℕ) (g : Module.End ℤ\_[l] X), GenBy l R X → Unip g → ∀ m : ℕ,}\leanquoteline{6}{cokLen (g \textasciicircum{} l \textasciicircum{} m - 1) ≤ cokLen (g \textasciicircum{} l \textasciicircum{} jZero l R - 1) + ((R * (m - jZero l R) : ℕ) : ℕ∞)}}
\expandafter\newcommand\csname LQ@step0\endcsname{\leanquoteline{0}{theorem step0 (hμ : HasMuQ l L₀) \{γ : G L₀ Ω\} (hγ : IsGen l ⊤ γ) \{e : ℤ\} (he : e ≠ 0) (m : ℕ) (L₁ : Type) [Field L₁]}\leanquoteline{4}{[NumberField L₁] [Algebra L₀ L₁] (hL : IsLayer l m L₀ L₁) :}\leanquoteline{4}{shaInv L₁ l e = cokerLen (xTw l l (⊤ : Subgroup (G L₀ Ω)) e γ \textasciicircum{} l \textasciicircum{} m - 1)}}


\linespread{1.08}                       % 11 pt on 14.7 pt
\setlength{\parindent}{0pt}
\setlength{\parskip}{5.5pt plus 1.5pt minus 0.5pt}
\setlength{\abovedisplayskip}{7.5pt plus 2pt minus 3pt}
\setlength{\belowdisplayskip}{7.5pt plus 2pt minus 3pt}
\setlength{\abovedisplayshortskip}{3pt plus 2pt}
\setlength{\belowdisplayshortskip}{6pt plus 2pt minus 2pt}
\emergencystretch=1.5em
\clubpenalty=10000 \widowpenalty=10000 \displaywidowpenalty=10000
\relpenalty=7000 \binoppenalty=9000
\raggedbottom
\pagestyle{fancy}\fancyhf{}\renewcommand{\headrulewidth}{0pt}
\fancyfoot[C]{\sffamily\fontsize{7.6}{9}\selectfont\color{inkfaint}\thepage\,/\,\pageref*{LastPage}}
\setlist{topsep=2pt,partopsep=0pt,parsep=0pt,itemsep=4pt,leftmargin=2.1em,labelsep=0.55em}

\newcommand{\Sha}{\text{\normalfont Ш}}   % the letter only; the notations built on it come from the shared table
\newcommand{\Q}{\mathbb{Q}}
\newcommand{\Z}{\mathbb{Z}}
\newcommand{\F}{\mathbb{F}}
\newcommand{\len}{\operatorname{len}}
\newcommand{\cok}{\operatorname{cok}}
\newcommand{\End}{\operatorname{End}}
\newcommand{\Hom}{\operatorname{Hom}}
\newcommand{\Gal}{\operatorname{Gal}}
\newcommand{\diag}{\operatorname{diag}}
\newcommand{\kw}[1]{\emph{#1}}
\newcommand{\m}[1]{\mbox{$#1$}}        % a short formula that must not be broken across lines
% The notations of the functions that the archive defines are made by the build from digests/notation.json and
% this note's accounts.json, into data.tex: one slot for each explicit argument. This holds for the declarations of
% the concordance (\tw, \mz, \vl, \shaInv, \VLen, \cZero, \layU, \XDual, \dTw, \xTw, \ShaD) and for the
% definitions that the text uses without citing them (\GalG, \galOf, \cyc, \QZl, \princ, \levU, \cycU, \Un,
% \Dinf, \Xinf, \xRank, \qzG, \qzTw, \ShaOf). No macro of this preamble carries an argument of a function: the
% prime in the algebraic closure of Q_l and in the logarithm is written out in the text.
\newcommand{\head}[1]{\par\addvspace{17pt plus 4pt minus 2pt}\Needspace{6\baselineskip}%
  \noindent{\fontsize{13.2}{16}\selectfont\bfseries\color{ink}#1}\par\nopagebreak\vspace{-1pt}}
\newcommand{\sect}[2]{\refstepcounter{section}\label{#2}\par\addvspace{19pt plus 4pt minus 2pt}\Needspace{7\baselineskip}%
  \noindent\llap{{\fontsize{13.2}{16}\selectfont\bfseries\color{inkfaint}\thesection}\hspace{9pt}}%
  {\fontsize{13.2}{16}\selectfont\bfseries\color{ink}#1}\par\nopagebreak\vspace{-1pt}}
\newcommand{\subhead}[1]{\par\addvspace{9pt plus 2pt}\Needspace{4\baselineskip}\noindent\textbf{#1}\enspace\ignorespaces}
% where a statement lives in the Lean archive
\newcommand{\leantag}{{\sffamily\bfseries\fontsize{6.3}{8}\selectfont\color{leancol}\addfontfeatures{LetterSpace=7}LEAN}}
\newcommand{\sep}{\unskip\nobreak\hspace{5pt}{\color{inkfaint}\textperiodcentered}\hspace{5pt plus 2pt}\ignorespaces}
\newcommand{\leanline}[1]{\par\nopagebreak\vspace{-2.5pt}{\fontsize{8.9}{12.4}\selectfont\color{inksoft}\raggedright
  \noindent\llap{\leantag\hspace{9pt}}#1\par}\vspace{1pt}}
% a displayed statement in a tinted box: #1 colour, #2 text
\newcommand{\statement}[2]{\par\addvspace{7pt}\begingroup\setlength{\fboxsep}{8.5pt}\setlength{\fboxrule}{0.6pt}%
  \noindent\fcolorbox{#1!30!white}{#1!3!white}{\parbox{\dimexpr\linewidth-2\fboxsep-2\fboxrule\relax}{#2}}%
  \endgroup\par\addvspace{7pt}}
\newcommand{\nt}[2]{\par\noindent\hangindent=1.1em\hangafter=1 #1\enspace{\color{inksoft}#2}\par\vspace{1pt}}
\newcommand{\legdot}[1]{\tikz[baseline=-0.62ex]\fill[#1] (0,0) circle (2.5pt);\hspace{3pt}}
\newcommand{\crow}[4]{\csname CS@#1\endcsname & #2 & \textcolor{inksoft}{#3} & \textcolor{inksoft}{#4}\\}
\newcommand{\leanquoteline}[2]{\par\noindent\hangindent=\dimexpr#1\fontdimen2\font+2.2em\relax\hangafter=1
  \hspace*{#1\fontdimen2\font}#2}
\newcommand{\quoted}[1]{\par\vspace{3pt}{\ttfamily\fontsize{7.5}{10.4}\selectfont\color{leancol}\raggedright
  \setlength{\parskip}{0pt}\leanquote{#1}\par}\vspace{4pt}}
\newcommand{\bx}[3]{{\sffamily\bfseries\fontsize{7.3}{9.2}\selectfont\color{#1!82!black}#2}\\[1.5pt]#3}
\newcommand{\cs}[2]{\expandafter\newcommand\csname CS@#1\endcsname{#2}}

% theorems
\newtheoremstyle{plainx}{7pt plus 2pt}{7pt plus 2pt}{\itshape}{}{\bfseries}{.}{0.5em}{}
\newtheoremstyle{uprightx}{7pt plus 2pt}{7pt plus 2pt}{}{}{\bfseries}{.}{0.5em}{}
\theoremstyle{plainx}
\newtheorem{theorem}{Theorem}[section]
\newtheorem{lemma}[theorem]{Lemma}
\newtheorem{proposition}[theorem]{Proposition}
\newtheorem{corollary}[theorem]{Corollary}
\theoremstyle{uprightx}
\newtheorem{definition}[theorem]{Definition}
\newtheorem{setting}[theorem]{Setting}
\newtheorem{remark}[theorem]{Remark}
\newtheorem{example}[theorem]{Example}
\newenvironment{conventions}{\par\addvspace{5pt}\noindent\textbf{Conventions.}\enspace\ignorespaces}{\par\addvspace{3pt}}
\renewcommand{\qedsymbol}{{\setlength{\fboxsep}{0pt}\setlength{\fboxrule}{0.5pt}\fbox{\rule{0pt}{0.46em}\hspace{0.46em}}}}
\makeatletter
\renewenvironment{proof}[1][Proof]{\par\pushQED{\qed}\normalfont\topsep0pt\partopsep0pt
  \trivlist\item[\hskip\labelsep\itshape #1.]\ignorespaces}{\popQED\endtrivlist\@endpefalse}
\makeatother

% where each declaration of the concordance is treated in the text
\cs{vl}{Def.~\ref{d:val}} \cs{GenBy}{Def.~\ref{d:gen}} \cs{cokLen}{Def.~\ref{d:cok}} \cs{InPow}{Def.~\ref{d:cok}}
\cs{Unip}{Def.~\ref{d:unip}} \cs{IsLevel}{Def.~\ref{d:level}} \cs{twLen}{Def.~\ref{d:tau}} \cs{jZero}{Def.~\ref{d:tau}}
\cs{range}{Lemma~\ref{l:len}(a)} \cs{lenComp}{Lemma~\ref{l:len}(b)} \cs{mono}{Lemma~\ref{l:len}(c)}
\cs{smith}{Lemma~\ref{l:smith}} \cs{lift}{Prop.~\ref{p:lift}}
\cs{lte}{Lemma~\ref{l:lte}} \cs{lteNorm}{Lemma~\ref{l:lte}} \cs{units}{Lemma~\ref{l:units}}
\cs{helly}{Lemma~\ref{l:helly}} \cs{ball}{Prop.~\ref{p:ball}} \cs{detFactor}{Lemma~\ref{l:eigen}}
\cs{normVal}{Prop.~\ref{p:det}, proof} \cs{detBound}{Prop.~\ref{p:det}} \cs{level}{Cor.~\ref{c:level}}
\cs{cokL}{Lemma~\ref{l:small}(a)} \cs{nilLen}{Lemma~\ref{l:small}(b)} \cs{unipPow}{Lemma~\ref{l:unip}(a)}
\cs{frob}{Lemma~\ref{l:cong}(a)} \cs{powIn}{Lemma~\ref{l:cong}(b), (c)} \cs{phi}{Lemma~\ref{l:step}}
\cs{step}{Lemma~\ref{l:step}} \cs{twist}{Lemma~\ref{l:unip}(b)}
\cs{growth}{Prop.~\ref{p:growth}} \cs{twgrowth}{Theorem~B} \cs{centres}{Theorem~A}
\cs{centresnorm}{Cor.~C, the bound} \cs{near1}{Lemma~\ref{l:sep}} \cs{exc}{Cor.~\ref{c:exc}}
\cs{qzInj}{Prop.~\ref{p:dual}(a)} \cs{torsFin}{Prop.~\ref{p:dual}(b)} \cs{dualDiv}{Prop.~\ref{p:dual}(c)}
\cs{dualUnip}{Prop.~\ref{p:dual}(d)}
\cs{XDual}{Def.~\ref{d:dinf}} \cs{IsTower}{Theorem~D} \cs{IsLayer}{Def.~\ref{d:layer}}
\cs{shaInv}{Def.~\ref{d:sha}} \cs{VLen}{Def.~\ref{d:ldata}} \cs{ABound}{Def.~\ref{d:abound}}
\cs{SecondBound}{Def.~\ref{d:abound}} \cs{FirstBound}{Def.~\ref{d:first}} \cs{dinfNil}{Input~1}
\cs{xUnip}{Cor.~\ref{c:unip}} \cs{step0}{Input~2} \cs{Descends}{Input~2} \cs{descends}{Input~2}
\cs{cycRange}{Lemma~\ref{l:layers}(a)} \cs{genLevel}{Lemma~\ref{l:layers}(c)} \cs{shaLeV}{Input~3}
\cs{aDesc}{Prop.~\ref{p:adesc}} \cs{aTower}{Theorem~D}
\cs{isTowerIff}{Lemma~\ref{l:rank}} \cs{FW}{Theorem~\ref{t:fw}} \cs{FWthm}{Theorem~\ref{t:fw}}
\cs{towerFixed}{Prop.~\ref{p:family}} \cs{fixedIn}{Prop.~\ref{p:fixed}(a)} \cs{twoIn}{Prop.~\ref{p:fixed}(b)}
\cs{LemmaK}{Def.~\ref{d:K}} \cs{window}{Lemma~\ref{l:window}} \cs{thmBab}{Prop.~\ref{p:window}(a)}
\cs{thmBabK}{Prop.~\ref{p:window}(b)}
\cs{normDvd}{Lemma~\ref{l:lte}, proof} \cs{normAddVal}{Lemma~\ref{l:lte}, proof}
\cs{dvdNorm}{Prop.~\ref{p:det}, proof} \cs{vlProd}{Prop.~\ref{p:det}, proof} \cs{nilSucc}{Lemma~\ref{l:small}(b), proof}
\cs{nilLt}{Lemma~\ref{l:small}(b), proof} \cs{vlNorm}{Cor.~C, the bound} \cs{memLev}{Cor.~\ref{c:unip}}
\cs{eqTwLen}{Input~2} \cs{qzSep}{Lemma~\ref{l:char}(a)} \cs{qzAnn}{Lemma~\ref{l:char}(a), proof}
\cs{exactDual}{Lemma~\ref{l:char}(b)} \cs{simpleDual}{Lemma~\ref{l:char}(c), proof} \cs{leDual}{Lemma~\ref{l:char}(c), proof}
\cs{lenDual}{Lemma~\ref{l:char}(c)} \cs{dualKer}{Lemma~\ref{l:cokdual}} \cs{cokerDual}{Lemma~\ref{l:cokdual}}
\cs{pgIndex}{Lemma~\ref{l:princ}(a)} \cs{pgExists}{Lemma~\ref{l:princ}(c), proof} \cs{pgEmult}{Lemma~\ref{l:princ}(b)}
\cs{princGen}{Lemma~\ref{l:princ}} \cs{princClosed}{Lemma~\ref{l:princ}(c), proof} \cs{lteOdd}{Lemma~\ref{l:princ}(b), proof}
\cs{lteSq}{Lemma~\ref{l:princ}(b), proof} \cs{cZero}{Def.~\ref{d:czero}} \cs{layU}{Def.~\ref{d:lay}}
\cs{muQ}{Lemma~\ref{l:roots}} \cs{galOfCyc}{Lemma~\ref{l:roots}(b)} \cs{galEmbTop}{Lemma~\ref{l:layers}(a), proof}
\cs{luKey}{Lemma~\ref{l:layers}(b)} \cs{layerUniq}{Lemma~\ref{l:layers}(a), (b)} \cs{layerGen}{Lemma~\ref{l:layers}(a), (c)}
\cs{IsGen}{Def.~\ref{d:isgen}} \cs{dTw}{Def.~\ref{d:tw}} \cs{xTw}{Def.~\ref{d:tw}} \cs{xTwApply}{Lemma~\ref{l:twdual}}
\cs{ShaD}{Def.~\ref{d:tw}} \cs{desc}{Input~2a} \cs{shaLenInv}{Def.~\ref{d:ldata}} \cs{shaDGen}{Input~2b}

\begin{document}
\color{ink}

% ------------------------------------------------------------------------------------- title
{\raggedright
{\sffamily\bfseries\fontsize{7.2}{9}\selectfont\color{leancol}\addfontfeatures{LetterSpace=8}DEFORMALIZATION 1\par}
\vspace{5pt}
{\fontsize{21.5}{25.5}\selectfont\bfseries Twisted cokernels in a cyclotomic tower\par}
\vspace{5pt}
{\fontsize{12.4}{16}\selectfont\itshape\color{inksoft} One cluster of the Lean archive behind the proof of the Riemann
hypothesis announced by L.\,Alpöge, read back into mathematics\par}}
\vspace{8pt}
{\color{rulegrey}\hrule height 0.6pt}
\vspace{7pt}
\thispagestyle{empty}

\textbf{Summary.} Let $X$ be a module over the $\ell$-adic integers generated by $R$ elements, $\gamma$ an
endomorphism of~$X$ that is unipotent modulo~$\ell$, and \m{u\ne1} an element of \m{1+\ell\Z_\ell} (of
\m{1+4\Z_2} if \m{\ell=2}). We study the lengths $\tw{\gamma}{u}{h}{m}$ of
the cokernels of $(u^h\gamma)^{\ell^m}-1$ as the \kw{twist}~$h$ and the \kw{layer}~$m$ vary. There are $R$ points
\m{c_1,\dots,c_R} of~$\Z_\ell$, the \kw{centres}, such that $\tw{\gamma}{u}{h}{m}$ exceeds a quantity depending only on the layer by at most
$\sum_i \vl{\ell}{c_i-h}$, where $\vl{\ell}{x}$ is the $\ell$-adic valuation of~$x$. Hence among the twists \m{1\le h\le Z} at most~$R$ are exceptional, however large $Z$ is.
Applied to the Iwasawa module of the cyclotomic $\Z_\ell$-extension of a number field, this bounds the locally
trivial cohomology of the Tate twists of~$\QZl{\ell}$ along the tower. The proof is linear algebra over~$\Z_\ell$ and one
compactness argument; it uses neither the structure theorem for Iwasawa modules nor the $\ell$-adic logarithm.

\statement{leancol}{\fontsize{9.6}{13.2}\selectfont\textbf{What this note is.} On 20 November 2026 L.\,Alpöge
announced a proof of the Riemann hypothesis, with a Lean archive \cite{Al}. That is a claim, and nothing here asserts
that the hypothesis is proved. This note takes one cluster of the archive, the \ClusterSize{} declarations that the
embedding atlas~\cite{At} groups as \href{\AtlasClusterUrl}{\emph{Assemble the arithmetic tower window}}, together
with the directory on which they rest (\DirFiles{} files, \DirLines{} lines, \DirDecls{} theorems and lemmas), and
rewrites it as mathematics. The algebra is proved in full (§\ref{s:len}--§\ref{s:exc} and
Proposition~\ref{p:dual}). Of the three arithmetic inputs (§\ref{s:arith}), two are stated with a sketch, and the
third, the descent to a layer of the tower, is proved in §\ref{s:descent} from two statements of a neighbouring
directory. Every statement was
compared with its formal counterpart at commit \href{\PinUrl}{{\leanfont\Pin}} of the public mirror; Appendix~A is
the concordance of \ConcDecls{} declarations. Formal status: \NChecked{} of the \ClusterSize{} declarations were
replayed by name by the reading project (type, axioms, kernel), the fourth lies in the closure of one of them. The complete proof, of which the cluster is a part, is the
subject of the replays recorded in the ledger of this project.}

% ------------------------------------------------------------------------------------- 1
\sect{The question}{s:intro}
In the announced argument a chain of theta lifts is compared with a sum over primes, and each lift loses a factor.
For a prime~$\ell$ the loss is governed by groups of Galois cohomology: for a number field~$E$ and an integer
\m{h\ge1}, by the length of the group $\Sha^1_\ell(E,h)$ of classes with coefficients in the $(-h)$-th Tate twist of~$\QZl{\ell}$ that are unramified
outside~$\ell$ and trivial at the places above~$\ell$ (the definition is recalled in §\ref{s:arith}). The proof needs
the total length over all twists \m{1\le h\le Z}, for a very large~$Z$, to be small on a whole range of scales;
see §6 of the guide~\cite{Gu}. A bound for each $h$ that grows with $Z$ is useless. What one can hope for is a bound
that holds for all twists except a set of \kw{exceptional} ones, \emph{whose number does not depend on~$Z$}.

Along the cyclotomic tower over~$E$ these lengths are lengths of cokernels of one operator on one module. This note
proves the statement about exceptional twists for such operators, in the generality in which the archive proves it.

\subhead{The simplest case.} Let $X=\Z_\ell$, with $\ell$ odd, let $u=1+\ell$, and let $\gamma$ be multiplication by
a principal unit, which we write as $u^{-c}$ with \m{c\in\Z_\ell}. For integers \m{h,m\ge0} the operator
$(u^h\gamma)^{\ell^m}-1$ is multiplication by \m{u^{\ell^m(h-c)}-1}, an element of valuation \m{1+m+\vl{\ell}{h-c}}.
So the cokernel has length
\[
\tw{\gamma}{u}{h}{m}=1+m+\vl{\ell}{h-c}.
\]
It is large only when the integer $h$ is $\ell$-adically close to~$c$. Now two integers \m{1\le h<h'\le Z} cannot
both satisfy \m{\vl{\ell}{h-c}>\log_\ell Z}: their difference would be divisible by a power of~$\ell$ larger than~$Z$.
So at most one twist \m{h\le Z} is exceptional, whatever~$Z$. The general case has $R$ eigenvalues in place of
one, hence $R$ centres; the work lies in doing without a diagonal form, and in allowing torsion in~$X$.

% ------------------------------------------------------------------------------------- 2
\sect{Setting and results}{s:results}
\begin{conventions}\phantomsection\label{c:conventions}
Throughout, $\ell$ is a prime number, $\Z_\ell$ is the ring of $\ell$-adic integers, $\Q_\ell$ its field of
fractions and $\overline{\Q}_\ell$ an algebraic closure of~$\Q_\ell$. A \kw{module} is a module over~$\Z_\ell$,
unless another ring is named, and the letters $X$, $M$ and~$D$ stand for modules. The \kw{length} $\len M$ of a
module~$M$ is the supremum of the lengths of its chains of submodules, a non-negative integer or~$\infty$; it is
additive in short exact sequences, so the length of a submodule or of a quotient of~$M$ is at most $\len M$. A
finite module of order $\ell^n$ has length~$n$. We write $\End(X)$ for the ring of endomorphisms of~$X$, $1$ for
the identity, and identify a scalar with the multiplication by it. In §\ref{s:arith} and §\ref{s:descent} the
letters $F$ and~$L$ stand for number fields, and a field with a bar for an algebraic closure of the field.
\end{conventions}

\begin{definition}[valuation]\label{d:val}
For \m{x\in\Z_\ell} the \kw{$\ell$-adic valuation} $\vl{\ell}{x}$ is the largest integer \m{n\ge0} with
\m{x\in\ell^n\Z_\ell}, and \m{\vl{\ell}{0}=\infty}; thus \m{\vl{\ell}{\ell}=1}. It is a function of the prime and
of~$x$, and both are written every time. The same notation is used for the extension of the valuation to
\m{z\in\overline{\Q}_\ell}, and \m{|z|=\ell^{-\vl{\ell}{z}}} is the absolute value. It is ultrametric:
\m{|z+z'|\le\max(|z|,|z'|)}.
\end{definition}
\leanline{\lean{vl}}

\begin{definition}[generators]\label{d:gen}
Let $X$ be a module and \m{R\ge0} an integer. The module $X$ is \kw{generated by $R$ elements} if there is a
surjective homomorphism of modules \m{\pi\colon\Z_\ell^R\to X}. No such $R$ is fixed in this note, and every
statement names its own.
\end{definition}
\leanline{\lean{GenBy}}

\begin{definition}[cokernels and congruences]\label{d:cok}
Let $X$ be a module and \m{f\in\End(X)}. The cokernel of~$f$ is
\[
\cok f=X/f(X),
\]
and $\len\cok f$, an element of \m{\{0,1,2,\dots,\infty\}}, is the length of the cokernel of~$f$. For an integer
\m{n\ge0} we write \m{f\equiv0\pmod{\ell^n}} if \m{f(X)\subseteq\ell^nX}, and \m{f\equiv g\pmod{\ell^n}}, for
\m{g\in\End(X)}, if \m{f-g\equiv0\pmod{\ell^n}}. The set $I_n$ of the \m{f\in\End(X)} with
\m{f\equiv0\pmod{\ell^n}} is a two-sided ideal, \m{I_iI_j\subseteq I_{i+j}}, and it contains $\ell^n\End(X)$. (If
$X$ has torsion, $I_n$ can be larger than $\ell^n\End(X)$.)
\end{definition}
\leanline{\lean{cokLen}\sep\lean{InPow}}

\begin{definition}\label{d:unip}
Let $X$ be a module. An endomorphism $\gamma$ of $X$ is \kw{unipotent modulo~$\ell$} if
\m{(\gamma-1)^k\equiv0\pmod\ell} for some integer \m{k\ge1}.
\end{definition}
\leanline{\lean{Unip}}

\begin{definition}\label{d:level}
An element $u$ of $\Z_\ell$ has \kw{level}~$a$, for an integer \m{a\ge1}, if \m{\vl{\ell}{u-1}=a}, and \m{a\ge2} in case
\m{\ell=2}. Such a $u$ is a unit, and it generates the group \m{1+\ell^a\Z_\ell} topologically: this group is the
closure of the set of powers of~$u$ (Lemma~\ref{l:princ}).
\end{definition}
\leanline{\lean{IsLevel}}

\begin{definition}\label{d:tau}
Let $X$ be a module, \m{\gamma\in\End(X)} and \m{u\in\Z_\ell}. For integers \m{h,m\ge0}, the \kw{twisted length}
of the pair $(\gamma,u)$ at twist~$h$ and layer~$m$ is
\[
\tw{\gamma}{u}{h}{m}=\len\cok\bigl((u^h\gamma)^{\ell^m}-1\bigr)\ \in\ \{0,1,2,\dots,\infty\}.
\]
It is a function of four arguments, and all four are written every time: the endomorphism and the scalar are
never dropped from the notation. The definition asks nothing of $\gamma$ and~$u$; unipotence and the level are
hypotheses of the theorems below, not of the definition.

For integers \m{R\ge0} the \kw{reference layer} is \m{\mz{\ell}{R}=\lfloor\log_\ell R\rfloor+3}, with
\m{\mz{\ell}{0}=3}. It depends on the prime and on~$R$ only.
\end{definition}
\leanline{\lean{twLen}\sep\lean{jZero}}

\subhead{Results.} The three statements that follow have the same hypotheses. They are written out once, and
each statement cites them.

\begin{setting}\label{set:alg}
$X$ is a module generated by $R$ elements, for an integer \m{R\ge0}; $\gamma$ is an endomorphism of~$X$ that is
unipotent modulo~$\ell$; and $u$ is an element of~$\Z_\ell$ of level~$a$, for an integer \m{a\ge1}.
\end{setting}

\statement{arith}{\phantomsection\label{t:A}\textbf{Theorem A} (centres).\enspace\emph{In Setting~\ref{set:alg} there are
\m{c_1,\dots,c_R\in\Z_\ell}, chosen for the pair $(\gamma,u)$ and the same for all twists and layers, such that for
all integers \m{h,m\ge0}}
\[
\tw{\gamma}{u}{h}{m}\ \le\ \sum_{i=1}^{R}\Bigl(\max\bigl(m,\mz{\ell}{R}\bigr)+a+\vl{\ell}{c_i-h}\Bigr).
\]
\phantomsection\label{t:B}\textbf{Theorem B} (growth).\enspace\emph{In Setting~\ref{set:alg}, for all integers \m{h,m\ge0}}
\[
\tw{\gamma}{u}{h}{m}\ \le\ \tw{\gamma}{u}{h}{\mz{\ell}{R}}+R\cdot\max\bigl(m-\mz{\ell}{R},0\bigr).
\]}
\statement{arith}{\phantomsection\label{t:C}\textbf{Corollary C} (exceptional twists).\enspace\emph{In Setting~\ref{set:alg}, let
\m{c_1,\dots,c_R} be centres of the pair $(\gamma,u)$ as in Theorem~A, and let \m{Z\ge1} be an integer. The set
$\mathcal E_Z$ of integers \m{1\le h\le Z} such that \m{\vl{\ell}{c_i-h}>\lfloor\log_\ell Z\rfloor} for some~$i$ has at most $R$ elements. For
every integer \m{1\le h\le Z} outside~$\mathcal E_Z$ and every \m{m\ge0}}
\[
\tw{\gamma}{u}{h}{m}\ \le\ R\,\bigl(\max\bigl(m,\mz{\ell}{R}\bigr)+a+\lfloor\log_\ell Z\rfloor\bigr).
\]}
\leanline{A \lean{centres}\sep B \lean{twgrowth}\sep C: the bound \lean{centresnorm}, with \lean{vlNorm}; the
count is proved in the archive for the lengths of Theorem~D (Corollary~\ref{c:exc})}

The order of the choices matters. The module, the endomorphism~$\gamma$ and the scalar~$u$ are given first; the
centres are then chosen for them, and a different $\gamma$ or a different $u$ has different centres. Only after
that do the twist~$h$ and the layer~$m$ vary. In §\ref{s:arith} the pair $(\gamma,u)$ is a particular one: the
action of an element $\sigma$ of a Galois group on the Iwasawa module, and the value of the cyclotomic character
at~$\sigma$.

The two theorems complement each other. Theorem~A is strong when $h$ is far from every centre, and says nothing
when $h$ is a centre: then its right-hand side is infinite, and $\tw{\gamma}{u}{h}{m}$ may be infinite too
(Example~\ref{e:one}). Theorem~B gives no information on any single layer, but it shows that all the layers are
controlled by one of them, also at a centre. In the arithmetic application that one layer is bounded by a duality
theorem, which is why both are needed.

The arithmetic statement concerns a number field~$L$ that contains the $\ell$-th roots of unity (the fourth roots
if \m{\ell=2}), the layers~$L_m$ of its cyclotomic $\Z_\ell$-extension, and an Iwasawa module: the Galois group of
the maximal abelian pro-$\ell$ extension of \m{L_\infty=\bigcup_mL_m} that is unramified everywhere and in which
every prime above~$\ell$ splits completely. In the archive this module is a function of three arguments, written
$\Xinf{\ell}{N}{U}$ here (Definition~\ref{d:dinf}): the prime, an integer~$N$ whose prime divisors are the places
where complete splitting is asked, and a subgroup~$U$ of the Galois group. The module of the theorem is the one
with \m{N=\ell} and with $U$ the whole group. The fields are set up in Setting~\ref{set:tower} of §\ref{s:arith},
the three statements of Galois cohomology that the archive assumes are listed in Setting~\ref{set:coh}, and the
theorem cites both.

\statement{arith}{\phantomsection\label{t:D}\textbf{Theorem D} (class groups in the tower).\enspace\emph{In
Setting~\ref{set:tower}, with the number field~$L$, its algebraic closure~$\overline L$ and the layers~$L_m$ of that
setting, assume the three hypotheses on Galois cohomology of Setting~\ref{set:coh}. Write
\m{G=\GalG{L}{\overline L}} and \m{a=\cZero{\ell}{L}} (Definition~\ref{d:czero}), and suppose that the module
$\Xinf{\ell}{\ell}{G}$ is generated by \m{R\ge1} elements. Then}
\begin{enumerate}[label=\textup{(\roman*)},topsep=3pt,itemsep=1pt]
\item \emph{there are \m{c_1,\dots,c_R\in\Z_\ell}, depending on~$L$, such that for all integers \m{h\ge1} and
\m{m,B\ge0} with \m{\vl{\ell}{c_i-h}\le B} for all~$i$}
\[
\shaInv{L_m}{\ell}{h}\ \le\ R\,\bigl(\max\bigl(m,\mz{\ell}{R}\bigr)+a+B\bigr);
\]
\item \emph{for all integers \m{h\ge1} and \m{m\ge0}, writing \m{n=\mz{\ell}{R}} for the reference layer, and
for every local datum $\mathcal D$ of the field~$L_n$ at~$\ell$ (Definition~\ref{d:ldata}),}
\[
\shaInv{L_m}{\ell}{h}\ \le\ \VLen{\mathcal D}{h}+R\cdot\max(m-n,0).
\]
\end{enumerate}
\emph{In particular the tower bound of Definition~\ref{d:abound} holds for $(\ell,R,L)$.}}
\leanline{\lean{aTower}\sep hypothesis \lean{IsTower}\sep the conclusion in the form of the
archive: \lean{ABound}, \lean{SecondBound} (Definition~\ref{d:abound})\sep the same conclusion from one
descending element: \lean{aDesc} (Proposition~\ref{p:adesc})}

Here $\shaInv{F}{\ell}{h}$ is the length of the group of classes of §\ref{s:intro} for a number field~$F$
(Definition~\ref{d:sha}), and $\VLen{\mathcal D}{h}$ is the length of a second cohomology group
(Definition~\ref{d:ldata}). The three hypotheses of Setting~\ref{set:coh} (global duality of Poitou and Tate,
finiteness of local cohomology, local Tate duality) enter only through Input~3 of §\ref{s:arith}.
\subhead{Plan of the proof.} The diagram shows the steps in the order in which they are used. The upper row
produces the centres at one layer; the middle row compares the layers; the lower row is the dictionary with
arithmetic.

\par\vspace{5pt}
\noindent\makebox[\linewidth][c]{%
\begin{tikzpicture}[x=1cm,y=1cm,
  box/.style 2 args={draw=#1!55!white,line width=0.6pt,rounded corners=3pt,fill=#1!4!white,minimum height=1.62cm,
    text width=#2,inner xsep=4.5pt,inner ysep=3.5pt,align=flush left,font=\fontsize{7.9}{9.6}\selectfont},
  flow/.style={-{Stealth[length=4.6pt,width=4pt]},line width=0.75pt,draw=ink!72}]
  \node[box={est}{3.05cm}] (a1) at (0,0)     {\bx{est}{§3\quad Presentation}{cokernel length \m{\le} valuation of a determinant}};
  \node[box={est}{3.05cm}] (a2) at (3.95,0)  {\bx{est}{§4\quad Eigenvalues}{\m{\det(tS-1)=\prod_i(tz_i-1)}}};
  \node[box={est}{3.05cm}] (a3) at (7.9,0)   {\bx{est}{§4\quad One centre each}{exponent lifting, ultrametric inequality, compactness}};
  \node[box={arith}{3.05cm}] (a4) at (11.85,0) {\bx{arith}{Theorem A, Corollary C}{one set of centres for all layers; at most $R$ exceptional twists (§6)}};
  \node[box={ana}{3.05cm}] (b1) at (0,-2.05)    {\bx{ana}{§5\quad Unipotence}{\m{(\gamma-1)^R\equiv0} modulo~$\ell$; stable under twists}};
  \node[box={ana}{3.05cm}] (b2) at (3.95,-2.05) {\bx{ana}{§5\quad Congruences}{\m{g^{\ell^n}\equiv1} modulo~$\ell^2$ from a small layer on}};
  \node[box={ana}{3.05cm}] (b3) at (7.9,-2.05)  {\bx{ana}{§5\quad One step}{from layer $n$ to \m{n+1} the length grows by at most~$R$}};
  \node[box={arith}{3.05cm}] (b4) at (11.85,-2.05) {\bx{arith}{Theorem B}{every layer is controlled by the layer~$\mz{\ell}{R}$}};
  \node[box={inksoft}{3.05cm}] (c1) at (0,-4.1)    {\bx{inksoft}{§7\quad Dual module}{locally nilpotent on a discrete module, so unipotent on its dual}};
  \node[box={inksoft}{3.05cm}] (c2) at (3.95,-4.1) {\bx{inksoft}{§7, §8\quad Descent}{the length of the classes over a layer is a twisted length on the Iwasawa module}};
  \node[box={inksoft}{3.05cm}] (c3) at (7.9,-4.1)  {\bx{inksoft}{§7\quad Poitou--Tate}{at the layer~$\mz{\ell}{R}$ the length is at most that of a second cohomology group}};
  \node[box={arith}{3.05cm}] (c4) at (11.85,-4.1)  {\bx{arith}{Theorem D}{class groups in the cyclotomic tower}};
  \draw[flow] (a1) -- (a2); \draw[flow] (a2) -- (a3); \draw[flow] (a3) -- (a4);
  \draw[flow] (b1) -- (b2); \draw[flow] (b2) -- (b3); \draw[flow] (b3) -- (b4);
  \draw[flow] (c1) -- (c2); \draw[flow] (c2) -- (c3); \draw[flow] (c3) -- (c4);
  \draw[flow] (b4) -- (a4); \draw[flow] (b4) -- (c4);
  \draw[flow,rounded corners=3pt] (a4.east) -- ++(0.3,0) |- (c4.east);
\end{tikzpicture}}
\par\vspace{3pt}

\subhead{The cluster on the map.} The embedding atlas~\cite{At} places the \AtlasDecls{} written theorems and lemmas
of the archive on a map of mathematical subjects and groups them into \AtlasBundles{} clusters. The cluster treated
here has \ClusterSize{} members: Theorem~D in two forms, and two statements that assemble it into the window of the
guide (§\ref{s:use}). The \DirDecls{} theorems and lemmas of the directory that proves Theorem~D are filed by the
atlas under \SubjectCount{} subjects:

\par\vspace{2pt}
{\fontsize{9.2}{12}\selectfont\noindent\hspace*{2.1em}\begin{tabular}{@{}l@{\hspace{14pt}}r@{}}
\SubjectRows
\end{tabular}\par}
\vspace{2pt}

The first two rows, more than half of the directory, are the algebra of §\ref{s:len}--§\ref{s:exc} and of
Proposition~\ref{p:dual}. The statement is one of Galois cohomology and Iwasawa theory, and almost all of its
proof is linear algebra and congruences. A map of subjects shows where the pieces lie; the
order in which they are used is the diagram above. (Four of the algebraic lemmas of §\ref{s:growth} and
§\ref{s:arith} are filed under automorphic representations, by the context of their module.)

% ------------------------------------------------------------------------------------- 3
\sect{Lengths and determinants}{s:len}
\begin{lemma}\label{l:len}
Let $X$ be a module over~$\Z_\ell$ and let \m{f,g\in\End(X)}.
\begin{enumerate}[label=\textup{(\alph*)},topsep=2pt,itemsep=2pt]
\item If \m{g(X)\subseteq f(X)}, then \m{\len\cok f\le\len\cok g}.
\item \m{\len\cok(fg)\le\len\cok f+\len\cok g}.
\item If \m{i,j\ge0} are integers and $i$ divides~$j$, then \m{\len\cok(g^i-1)\le\len\cok(g^j-1)}.
\end{enumerate}
\end{lemma}
\begin{proof}
(a) $\cok f$ is a quotient of $\cok g$. (b) The module $X/fg(X)$ has the submodule $f(X)/fg(X)$ with quotient
$X/f(X)$, and $f$ induces a surjection of $X/g(X)$ onto $f(X)/fg(X)$. (c) If \m{j=ik}, the factorization
\m{g^j-1=(g^i-1)(1+g^i+\dots+g^{i(k-1)})} shows that the image of \m{g^j-1} lies in that of \m{g^i-1}; apply~(a).
For \m{j=0} the right-hand side is the length of~$X$.
\end{proof}
\leanline{(a) \lean{range}\sep (b) \lean{lenComp}, which also has the same inequality for kernels and holds over
any commutative ring\sep (c) \lean{mono}}

\begin{lemma}\label{l:smith}
For a square matrix $A$ with entries in~$\Z_\ell$, of size~$R$,\quad
\m{\len\,\Z_\ell^R/A\Z_\ell^R=\vl{\ell}{\det A}.}
\end{lemma}
\begin{proof}
By the theory of elementary divisors there are invertible matrices $P$, $Q$ over~$\Z_\ell$ with
\m{PAQ=\diag(d_1,\dots,d_R)}. The cokernel of~$A$ is isomorphic to \m{\bigoplus_i\Z_\ell/d_i\Z_\ell}, the module
$\Z_\ell/d_i\Z_\ell$ has length $\vl{\ell}{d_i}$, and \m{\vl{\ell}{\det A}=\sum_i\vl{\ell}{d_i}}. Both sides are infinite when \m{\det A=0}.
\end{proof}
\leanline{\lean{smith}, there for a square matrix over any discrete valuation ring}

\begin{proposition}[presentation]\label{p:lift}
Let $X$ be a module generated by $R$ elements and let \m{g\in\End(X)}. There is a square matrix $S$
over~$\Z_\ell$, of size~$R$, such that for every \m{t\in\Z_\ell}
\[
\len\cok(tg-1)\ \le\ \vlb{\ell}{\det(tS-1)}.
\]
\end{proposition}
\begin{proof}
Let \m{\pi\colon\Z_\ell^R\to X} be a surjection. Lifting the images under \m{g\circ\pi} of the basis vectors gives a
matrix $S$ with \m{\pi\circ S=g\circ\pi}. Then \m{\pi\circ(tS-1)=(tg-1)\circ\pi}, so $\pi$ maps the image of $tS-1$
into the image of $tg-1$ and induces a surjection of \m{\Z_\ell^R/(tS-1)\Z_\ell^R} onto $\cok(tg-1)$. Apply
Lemma~\ref{l:smith}.
\end{proof}
\leanline{\lean{lift}}

\begin{remark}\label{r:fitting}
The inequality is an equality when $X$ is free of rank~$R$, and can be far from one otherwise. For
\m{X=\Z_\ell/\ell\Z_\ell}, \m{g=1} and \m{t\equiv1\pmod\ell} the left-hand side is~$1$ and the right-hand side, for the lift \m{S=1}, is
\m{\vl{\ell}{t-1}}. Determinants see a presentation of the module, not the module.
\end{remark}

% ------------------------------------------------------------------------------------- 4
\sect{A centre for every eigenvalue}{s:centres}
\begin{lemma}[lifting the exponent]\label{l:lte}
Let \m{u\in\Z_\ell} have level~$a$. Then \m{\vl{\ell}{u^n-1}=a+\vl{\ell}{n}} for every integer~$n$.
Consequently, for all integers \m{j,n,n'\ge0},
\[
\bigl|u^{\ell^jn}-u^{\ell^jn'}\bigr|=\ell^{-(j+a)}\,|n-n'|.
\]
The first formula is the case of the discrete valuation ring~$\Z_\ell$, which has characteristic zero and whose
maximal ideal contains the prime~$\ell$, of a statement of the archive on such rings: there the hypothesis on a
unit~$u$ congruent to~$1$ is \m{\vl{\ell}{\ell}<(\ell-1)\,\vl{\ell}{u-1}}, with the valuation of the ring, and
for~$\Z_\ell$ it says that $u$ has a level.
\end{lemma}
\begin{proof}
Write \m{u=1+x}, so \m{\vl{\ell}{x}=a}. If $\ell$ does not divide $n$, then \m{u^n-1=nx+x^2(\cdots)} has valuation~$a$.
For \m{n=\ell},
\[
u^\ell-1=\ell x+\sum_{1<i<\ell}\binom{\ell}{i}x^i+x^\ell ,
\]
where the first term has valuation \m{a+1}, the middle terms have valuation at least \m{2a+1}, and the last has
valuation~$\ell a$, which exceeds \m{a+1} because \m{(\ell-1)a>1}. This is where \m{a\ge2} is needed for \m{\ell=2}.
So $u^\ell$ has level \m{a+1}, and induction on $\vl{\ell}{n}$ gives the formula for \m{n>0}. For \m{n<0} use
\m{u^n-1=-u^n(u^{-n}-1)}; for \m{n=0} both sides are infinite. The second formula follows from
\m{u^{\ell^jn}-u^{\ell^jn'}=u^{\ell^jn'}(u^{\ell^j(n-n')}-1)}.
\end{proof}
\leanline{\lean{lte}\sep\lean{lteNorm}, with \lean{normDvd} and \lean{normAddVal}}

\begin{lemma}\label{l:units}
Let $t$ and~$t'$ be units of~$\Z_\ell$ and \m{z\in\overline{\Q}_\ell}. Then\quad \m{|t-t'|\le\max\bigl(|tz-1|,\,|t'z-1|\bigr).}
\end{lemma}
\begin{proof}
We have \m{(tz-1)-(t'z-1)=(t-t')z}. If \m{|z|\ge1}, then \m{|t-t'|\le|(t-t')z|}, and the ultrametric inequality
gives the claim. If \m{|z|<1}, then \m{|tz|<1}, so \m{|tz-1|=1\ge|t-t'|}.
\end{proof}
\leanline{\lean{units}}

\begin{lemma}\label{l:helly}
For \m{n=0,1,2,\dots} let \m{b_n\in\Z_\ell} and a real number \m{r_n\ge0} be given, such that
\[
|b_n-b_{n'}|\le\max(r_n,r_{n'})\qquad\text{for all }n,n'.
\]
Then there is \m{c\in\Z_\ell} with \m{|c-b_n|\le r_n} for all~$n$.
\end{lemma}
\begin{proof}
The closed balls \m{\{y:|y-b_n|\le r_n\}} are closed subsets of the compact space~$\Z_\ell$. Any finitely many of
them have a common point, namely the centre $b_n$ of one of smallest radius among them. Hence all of them have a
common point.
\end{proof}
\leanline{\lean{helly}}

\begin{proposition}[one centre]\label{p:ball}
Let \m{u\in\Z_\ell} have level~$a$, let \m{j\ge0} be an integer and \m{z\in\overline{\Q}_\ell}. There is
\m{c\in\Z_\ell} such that for every integer \m{h\ge0}
\[
\ell^{-(j+a)}\,|c-h|\ \le\ \bigl|u^{\ell^jh}z-1\bigr|.
\]
\end{proposition}
\begin{proof}
Put \m{b_h=h} and \m{r_h=\ell^{j+a}\,|u^{\ell^jh}z-1|}. By Lemma~\ref{l:lte} and Lemma~\ref{l:units}, applied to
the units $u^{\ell^jh}$ and $u^{\ell^jh'}$,
\[
|h-h'|=\ell^{j+a}\,\bigl|u^{\ell^jh}-u^{\ell^jh'}\bigr|\le\max(r_h,r_{h'}).
\]
Lemma~\ref{l:helly} gives~$c$.
\end{proof}
\leanline{\lean{ball}}

The proposition says that the function \m{h\mapsto u^{\ell^jh}z-1}, on the integers, is bounded below by a
multiple of the distance to one point. The function extends continuously to~$\Z_\ell$. If the extension has a
zero~$c$, this zero is the only possible centre, and the inequality is an equality. In general there is no zero
in~$\Z_\ell$; then the function is bounded away from zero, and the centre is not unique.

\begin{lemma}\label{l:eigen}
Let $S$ be a square matrix of size $R$ over an algebraically closed field, with eigenvalues \m{z_1,\dots,z_R}
counted with multiplicity. Then \m{\det(tS-1)=\prod_{i=1}^R(tz_i-1)} for every element $t$ of the field.
\end{lemma}
\begin{proof}
$S$ is conjugate to an upper triangular matrix with diagonal \m{z_1,\dots,z_R}, and then $tS-1$ is conjugate to an
upper triangular matrix with diagonal \m{tz_1-1,\dots,tz_R-1}.
\end{proof}
\leanline{\lean{detFactor}}

\begin{proposition}\label{p:det}
Let \m{u\in\Z_\ell} have level~$a$, let \m{j\ge0}, and let $S$ be a square matrix of size $R$ over~$\Z_\ell$. There are
\m{c_1,\dots,c_R\in\Z_\ell} such that for every integer \m{h\ge0}
\[
\vlB{\ell}{\det\bigl(u^{\ell^jh}S-1\bigr)}\ \le\ \sum_{i=1}^{R}\bigl(j+a+\vl{\ell}{c_i-h}\bigr).
\]
\end{proposition}
\begin{proof}
Let \m{z_1,\dots,z_R\in\overline{\Q}_\ell} be the eigenvalues of~$S$ and let $c_i$ be a centre for $z_i$ as in
Proposition~\ref{p:ball}. By Lemma~\ref{l:eigen},
\[
\Bigl|\det\bigl(u^{\ell^jh}S-1\bigr)\Bigr|=\prod_{i=1}^{R}\bigl|u^{\ell^jh}z_i-1\bigr|\ \ge\
\prod_{i=1}^{R}\ell^{-(j+a)}\,|c_i-h| ,
\]
and taking valuations gives the claim.
\end{proof}
\leanline{\lean{detBound}\sep\lean{normVal}, with \lean{dvdNorm} and \lean{vlProd}}

\begin{corollary}[centres at one layer]\label{c:level}
Let $X$ be a module generated by $R$ elements, let $\gamma$ be any endomorphism of~$X$, let \m{u\in\Z_\ell}
have level~$a$ and let \m{j\ge0}. There are \m{c_1,\dots,c_R\in\Z_\ell} such that for every integer \m{h\ge0}
\[
\tw{\gamma}{u}{h}{j}\ \le\ \sum_{i=1}^{R}\bigl(j+a+\vl{\ell}{c_i-h}\bigr).
\]
\end{corollary}
\begin{proof}
Since \m{(u^h\gamma)^{\ell^j}=u^{\ell^jh}\gamma^{\ell^j}}, Proposition~\ref{p:lift} for \m{g=\gamma^{\ell^j}}
and \m{t=u^{\ell^jh}} bounds $\tw{\gamma}{u}{h}{j}$ by \m{\vl{\ell}{\det(u^{\ell^jh}S-1)}}, with a matrix $S$ that does not depend
on~$h$. Apply Proposition~\ref{p:det}.
\end{proof}
\leanline{\lean{level}}

Unipotence is not used here. But the matrix $S$, and with it the centres, depend on the layer~$j$. To have one
set of centres for all layers, the layers must be compared. This is the purpose of the next section.

% ------------------------------------------------------------------------------------- 5
\sect{Growth along the tower}{s:growth}
\begin{lemma}\label{l:small}
\begin{enumerate}[label=\textup{(\alph*)},topsep=2pt,itemsep=2pt]
\item If the module $X$ is generated by $R$ elements, then \m{\len X/\ell X\le R}; in the notation of
Definition~\ref{d:cok} this is \m{\len\cok(\ell\cdot1)\le R}.
\item Let $M$ be a module over a commutative ring, of length at most~$n$ for an integer \m{n\ge0}. A nilpotent
endomorphism $f$ of~$M$ satisfies \m{f^n=0}.
\end{enumerate}
\end{lemma}
\begin{proof}
(a) $X/\ell X$ is a quotient of \m{(\Z_\ell/\ell\Z_\ell)^R}. (b) As long as \m{f^i\ne0}, the image of $f^{i+1}$ is
strictly smaller than the image of~$f^i$: equality would give \m{f^i(M)=f^{i+j}(M)} for all~$j$, hence
\m{f^i=0}. If \m{f^n\ne0}, the images of \m{f^0,f^1,\dots,f^n} and the zero module would form a chain of length
\m{n+1}.
\end{proof}
\leanline{\lean{cokL}\sep\lean{nilLen}, with \lean{nilSucc} and \lean{nilLt}}

\begin{lemma}\label{l:unip}
Let $X$ be a module generated by $R$ elements and let \m{\gamma\in\End(X)} be unipotent modulo~$\ell$.
\begin{enumerate}[label=\textup{(\alph*)},topsep=2pt,itemsep=2pt]
\item \m{(\gamma-1)^R\equiv0\pmod\ell}.
\item For every \m{t\in\Z_\ell} with \m{t\equiv1\pmod\ell}, the endomorphism $t\gamma$ is unipotent modulo~$\ell$.
\end{enumerate}
\end{lemma}
\begin{proof}
(a) The endomorphism of $X/\ell X$ induced by \m{\gamma-1} is nilpotent, and $X/\ell X$ has length at most~$R$;
apply Lemma~\ref{l:small}. (b) Write \m{t\gamma-1=(\gamma-1)+(t-1)\gamma}. The two summands commute and the second
lies in~$I_1$. In the binomial expansion of \m{(t\gamma-1)^k} every term except \m{(\gamma-1)^k} contains the second
summand, so it lies in the ideal~$I_1$; and \m{(\gamma-1)^k\in I_1} for a suitable~$k$.
\end{proof}
\leanline{\lean{unipPow}\sep\lean{twist}}

\begin{lemma}[congruences]\label{l:cong}
Let $X$ be a module, let \m{g,w\in\End(X)} and let \m{k,n\ge0} be integers.
\begin{enumerate}[label=\textup{(\alph*)},topsep=2pt,itemsep=2pt]
\item \m{(g-1)^{\ell^k}\equiv g^{\ell^k}-1\pmod\ell}.
\item If \m{w\equiv1\pmod{\ell^n}}, then \m{w^i\equiv1\pmod{\ell^n}} for every \m{i\ge0}.
\item If \m{w\equiv1\pmod\ell}, then \m{w^\ell\equiv1\pmod{\ell^2}}.
\end{enumerate}
\end{lemma}
\begin{proof}
(a) The polynomial \m{(T-1)^{\ell^k}-(T^{\ell^k}-1)} has coefficients divisible by~$\ell$. (b) We have
\m{w^i-1=(1+w+\dots+w^{i-1})(w-1)}, and $I_n$ is an ideal. (c) Put \m{f=w-1\in I_1}. In
\m{w^\ell-1=\sum_{i=1}^{\ell}\binom{\ell}{i}f^i} the terms with \m{i<\ell} lie in \m{\ell I_1\subseteq I_2}, and
\m{f^\ell\in I_\ell\subseteq I_2}.
\end{proof}
\leanline{\lean{frob}\sep\lean{powIn}}

\begin{lemma}[one step]\label{l:step}
Let $X$ be a module generated by $R$ elements and let \m{w\in\End(X)} with \m{w\equiv1\pmod{\ell^2}}. Put
\m{\nu=1+w+\dots+w^{\ell-1}}. Then \m{\ell X\subseteq\nu(X)}, and
\[
\len\cok(w^\ell-1)\ \le\ \len\cok(w-1)+R.
\]
\end{lemma}
\begin{proof}
We have \m{w^\ell-1=\nu\,(w-1)}. By Lemma~\ref{l:cong}(b) each $w^i$ is
congruent to $1$ modulo~$\ell^2$, so \m{\nu\equiv\ell\pmod{\ell^2}}: for every \m{x\in X} there is $y$ with
\m{\ell x=\nu x-\ell^2y}. Thus \m{\ell X\subseteq\nu(X)+\ell\cdot\ell X}. The module $\ell X$ is finitely
generated, so Nakayama's lemma gives \m{\ell X\subseteq\nu(X)}. By Lemma~\ref{l:len}(a) and
Lemma~\ref{l:small}(a), \m{\len\cok\nu\le\len X/\ell X\le R}, and Lemma~\ref{l:len}(b) finishes the proof.
\end{proof}
\leanline{the inclusion \lean{phi}, there for every finitely generated module\sep the inequality \lean{step}}

\begin{proposition}[growth]\label{p:growth}
Let $X$ be a module generated by $R$ elements and let \m{g\in\End(X)} be unipotent modulo~$\ell$. Then for
every integer \m{m\ge0}
\[
\len\cok\bigl(g^{\ell^m}-1\bigr)\ \le\ \len\cok\bigl(g^{\ell^{\mz{\ell}{R}}}-1\bigr)+R\cdot\max(m-\mz{\ell}{R},0).
\]
\end{proposition}
\begin{proof}
For \m{m\le \mz{\ell}{R}} the exponent $\ell^m$ divides $\ell^{\mz{\ell}{R}}$, and the claim is
Lemma~\ref{l:len}(c).

Let \m{k=\mz{\ell}{R}-2}, so that \m{R\le\ell^k}. Then, modulo the indicated powers of~$\ell$,
\[
(g-1)^{\ell^k}\equiv0\ (\ell),\qquad g^{\ell^k}\equiv1\ (\ell),\qquad g^{\ell^{k+1}}\equiv1\ (\ell^2),\qquad
g^{\ell^n}\equiv1\ (\ell^2)\quad\text{for all }n\ge k+1,
\]
by Lemma~\ref{l:unip}(a) and parts (a), (c), (b) of Lemma~\ref{l:cong}, in this order. For such $n$, Lemma~\ref{l:step} with \m{w=g^{\ell^n}} gives
\[
\len\cok\bigl(g^{\ell^{n+1}}-1\bigr)\ \le\ \len\cok\bigl(g^{\ell^n}-1\bigr)+R ,
\]
and the claim for \m{m>\mz{\ell}{R}} follows by induction, starting from \m{n=\mz{\ell}{R}}.
\end{proof}
\leanline{\lean{growth}, with \lean{mono} for the layers below the reference layer}

\begin{proof}[Proof of Theorems A and B]
Since \m{u\equiv1\pmod\ell}, also \m{u^h\equiv1\pmod\ell}, so \m{u^h\gamma} is unipotent modulo~$\ell$ by
Lemma~\ref{l:unip}(b). Proposition~\ref{p:growth} for \m{g=u^h\gamma} is Theorem~B. For Theorem~A take the centres of
Corollary~\ref{c:level} at the layer \m{j=\mz{\ell}{R}}; then by Theorem~B
\[
\begin{aligned}
\tw{\gamma}{u}{h}{m}\ &\le\ \sum_{i=1}^{R}\bigl(\mz{\ell}{R}+a+\vl{\ell}{c_i-h}\bigr)+R\cdot\max\bigl(m-\mz{\ell}{R},0\bigr)\\
&=\ \sum_{i=1}^{R}\bigl(\max\bigl(m,\mz{\ell}{R}\bigr)+a+\vl{\ell}{c_i-h}\bigr).
\end{aligned}\qedhere
\]
\end{proof}

% ------------------------------------------------------------------------------------- 6
\sect{Exceptional twists}{s:exc}
\begin{lemma}[separation]\label{l:sep}
Let $K$ be a field with an ultrametric absolute value \m{|\cdot|} that extends the $\ell$-adic absolute value
of~$\Q_\ell$, let \m{c\in K}, and let \m{Z\ge1} be an integer. If integers $h$ and~$h'$ with \m{1\le h\le Z} and
\m{1\le h'\le Z} satisfy
\[
|c-h|<\ell^{-\lfloor\log_\ell Z\rfloor}\qquad\text{and}\qquad|c-h'|<\ell^{-\lfloor\log_\ell Z\rfloor},
\]
then \m{h=h'}. In particular, for \m{c\in\Z_\ell} at most one integer $h$ with \m{1\le h\le Z} satisfies
\m{\vl{\ell}{c-h}>\lfloor\log_\ell Z\rfloor}.
\end{lemma}
\begin{proof}
By the ultrametric inequality \m{|h'-h|\le\max(|c-h|,|c-h'|)<\ell^{-\lfloor\log_\ell Z\rfloor}}. The absolute value
of~$K$ restricts to the $\ell$-adic one on the integers, so \m{\vl{\ell}{h'-h}>\lfloor\log_\ell Z\rfloor}: the
integer \m{h'-h} is divisible by \m{\ell^{\lfloor\log_\ell Z\rfloor+1}}, which is larger than~$Z$. But \m{h'-h} lies
strictly between $-Z$ and~$Z$. So \m{h'-h=0}. The last sentence is the case \m{K=\Q_\ell}, where
\m{|c-h|=\ell^{-\vl{\ell}{c-h}}}.
\end{proof}
\leanline{\lean{near1}}

\begin{proof}[Proof of Corollary C]
By Lemma~\ref{l:sep} each centre contributes at most one element to~$\mathcal E_Z$. If \m{1\le h\le Z} is not
in~$\mathcal E_Z$, then \m{\vl{\ell}{c_i-h}\le\lfloor\log_\ell Z\rfloor} for all~$i$, and Theorem~A gives the bound.
\end{proof}

In the archive the count of the exceptional twists is stated for the arithmetic lengths of Theorem~D, not for
abstract modules, and the centres there lie in a field~$K$ as in Lemma~\ref{l:sep}; this form is
Corollary~\ref{c:exc} of §\ref{s:arith}, and the argument is the one above.

% ------------------------------------------------------------------------------------- 7
\sect{Class groups in the cyclotomic tower}{s:arith}
This section defines the objects of Theorem~D, each with all the arguments of its formal definition, states three
inputs and proves the theorem. The definitions are made for any number field, as in the archive; the theorem uses
them for the field of Setting~\ref{set:tower}.

\subhead{Galois groups.} Let $F$ be a number field and $\overline F$ an algebraic closure of~$F$. The group
\m{\GalG{F}{\overline F}} of the automorphisms of~$\overline F$ over~$F$ carries the Krull topology; the same
notation is used for any Galois extension of~$F$ in the place of~$\overline F$. For a subfield~$E$
of~$\overline F$, which need not contain~$F$, we write \m{\galOf{F}{E}{\overline F}} for the subgroup of the
elements of \m{\GalG{F}{\overline F}} that fix every element of~$E$. The map
\[
\cyc{F}{\overline F}{\ell}{\cdot}\colon\ \GalG{F}{\overline F}\longrightarrow\Z_\ell^\times
\]
is the $\ell$-adic cyclotomic character: \m{\sigma(\zeta)=\zeta^{u}} with \m{u=\cyc{F}{\overline F}{\ell}{\sigma}},
for every root of unity \m{\zeta\in\overline F} of $\ell$-power order. The field, its closure and the prime are part
of the notation and are written every time; for a subgroup~$V$ we write \m{\cyc{F}{\overline F}{\ell}{V}} for its
image. For an integer \m{n\ge0} let
\[
\princ{\ell}{n}=\bigl\{u\in\Z_\ell^\times:\ u\equiv1\pmod{\ell^n}\bigr\},
\]
so that \m{\princ{\ell}{0}=\Z_\ell^\times} and \m{\princ{\ell}{n}=1+\ell^n\Z_\ell} for \m{n\ge1}. These are closed
subgroups of~$\Z_\ell^\times$. For a subgroup $U$ of \m{\GalG{F}{\overline F}} put
\[
\levU{\ell}{U}{n}=\bigl\{\sigma\in U:\ \cyc{F}{\overline F}{\ell}{\sigma}\in \princ{\ell}{n}\bigr\},\qquad
\cycU{\ell}{U}=\bigl\{\sigma\in U:\ \cyc{F}{\overline F}{\ell}{\sigma}=1\bigr\}.
\]
For the whole group in the place of~$U$ the first of them is also written \m{\Un{F}{\overline F}{\ell}{n}}.

\begin{definition}[topological generators]\label{d:isgen}
Let $F$ be a number field with an algebraic closure~$\overline F$, let $U$ be a subgroup of
\m{\GalG{F}{\overline F}} and \m{\sigma\in\GalG{F}{\overline F}}. We say that $\sigma$ is a \kw{topological generator
of~$U$ for the cyclotomic character of~$\ell$} if \m{\sigma\in U} and the image
\m{\cyc{F}{\overline F}{\ell}{U}} lies in the closure of the set of the powers
\m{\cyc{F}{\overline F}{\ell}{\sigma}^n}, \m{n\in\Z}.
\end{definition}
\leanline{\lean{IsGen}}

\begin{definition}[the exponent]\label{d:czero}
For a field~$L$, the number \m{\cZero{\ell}{L}} is the largest integer \m{n\ge0} such that $L$ contains a
primitive root of unity of order~$\ell^n$. A number field contains finitely many roots of unity, so for it such an
integer exists. (The archive defines \m{\cZero{\ell}{L}} as a supremum, for any natural number in place of the prime
and any field, and the supremum is read as~$0$ when the orders are not bounded.)
\end{definition}
\leanline{\lean{cZero}}

\begin{definition}[characters with local conditions]\label{d:dinf}
Let $F$ be a number field with an algebraic closure~$\overline F$, let \m{N\ge0} be an integer, $U$ a subgroup of
\m{\GalG{F}{\overline F}}, $\Lambda$ a commutative ring and $M$ a module over~$\Lambda$. A \kw{valuation ring}
of~$\overline F$ is a subring $\mathcal O$ such that \m{y\in\mathcal O} or \m{y^{-1}\in\mathcal O} for every
non-zero \m{y\in\overline F}. Its \kw{decomposition group} consists of the \m{\sigma\in\GalG{F}{\overline F}} with
\m{\sigma(\mathcal O)=\mathcal O}, and its \kw{inertia group} of those among them for which \m{\sigma y-y} lies in
the maximal ideal of~$\mathcal O$ for every \m{y\in\mathcal O}. We write
\[
\XDual{N}{U}{\Lambda}{M}
\]
for the $\Lambda$-module of the maps \m{\varphi\colon U\to M} with the following four properties.
\begin{enumerate}[label=\textup{(\arabic*)},topsep=2pt,itemsep=1pt]
\item $\varphi$ is locally constant.
\item \m{\varphi(\rho\rho')=\varphi(\rho)+\varphi(\rho')} for all \m{\rho,\rho'\in U}.
\item $\varphi$ vanishes on the intersection of~$U$ with the inertia group of every valuation ring of~$\overline F$.
\item $\varphi$ vanishes on the intersection of~$U$ with the decomposition group of every valuation ring
of~$\overline F$ whose maximal ideal contains~$N$.
\end{enumerate}
For the prime~$\ell$ we put
\[
\Dinf{\ell}{N}{U}=\XDual{N}{\cycU{\ell}{U}}{\Z_\ell}{\QZl{\ell}},\qquad
\Xinf{\ell}{N}{U}=\Hom_{\Z_\ell}\bigl(\Dinf{\ell}{N}{U},\,\QZl{\ell}\bigr).
\]
Each of them is a function of three arguments, and all three are written every time. If $U$ is a normal subgroup,
\m{\GalG{F}{\overline F}} acts on both modules:
\[
(\sigma\varphi)(\rho)=\varphi(\sigma^{-1}\rho\sigma),\qquad(\sigma x)(\varphi)=x(\sigma^{-1}\varphi)
\qquad\text{for }\varphi\in \Dinf{\ell}{N}{U},\ x\in \Xinf{\ell}{N}{U}.
\]
\end{definition}
\leanline{\lean{XDual}}

The \kw{Iwasawa module} of Theorem~D is \m{\Xinf{\ell}{\ell}{\GalG{L}{\overline L}}}, for the field~$L$ of
Setting~\ref{set:tower}: here \m{N=\ell} and $U$ is the whole group. For this~$U$ the group $\cycU{\ell}{U}$ is
\m{\galOf{L}{L_\infty}{\overline L}}, and \m{\Dinf{\ell}{\ell}{\GalG{L}{\overline L}}} is the group of the continuous
homomorphisms \m{\galOf{L}{L_\infty}{\overline L}\to\QZl{\ell}} that vanish on every inertia group and on the
decomposition groups of the places above~$\ell$. By Pontryagin duality its dual is the Galois group of the maximal
abelian pro-$\ell$ extension of~$L_\infty$ that is unramified everywhere and in which the primes above~$\ell$ split
completely, as in §\ref{s:results}; the archive works with the definition as a dual, which needs no topology on the
module. With \m{N=1} no decomposition group enters, and \m{\Xinf{\ell}{1}{\GalG{L}{\overline L}}} is the Galois group
of the maximal unramified abelian pro-$\ell$ extension of~$L_\infty$.

\begin{setting}\label{set:tower}
$L$ is a number field and $\overline L$ is an algebraic closure of~$L$. The field $L$ contains the $\ell$-th roots
of unity, and the fourth roots of unity if \m{\ell=2}. Such a field is totally complex (Lemma~\ref{l:roots}(d));
the archive lists this among its hypotheses. The field \m{L_\infty\subseteq\overline L} is generated
over~$L$ by all roots of unity of $\ell$-power order, and for an integer \m{m\ge0} the field $L_m$ is the subfield
of~$L_\infty$ of degree $\ell^m$ over~$L$.
\end{setting}

In Setting~\ref{set:tower} the image of \m{\GalG{L}{\overline L}} under the cyclotomic character is the group
$\princ{\ell}{a}$ with \m{a=\cZero{\ell}{L}}, a group isomorphic to~$\Z_\ell$, and the field $L_m$ exists and is unique
(Lemma~\ref{l:layers}).

\begin{definition}[layers]\label{d:layer}
Let \m{j\ge0} be an integer and let $L'$ be an extension of a field~$L$. The field $L'$ is a \kw{layer of
index~$j$ over~$L$ for the prime~$\ell$} if, for some integer \m{k\ge0}, it contains a primitive root of unity of
order~$\ell^k$ and is generated over~$L$ by the roots of unity of order dividing~$\ell^k$, and if the degree of
$L'$ over~$L$ is~$\ell^j$.
\end{definition}
\leanline{\lean{IsLayer}}

In Setting~\ref{set:tower} the field $L_m$ is a layer of index~$m$ over~$L$ for~$\ell$, and a number field that is a
layer of index~$m$ over~$L$ for~$\ell$ is isomorphic to~$L_m$ over~$L$ (Lemma~\ref{l:layers}). The statements of
this note are made for the fields~$L_m$; the archive makes them for every such layer.

\begin{definition}[the classes and their length]\label{d:sha}
Let $F$ be a number field, $\overline F$ an algebraic closure of~$F$ and $h$ an integer. The Tate twist
\m{\qzG{F}{\overline F}{\ell}{h}} is the group $\QZl{\ell}$ on which \m{\sigma\in\GalG{F}{\overline F}} acts as the
multiplication by \m{\cyc{F}{\overline F}{\ell}{\sigma}^{-h}}; it is a function of the field, its closure, the
prime and the twist. A valuation ring of~$\overline F$ lies
\kw{above} a prime~$\mathfrak p$ of~$F$ if its maximal ideal meets the ring of integers of~$F$ in~$\mathfrak p$.
The group $\Sha^1_\ell(F,h)$ consists of the classes in
\m{H^1\bigl(\GalG{F}{\overline F},\qzG{F}{\overline F}{\ell}{h}\bigr)} whose restriction to the decomposition group of every
valuation ring above a prime of~$F$ that divides~$\ell$ is trivial, and whose restriction to the inertia group of
every valuation ring above a prime of~$F$ that does not divide~$\ell$ is trivial. Cohomology is that of continuous
cochains, for the discrete topology on the coefficients. The length of this group is written
\[
\shaInv{F}{\ell}{h}\ \in\ \{0,1,2,\dots,\infty\};
\]
it is a function of the field, the prime and the twist.
\end{definition}
\leanline{\lean{shaInv}}

\begin{definition}[local data]\label{d:ldata}
Let $F$ be a number field. A \kw{local datum of~$F$ at~$\ell$} is a pair $\mathcal D$ of an algebraic closure
$\overline F$ of~$F$ and of a choice, for every prime~$\mathfrak p$ of~$F$ that divides~$\ell$, of one valuation
ring of~$\overline F$ above~$\mathfrak p$. Let $\Gamma_{\mathcal D}$ be the quotient of \m{\GalG{F}{\overline F}} by the
closed normal subgroup generated by the inertia groups of the valuation rings above the primes of~$F$ that do not
divide~$\ell$: the Galois group of the maximal extension of~$F$ in~$\overline F$ that is unramified outside~$\ell$
and the infinite places. The cyclotomic character \m{\cyc{F}{\overline F}{\ell}{\cdot}} factors
through~$\Gamma_{\mathcal D}$. For an integer~$h$ let $\Z_\ell(h)$ be the group $\Z_\ell$ on which $\Gamma_{\mathcal D}$ acts
through the $h$-th power of this character, and let \m{\qzTw{\mathcal D}{h}} be the group $\QZl{\ell}$ on
which $\Gamma_{\mathcal D}$ acts through its $(-h)$-th power, as in Definition~\ref{d:sha}. We write \m{H^i(\mathcal D,\cdot)} for the continuous cohomology of~$\Gamma_{\mathcal D}$, and
\[
\ShaOf{\mathcal D}{h}=\ker\Bigl(H^1\bigl(\mathcal D,\qzTw{\mathcal D}{h}\bigr)\longrightarrow
\prod_{\mathfrak p\mid\ell}H^1\bigl(F_{\mathfrak p},\qzTw{\mathcal D}{h}\bigr)\Bigr),
\]
where the cohomology of the completion $F_{\mathfrak p}$ is that of the decomposition group of the valuation ring
chosen above~$\mathfrak p$. The length of the second cohomology group is written
\[
\VLen{\mathcal D}{h}\ \in\ \{0,1,2,\dots,\infty\};
\]
it is a function of the datum and of the twist. The two descriptions of the classes have the same length: for every
local datum $\mathcal D$ of~$F$ at~$\ell$ and every integer~$h$,
\[
\len\ShaOf{\mathcal D}{h}=\shaInv{F}{\ell}{h}.
\]
This is a theorem of the archive, proved in a neighbouring directory by inflation; it is not deformalized here.
\end{definition}
\leanline{\lean{VLen}\sep the equality of the two lengths \lean{shaLenInv}}

\begin{setting}[three hypotheses on Galois cohomology]\label{set:coh}
The following three statements hold in the form in which the archive states them, for every prime number in the
place of~$\ell$. They are hypotheses of the cluster. They are not deformalized in this note.
\begin{enumerate}[label=\textup{(\alph*)},topsep=2pt,itemsep=1pt]
\item \emph{The duality of Poitou and Tate} \cite[VIII, §6]{NSW}, for a number field that is totally complex if the
prime is~$2$, a finite set of its primes that contains those above the prime, and a finite module of order a power
of the prime over the Galois group of the maximal extension unramified outside this set: the exact sequence of nine
terms, the duality of the locally trivial classes in degrees $1$ and~$2$, and the vanishing of the cohomology in
degrees at least~$3$.
\item \emph{Finiteness of local cohomology}: for a finite extension of the field of $\ell$-adic numbers and a finite
Galois module of order a power of~$\ell$, the cohomology is finite in degrees $0$, $1$,~$2$ and vanishes in degrees
at least~$3$.
\item \emph{Local Tate duality} for the same modules: the cup products with values in the second cohomology of the
roots of unity are perfect pairings.
\end{enumerate}
\end{setting}

\subhead{The algebraic bridge.} The module \m{\Dinf{\ell}{N}{U}} is discrete and \m{\Xinf{\ell}{N}{U}} is
compact. The next statement passes unipotence from one side to the other. It is pure algebra, and is proved in full
in the archive's directory.

\begin{proposition}[the dual module]\label{p:dual}
Let $D$ be a module and \m{g\in\End(D)}. Put
\[
X=\Hom_{\Z_\ell}(D,\QZl{\ell}),\qquad (g^\vee x)(\varphi)=x(g\varphi)\quad\text{for }x\in X,\ \varphi\in D ,
\]
and let $D[\ell]$ be the submodule of the elements of~$D$ killed by~$\ell$.
\begin{enumerate}[label=\textup{(\alph*)},topsep=2pt,itemsep=2pt]
\item The module $\QZl{\ell}$ is injective.
\item If $X$ is finitely generated, then $D[\ell]$ is finite.
\item If \m{x\in X} vanishes on~$D[\ell]$, that is \m{x(\varphi)=0} for all \m{\varphi\in D} with
\m{\ell\varphi=0}, then \m{x\in\ell X}.
\item Suppose that $X$ is finitely generated and that \m{g-1} is locally nilpotent: for every \m{\varphi\in D}
there is an integer \m{k\ge0} with \m{(g-1)^k\varphi=0}. Then $g^\vee$ is unipotent modulo~$\ell$.
\end{enumerate}
\end{proposition}
\begin{proof}
(a) The module $\QZl{\ell}$ is divisible, hence injective, since $\Z_\ell$ is a principal ideal domain.

(b) Let \m{x_1,\dots,x_r} generate~$X$. A non-zero \m{\varphi\in D[\ell]} spans a module
isomorphic to $\Z_\ell/\ell\Z_\ell$, which embeds into~$\QZl{\ell}$; by injectivity the embedding extends to an element
of~$X$ that does not vanish at~$\varphi$. Hence some $x_i$ does not vanish at~$\varphi$, and
\m{\varphi\mapsto(x_1(\varphi),\dots,x_r(\varphi))} embeds $D[\ell]$ into the finite group of $r$-tuples of elements
of $\QZl{\ell}$ killed by~$\ell$.

(c) If \m{x\in X} vanishes on $D[\ell]$, then
\m{\ell\varphi\mapsto x(\varphi)} is a well-defined homomorphism \m{\ell D\to\QZl{\ell}}. By injectivity it extends to some
\m{y\in X}, and \m{x=\ell y}.

(d) By (b) the module $D[\ell]$ is finite; choose \m{k\ge1} with \m{(g-1)^k=0} on it. For \m{x\in X} and
\m{\varphi\in D[\ell]},
\m{((g^\vee-1)^kx)(\varphi)=x((g-1)^k\varphi)=0}. So \m{(g^\vee-1)^kx\in\ell X} by~(c).
\end{proof}
\leanline{(a) \lean{qzInj}\sep (b) \lean{torsFin}\sep (c) \lean{dualDiv}\sep (d) \lean{dualUnip}}

\subhead{Three inputs.} Inputs 1 and~3 are proved in neighbouring directories of the archive, or assumed as
hypotheses of the cluster. We state them and indicate the standard arguments; they are not deformalized here.
Input~2 is proved in §\ref{s:descent}. Inputs 1 and~2 are stated for the modules of Theorem~D; the archive proves
Input~1 for every \m{\Dinf{\ell}{N}{U}} with $U$ closed and normal.

\statement{inksoft}{\phantomsection\label{i:1}\textbf{Input 1} (local nilpotence).\enspace\emph{In
Setting~\ref{set:tower}, write \m{G=\GalG{L}{\overline L}}, and let \m{\sigma\in G} satisfy
\m{\cyc{L}{\overline L}{\ell}{\sigma}\equiv1\pmod\ell}. Then \m{\sigma-1} is locally nilpotent on
\m{\Dinf{\ell}{\ell}{G}}: for every \m{\varphi\in \Dinf{\ell}{\ell}{G}} there is an integer \m{k\ge0} with
\m{(\sigma-1)^k\varphi=0}.}

\vspace{3pt}\phantomsection\label{i:2}\textbf{Input 2} (descent).\enspace\emph{In Setting~\ref{set:tower}, write
\m{G=\GalG{L}{\overline L}}, let \m{\sigma\in G} be a topological generator of~$G$ for the cyclotomic character
of~$\ell$ (Definition~\ref{d:isgen}), and put \m{u=\cyc{L}{\overline L}{\ell}{\sigma}}. For all integers
\m{h\ne0} and \m{m\ge0}, and for the layer $L_m$ of index~$m$ over~$L$,}
\[
\shaInv{L_m}{\ell}{h}=\len\cok\bigl((u^h\sigma)^{\ell^m}-1\bigr),
\]
\emph{where \m{u^h\sigma} is the operator on \m{\Xinf{\ell}{\ell}{G}} by which $\sigma$ acts on the twisted module
\m{\xTw{\ell}{\ell}{G}{h}} of Definition~\ref{d:tw}. For \m{h\ge1} the right-hand side is the twisted length
\m{\tw{\gamma}{u}{h}{m}} of the pair $(\gamma,u)$, where $\gamma$ is the action of~$\sigma$
on~\m{\Xinf{\ell}{\ell}{G}}.}

\vspace{3pt}\phantomsection\label{i:3}\textbf{Input 3} (Poitou--Tate).\enspace\emph{Assume the three hypotheses of
Setting~\ref{set:coh}. Let $F$ be a number field, and if \m{\ell=2} let $F$ have no real place. Then for every local
datum $\mathcal D$ of~$F$ at~$\ell$, formed in an algebraic closure~$\overline F$, and every integer~$h$,}
\[
\len\ShaOf{\mathcal D}{h}\ \le\ \VLen{\mathcal D}{h}.
\]}
\leanline{1 \lean{dinfNil}\sep 2 \lean{step0}; for \m{h\ge1} the statement is \lean{Descends}, proved in
\lean{descends}, and the last sentence is \lean{eqTwLen}\sep 3 \lean{shaLeV}}

\emph{On Input 1.} Write \m{G=\GalG{L}{\overline L}}. An element $\varphi$ of~\m{\Dinf{\ell}{\ell}{G}} has order a
power of~$\ell$, say~$\ell^k$, and is fixed by an open subgroup of~$G$ containing \m{\galOf{L}{L_\infty}{\overline L}},
hence by $\sigma^{\ell^n}$ for some~$n$. As in the proof of Lemma~\ref{l:cong}(a),
\m{(T-1)^{\ell^n}-(T^{\ell^n}-1)} is $\ell$ times a polynomial with integer coefficients, so
\m{(\sigma-1)^{\ell^n}\varphi} is $\ell$ times an element of the module generated by the conjugates
$\sigma^i\varphi$, each of which is again fixed by~$\sigma^{\ell^n}$; repeating this $k$ times gives
\m{(\sigma-1)^{k\ell^n}\varphi=0}.

\emph{On Input 2.} Write \m{G=\GalG{L}{\overline L}} again. The length of a group of cohomology classes over the
layer~$L_m$ is expressed through one operator on the module~\m{\Xinf{\ell}{\ell}{G}}, which belongs to the top of
the tower. The classes over~$L_m$ are
first identified with the elements of \m{\Dinf{\ell}{\ell}{G}} that are invariant under a twisted action of
\m{\galOf{L}{L_m}{\overline L}}; this is the cohomological step. The invariants are the kernel of one operator
on~\m{\Dinf{\ell}{\ell}{G}}, and duality turns the kernel into a cokernel on~\m{\Xinf{\ell}{\ell}{G}}. The proof
is given in §\ref{s:descent}.

\emph{On Input 3.} By the duality theorem of Poitou and Tate \cite[VIII, §6]{NSW}, $\ShaOf{\mathcal D}{h}$ is dual to
the corresponding subgroup of \m{H^2(\mathcal D,\Z_\ell(h+1))}; the coefficients are dual to each other:
the homomorphisms from \m{\qzTw{\mathcal D}{h}} to the group of the roots of unity of $\ell$-power order form the
module \m{\Z_\ell(h+1)}. In the archive the statement is deduced from the three hypotheses of Setting~\ref{set:coh},
which the cluster passes on unchanged.

\begin{corollary}[unipotence on the Iwasawa module]\label{c:unip}
In Setting~\ref{set:tower}, write \m{G=\GalG{L}{\overline L}} and suppose that \m{\Xinf{\ell}{\ell}{G}} is finitely
generated. If \m{\sigma\in G} satisfies \m{\cyc{L}{\overline L}{\ell}{\sigma}\equiv1\pmod\ell}, that is, if
\m{\sigma\in \levU{\ell}{G}{1}}, then the action of~$\sigma$ on \m{\Xinf{\ell}{\ell}{G}} is unipotent modulo~$\ell$.
\end{corollary}
\begin{proof}
The element $\sigma^{-1}$ satisfies the same congruence, so by Input~1 the operator \m{\sigma^{-1}-1} is locally
nilpotent on \m{\Dinf{\ell}{\ell}{G}}. By the formula \m{(\sigma x)(\varphi)=x(\sigma^{-1}\varphi)} of
Definition~\ref{d:dinf}, the action of~$\sigma$ on \m{\Xinf{\ell}{\ell}{G}} is~$g^\vee$, where $g$ is the action
of~$\sigma^{-1}$ on \m{\Dinf{\ell}{\ell}{G}}. Apply Proposition~\ref{p:dual}(d).
\end{proof}
\leanline{\lean{xUnip}, there for every \m{\Xinf{\ell}{N}{U}} with $U$ closed and normal and \m{\sigma\in \levU{\ell}{U}{1}}\sep
the reformulation of the congruence \lean{memLev}}

\begin{proposition}[Theorem D from one descending element]\label{p:adesc}
In Setting~\ref{set:tower}, assume the three hypotheses of Setting~\ref{set:coh}, write
\m{G=\GalG{L}{\overline L}}, and suppose that the module \m{\Xinf{\ell}{\ell}{G}} is generated by \m{R\ge1}
elements. Let \m{\sigma\in G} and an integer \m{a'\ge0} be such that
\begin{enumerate}[label=\textup{(\arabic*)},topsep=2pt,itemsep=1pt]
\item the element \m{u=\cyc{L}{\overline L}{\ell}{\sigma}} of~$\Z_\ell$ has level~$a'$;
\item the field $L$ contains a primitive root of unity of order~$\ell^{a'}$;
\item for all integers \m{h\ge1} and \m{m\ge0} and every number field~$L'$ that is a layer of index~$m$ over~$L$
for~$\ell$, \m{\shaInv{L'}{\ell}{h}=\len\cok\bigl((u^h\sigma)^{\ell^m}-1\bigr)}, where \m{u^h\sigma} is the operator
on \m{\Xinf{\ell}{\ell}{G}} of Input~2.
\end{enumerate}
Then the tower bound of Definition~\ref{d:abound} holds for $(\ell,R,L)$, with the field \m{K=\Q_\ell}, the
exponent~$a'$ and $R$ centres in~$\Z_\ell$.
\end{proposition}

\begin{proof}[Proof of Proposition~\ref{p:adesc} and of Theorem D]
For the proposition, the level~$a'$ is at least~$1$, so \m{u\equiv1\pmod\ell}. The module
\m{\Xinf{\ell}{\ell}{G}} is generated by $R$ elements, so it is finitely generated, and by
Corollary~\ref{c:unip} the action $\gamma$ of~$\sigma$ on it is unipotent modulo~$\ell$. So
Setting~\ref{set:alg} holds for \m{X=\Xinf{\ell}{\ell}{G}}, this pair $(\gamma,u)$ and the level~$a'$. The operator
\m{u^h\sigma} is \m{u^h\gamma}, so by (3) and Definition~\ref{d:tau},
\m{\shaInv{L'}{\ell}{h}=\tw{\gamma}{u}{h}{m}} for \m{h\ge1} and every such layer~$L'$. Theorem~A gives centres
\m{c_1,\dots,c_R\in\Z_\ell} with
\[
\shaInv{L'}{\ell}{h}\ \le\ R\,\bigl(\max\bigl(m,\mz{\ell}{R}\bigr)+a'+B\bigr)
\]
whenever \m{\vl{\ell}{c_i-h}\le B} for all~$i$, that is, \m{|c_i-h|\ge\ell^{-B}}. This is the first part of the
tower bound, with \m{K=\Q_\ell} and \m{R'=R}. For the second bound write \m{n=\mz{\ell}{R}}, let $L''$ be a number
field that is a layer of index~$n$ over~$L$ for~$\ell$, and let $\mathcal D$ be a local datum of~$L''$ at~$\ell$. By
Theorem~B, \m{\shaInv{L'}{\ell}{h}\le\tw{\gamma}{u}{h}{n}+R\cdot\max(m-n,0)}. By (3) the term
\m{\tw{\gamma}{u}{h}{n}} equals \m{\shaInv{L''}{\ell}{h}}, which is \m{\len\ShaOf{\mathcal D}{h}} by
Definition~\ref{d:ldata}, and this is at most \m{\VLen{\mathcal D}{h}} by Input~3 for the field \m{F=L''}, which
is totally complex because it contains~$L$.

For Theorem~D write \m{a=\cZero{\ell}{L}}. By Lemma~\ref{l:layers}(c) there is \m{\sigma\in G} that is a topological
generator of~$G$ for the cyclotomic character of~$\ell$; then \m{u=\cyc{L}{\overline L}{\ell}{\sigma}} has
level~$a$, and $L$ contains a primitive root of unity of order~$\ell^a$. By Input~2, condition~(3) holds for the
fields \m{L'=L_m}. A number field that is a layer of index~$m$ over~$L$ for~$\ell$ is isomorphic to~$L_m$ over~$L$
(Lemma~\ref{l:layers}(b)), and the length depends on the field only up to isomorphism; so (3) holds for every
layer. The proposition applies with \m{a'=a}: the tower bound holds for $(\ell,R,L)$, and for \m{L'=L_m} and
\m{L''=L_n} its two inequalities are (i) and~(ii).
\end{proof}
\leanline{\lean{aDesc}\sep\lean{aTower}}

\subhead{How the bound is used.}\phantomsection\label{s:use}
Three steps lead from Theorem~D to the window of the guide~\cite[§6]{Gu}; they complete the cluster. The statements
of this part are those of the archive, with their hypotheses. Only Corollary~\ref{c:exc} is proved here; the others
are not deformalized in this note, and the last one is stated without its objects, which belong to the guide.

\begin{definition}[the tower bound]\label{d:abound}
Let \m{R\ge0} be an integer and $L$ a field. \kw{The second bound holds for $(\ell,R,L)$} if for all integers
\m{h\ge1} and \m{m\ge0}, every number field~$L'$ that is a layer of index~$m$ over~$L$ for~$\ell$, every number
field~$L''$ that is a layer of index~$\mz{\ell}{R}$ over~$L$ for~$\ell$, and every local datum $\mathcal D$
of~$L''$ at~$\ell$,
\[
\shaInv{L'}{\ell}{h}\ \le\ \VLen{\mathcal D}{h}+R\cdot\max\bigl(m-\mz{\ell}{R},0\bigr).
\]
\kw{The tower bound holds for $(\ell,R,L)$} if the second bound holds for $(\ell,R,L)$ and there are a field~$K$
with an ultrametric absolute value that extends the $\ell$-adic absolute value of~$\Q_\ell$, integers \m{a'\ge0}
and \m{0\le R'\le R}, and elements \m{c_1,\dots,c_{R'}\in K}, such that $L$ contains a primitive root of unity of
order~$\ell^{a'}$ and such that, for all integers \m{h\ge1} and \m{m,B\ge0} with \m{|c_i-h|\ge\ell^{-B}} for
all~$i$, and every number field~$L'$ that is a layer of index~$m$ over~$L$ for~$\ell$,
\[
\shaInv{L'}{\ell}{h}\ \le\ R'\,\bigl(\max\bigl(m,\mz{\ell}{R}\bigr)+a'+B\bigr).
\]
\end{definition}
\leanline{\lean{SecondBound}\sep\lean{ABound}}

\begin{definition}[the first bound]\label{d:first}
Let \m{R,a',Z\ge0} be integers, $L$ a field and $\mathcal E$ a finite set of non-negative integers. \kw{The first
bound holds for $(\ell,R,a',L)$ up to~$Z$ outside~$\mathcal E$} if $\mathcal E$ has at most $R$ elements and, for
all integers~$h$ with \m{1\le h\le Z} that are not in~$\mathcal E$, all integers \m{m\ge0} and every number
field~$L'$ that is a layer of index~$m$ over~$L$ for~$\ell$,
\[
\shaInv{L'}{\ell}{h}\ \le\ R\,\bigl(\max\bigl(m,\mz{\ell}{R}\bigr)+a'+\lfloor\log_\ell Z\rfloor\bigr),
\]
where \m{\lfloor\log_\ell Z\rfloor} is read as~$0$ for \m{Z=0}.
\end{definition}
\leanline{\lean{FirstBound}}

\begin{corollary}[exceptional twists in the tower]\label{c:exc}
Let \m{R\ge0} be an integer and $L$ a field such that the tower bound holds for $(\ell,R,L)$. Then there is an
integer \m{a'\ge0} such that $L$ contains a primitive root of unity of order~$\ell^{a'}$ and such that for every
integer \m{Z\ge0} there is a finite set $\mathcal E$ of non-negative integers for which the first bound holds for
$(\ell,R,a',L)$ up to~$Z$ outside~$\mathcal E$.
\end{corollary}
\begin{proof}
Take $K$, $a'$, $R'$ and the centres $c_i$ of Definition~\ref{d:abound}, and let $\mathcal E$ be the set of the
integers \m{1\le h\le Z} with \m{|c_i-h|<\ell^{-\lfloor\log_\ell Z\rfloor}} for some~$i$. By Lemma~\ref{l:sep} each
centre contributes at most one element, so $\mathcal E$ has at most \m{R'\le R} elements. For \m{1\le h\le Z}
outside~$\mathcal E$ the tower bound applies with \m{B=\lfloor\log_\ell Z\rfloor}, and \m{R'\le R}.
\end{proof}
\leanline{\lean{exc}}

By Theorem~D the hypothesis of Corollary~\ref{c:exc} holds for the field~$L$ of Setting~\ref{set:tower}, and the
corollary is then Corollary~C for the lengths \m{\shaInv{L_m}{\ell}{h}}.

\emph{Finite generation.} The hypothesis of Theorem~D on the Iwasawa module is a condition on a finite group.

\begin{lemma}[finite generation]\label{l:rank}
Let $L$ be a number field with an algebraic closure~$\overline L$, write \m{G=\GalG{L}{\overline L}}, and let
\m{R\ge0} be an integer. The following are equivalent.
\begin{enumerate}[label=\textup{(\alph*)},topsep=2pt,itemsep=1pt]
\item The field $L$ is totally complex and contains the $\ell$-th roots of unity, and the fourth roots of unity
if \m{\ell=2}; \m{R\ge1}; and the module \m{\Xinf{\ell}{\ell}{G}} is generated by $R$ elements.
\item The field $L$ contains the $\ell$-th roots of unity, and the fourth roots of unity if \m{\ell=2};
\m{R\ge1}; and \m{\xRank{\ell}{\ell}{G}\le R}, where the left-hand side is the length of the submodule of the
elements of \m{\Dinf{\ell}{\ell}{G}} killed by~$\ell$.
\end{enumerate}
\end{lemma}
\leanline{\lean{isTowerIff}; the conditions in (a) are the hypotheses of Theorem~D on $(\ell,L,R)$}

\begin{theorem}[Ferrero and Washington, in the form of the archive]\label{t:fw}
Let $F$ be a number field that is Galois over~$\Q$ with an abelian Galois group and that contains the $\ell$-th
roots of unity, and the fourth roots of unity if \m{\ell=2}. There is an integer \m{r\ge0} such that for every
integer \m{n\ge0} and every number field~$F'$ over~$F$ that contains a primitive root of unity of order~$\ell^n$
and is generated over~$F$ by the roots of unity of order dividing~$\ell^n$, the quotient of the ideal class group
of~$F'$ by the subgroup of the $\ell$-th powers has at most $\ell^r$ elements.
\end{theorem}
\leanline{the statement \lean{FW}\sep its proof in the archive \lean{FWthm}}

This is the theorem of Ferrero and Washington~\cite{FW} that the Iwasawa invariant~$\mu$ of the cyclotomic
$\Z_\ell$-extension of an abelian number field vanishes, in the equivalent form that the $\ell$-rank of the class
groups is bounded along the cyclotomic tower \cite[Ch.\,13]{Wa}. The archive states it as a proposition and proves
it; the statements that use it take it as a hypothesis. From it the archive deduces a bound for the length in
Lemma~\ref{l:rank}(b), with one $R$ for all the fields isomorphic to a given abelian field:

\begin{proposition}[one $R$ for an abelian field]\label{p:family}
Assume the statement of Theorem~\ref{t:fw}. Let \m{\ell\ne2}, let \m{q\ge1} be an integer divisible by~$\ell$, and
let $\Delta$ be a subgroup of \m{(\Z/q\Z)^\times}, the Galois group of the field of the $q$-th roots of unity
over~$\Q$, all of whose elements are congruent to~$1$ modulo~$\ell$. There is an integer \m{R\ge0} such that for
every number field~$L$ isomorphic to the subfield fixed by~$\Delta$, with an algebraic closure~$\overline L$ and
\m{G=\GalG{L}{\overline L}}, the module \m{\Xinf{\ell}{1}{G}} is generated by $R$ elements and the conditions of
Lemma~\ref{l:rank}(a) hold for $(\ell,L,R)$.
\end{proposition}
\leanline{\lean{towerFixed}}

\begin{proposition}[the tower bound for the fields of the guide]\label{p:fixed}
Assume the statement of Theorem~\ref{t:fw}. Assume also that the conclusion of Theorem~D holds without further
hypotheses: for every prime number in the place of~$\ell$, every number field~$L$ and every integer~$R$ such that
the conditions of Lemma~\ref{l:rank}(a) hold, the tower bound holds for the prime, $R$ and~$L$.
\begin{enumerate}[label=\textup{(\alph*)},topsep=2pt,itemsep=1pt]
\item Let \m{\ell\ne2}, and let $q$ and~$\Delta$ be as in Proposition~\ref{p:family}. For every number field~$L$
isomorphic to the subfield fixed by~$\Delta$ there is an integer~$R$ such that the tower bound holds for
$(\ell,R,L)$.
\item For every integer \m{q\ge1} and every subgroup $\Delta$ of \m{(\Z/q\Z)^\times} there is an integer~$R$ such
that, for all prime numbers \m{p\equiv3\pmod 8} and \m{p'\equiv7\pmod{16}} that do not divide~$q$, there is a number
field~$L$ for which the tower bound holds for $(2,R,L)$ and which is a base field of the tower for the prime~$2$ in
the sense of the guide, for $q$, $\Delta$, $p$ and~$p'$. This last condition ties $L$ to the fields of §6 of the
guide; it is not deformalized here.
\end{enumerate}
\end{proposition}
\leanline{(a) \lean{fixedIn}\sep (b) \lean{twoIn}}

\emph{The reference layer.} For the second cohomology group in Theorem~D\,(ii) the cluster takes as an input that
it is finite, with a bound for its order in terms of the twist, the degree and the discriminant of the field:

\begin{definition}[the cohomology bound]\label{d:K}
Let $C$ be a real number. \kw{The cohomology bound holds with the constant~$C$} if for every prime number, written
$\ell$ again, every totally complex number field~$L$ that contains the $\ell$-th roots of unity, and the fourth
roots of unity if \m{\ell=2}, every local datum $\mathcal D$ of~$L$ at~$\ell$ and every integer \m{h\ge1}, the
group \m{H^2(\mathcal D,\Z_\ell(h+1))} is finite and
\[
\log\,\#H^2\bigl(\mathcal D,\Z_\ell(h+1)\bigr)\ \le\ \exp\Bigl(\bigl(2h\,[L:\Q]\bigr)^{C}\,\log\bigl(2\,|d_L|\bigr)\Bigr),
\]
where $\#$ is the number of elements, \m{[L:\Q]} the degree and $d_L$ the discriminant of~$L$.
\end{definition}
\leanline{\lean{LemmaK}}

So all the lengths in Theorem~D are finite, also at a centre, and the exceptional twists cost a bounded amount.

\emph{The window.} By Corollary~\ref{c:exc}, for each of finitely many primes~$\ell$ there are at most $R_\ell$
exceptional twists up to~$Z$, where $R_\ell$ is the number of generators for~$\ell$. Their union has at most
\m{\sum_\ell R_\ell} elements, whatever~$Z$. An exceptional twist spoils the estimate only on one interval of
scales, and an interval lemma finds a range of scales that avoids all of them:

\begin{lemma}[a window]\label{l:window}
Let \m{R\ge0} be an integer, let $\theta$ and~$\Theta$ be real numbers with \m{0<\theta\le\Theta}, let \m{\Xi>1} be
a real number, and let $\mathcal E$ be a finite set of non-negative integers with at most $R$ elements. Put
\m{J=2^{R+2}(\Theta/\theta)^R}. There is a real number~$Y$ with \m{\Xi^{1/J}\le Y} and \m{Y^2\le\Xi} such that for
every real number~$\omega$ with \m{Y\le\omega\le Y^2}, every \m{h\in\mathcal E} with \m{h\le\omega^\Theta}
satisfies \m{h\le\omega^\theta}.
\end{lemma}
\leanline{\lean{window}}

\begin{proposition}[the window of the guide]\label{p:window}
\begin{enumerate}[label=\textup{(\alph*)},topsep=2pt,itemsep=1pt]
\item Suppose that the cohomology bound of Definition~\ref{d:K} holds with some constant~$C$, and that the
conclusions (a) and~(b) of Proposition~\ref{p:fixed} hold, the first for every odd prime in the place of~$\ell$.
Then the statement of the window of the guide \cite[§6]{Gu} holds: there are constants, namely an integer for
every pair $(q,\Delta)$, an integer for every prime and every pair $(q,\Delta)$, and a positive real number, such
that for every real number \m{\Theta\ge1} the two families of bounds of the guide hold, with these constants, on a
window of scales.
\item Assume the three hypotheses of Setting~\ref{set:coh} and the statement of Theorem~\ref{t:fw}, and suppose that
the cohomology bound of Definition~\ref{d:K} holds with some constant~$C$. Then the conclusion of~(a) holds.
\end{enumerate}
\end{proposition}
\leanline{(a) \lean{thmBab}\sep (b) \lean{thmBabK}}

The statement of the window is not deformalized in this note: its fields and its two families of bounds are the
subject of §6 of the guide. Part~(b) follows from part~(a), Theorem~D and Proposition~\ref{p:fixed}.
\begin{remark}[the classical picture]\label{r:classical}
Let $X$ be a module that is free of rank~$R$, let \m{\gamma\in\End(X)} and let $u$ be a unit of~$\Z_\ell$, and let
\m{\Phi(T)=\det(T-(\gamma-1))} be the characteristic polynomial of \m{\gamma-1}. Factoring \m{(u^h\gamma)^{\ell^m}-1} over the roots of unity gives
\[
\tw{\gamma}{u}{h}{m}=\sum_{\zeta^{\ell^m}=1}\vlb{\ell}{\Phi(\zeta u^{-h}-1)}.
\]
At the layer~$0$ this reads \m{\tw{\gamma}{u}{h}{0}=\vl{\ell}{\Phi(u^{-h}-1)}}: the twists with a large length are the integers
close to a zero in~$\Z_\ell$ of the function \m{s\mapsto\Phi(u^{-s}-1)}. For the modules of this section, $\Phi$ is
the characteristic polynomial of Iwasawa's theory \cite{Iw}, \cite[Ch.\,13]{Wa}, and for abelian fields and odd~$\ell$
the main conjecture, proved by Mazur and Wiles~\cite{MW}, expresses the characteristic polynomials of closely related
modules through $\ell$-adic $L$-functions. In this language an exceptional twist is, roughly, an integer that is
$\ell$-adically very close to a zero of an $\ell$-adic $L$-function. None of this is used in the archive: the
centres there come from a compactness argument, and only their number matters.
\end{remark}

% ------------------------------------------------------------------------------------- 8
\sect{Descent to a layer}{s:descent}
This section proves Input~2. The cohomological part of the proof is one isomorphism, which the archive proves in
the neighbouring directory on towers of number fields. It is stated below as Input~2a, with a companion, Input~2b.
Everything else is proved here: duality for modules over~$\Z_\ell$, the group of principal units, and the Galois
theory of the tower.

\subhead{Characters.} For a module $M$ put \m{M^\vee=\Hom_{\Z_\ell}(M,\QZl{\ell})}, the module of \kw{characters}
of~$M$. Thus \m{X=D^\vee} in Proposition~\ref{p:dual}, and \m{\Xinf{\ell}{N}{U}=\Dinf{\ell}{N}{U}^\vee} in
Definition~\ref{d:dinf}. By Proposition~\ref{p:dual}(a) the module $\QZl{\ell}$ is injective: a homomorphism
from a submodule of~$M$ to~$\QZl{\ell}$ extends to a character of~$M$.

\begin{lemma}\label{l:char}
Let $M$ be a module over~$\Z_\ell$ and $W$ a submodule of~$M$, and form the modules of characters with values
in~$\QZl{\ell}$.
\begin{enumerate}[label=\textup{(\alph*)},topsep=2pt,itemsep=2pt]
\item For every non-zero \m{x\in M} there is a character $\chi$ of~$M$ with \m{\chi(x)\ne0}.
\item The restriction map \m{M^\vee\to W^\vee} is surjective. Its kernel is the image of the injective map
\m{(M/W)^\vee\to M^\vee} given by composition with the quotient map \m{M\to M/W}.
\item \m{\len M^\vee=\len M}.
\end{enumerate}
\end{lemma}
\begin{proof}
(a) The elements \m{t\in\Z_\ell} with \m{tx=0} form an ideal other than~$\Z_\ell$, so this ideal is $0$ or
$\ell^k\Z_\ell$ with \m{k\ge1}. Let $y$ be the class of $\ell^{-k}$ in~$\QZl{\ell}$, and the class of $\ell^{-1}$ if the
ideal is~$0$. Then \m{y\ne0}, and \m{ty=0} whenever \m{tx=0}. So \m{tx\mapsto ty} is a well-defined homomorphism
from the submodule $\Z_\ell x$ to~$\QZl{\ell}$, and a character $\chi$ that extends it has \m{\chi(x)=y}.

(b) Restriction is surjective because $\QZl{\ell}$ is injective. A character of~$M$ vanishes on~$W$ if and only if it
factors through \m{M/W}, and then it does so in exactly one way.

(c) Let $M$ have finite length~$n$. We use induction on~$n$; for \m{n=0} both modules are zero. Let \m{n\ge1}. A
composition series of~$M$ gives a submodule $W$ of length \m{n-1} such that \m{M/W} is simple. A simple module
over~$\Z_\ell$ is isomorphic to \m{\Z_\ell/\ell\Z_\ell}, so its characters correspond, by their value at the class
of~$1$, to the elements of~$\QZl{\ell}$ killed by~$\ell$. These form the group \m{\ell^{-1}\Z_\ell/\Z_\ell}, which has
length~$1$. By (b) and the induction hypothesis,
\m{\len M^\vee=\len(M/W)^\vee+\len W^\vee=1+(n-1)}.

Now let $M$ be any module. For a submodule $W$ let \m{W^\perp\subseteq M^\vee} be the submodule of the characters
that vanish on~$W$. If \m{W\subsetneq W'} are submodules of~$M$, then \m{{W'}^\perp\subsetneq W^\perp}: by (a) the
module \m{M/W} has a character that does not vanish on the image of~$W'$, and its composition with \m{M\to M/W}
lies in~$W^\perp$ and not in~${W'}^\perp$. So a chain of submodules of~$M$ gives a chain of the same length
in~$M^\vee$, and \m{\len M\le\len M^\vee}. If the length of~$M$ is infinite, so is the length of~$M^\vee$.
\end{proof}
\leanline{(a) \lean{qzSep}, with \lean{qzAnn}\sep (b) \lean{exactDual}\sep (c) \lean{lenDual},
with \lean{simpleDual} and \lean{leDual}}

\begin{lemma}[kernel and cokernel]\label{l:cokdual}
Let $D$ be a module over~$\Z_\ell$, let \m{g\in\End(D)}, and let $g^\vee$ be the endomorphism of \m{X=D^\vee} given by
\m{(g^\vee x)(\varphi)=x(g\varphi)}, as in Proposition~\ref{p:dual}. The image of~$g^\vee$ consists of the characters
of~$D$ that vanish on the kernel $\ker g$ of~$g$, the restriction of characters to any submodule of~$D$ is
surjective, and
\[
\len\cok g^\vee=\len\ker g .
\]
\end{lemma}
\begin{proof}
A character of the form \m{g^\vee x=x\circ g} vanishes on~$\ker g$. Conversely, let \m{y\in X} vanish on~$\ker g$.
Then \m{g\varphi\mapsto y(\varphi)} is a well-defined homomorphism \m{g(D)\to\QZl{\ell}}. It extends to some \m{x\in X},
and \m{y=g^\vee x}. So the image of~$g^\vee$ is the kernel of the restriction map \m{X\to(\ker g)^\vee}, which is
surjective by Lemma~\ref{l:char}(b). Hence $\cok g^\vee$ is isomorphic to $(\ker g)^\vee$, and
Lemma~\ref{l:char}(c) gives the lengths.
\end{proof}
\leanline{the image and the surjectivity \lean{dualKer}, there for a homomorphism between two modules\sep the
lengths \lean{cokerDual}}

\subhead{Principal units.} Recall from §\ref{s:arith} the groups \m{\princ{\ell}{0}=\Z_\ell^\times} and
\m{\princ{\ell}{n}=1+\ell^n\Z_\ell} for \m{n\ge1}. An element of~$\Z_\ell$ has level~$n$
(Definition~\ref{d:level}) if and only if it lies in $\princ{\ell}{n}$ and not in~$\princ{\ell}{n+1}$, where \m{n\ge1},
and \m{n\ge2} if \m{\ell=2}.

\begin{lemma}\label{l:princ}
Let \m{a\ge1} be an integer, and \m{a\ge2} if \m{\ell=2}.
\begin{enumerate}[label=\textup{(\alph*)},topsep=2pt,itemsep=2pt]
\item For \m{n\ge1} the index of $\princ{\ell}{n}$ in~$\Z_\ell^\times$ is \m{\ell^{n-1}(\ell-1)}. For \m{m\ge0} the
index of $\princ{\ell}{a+m}$ in~$\princ{\ell}{a}$ is~$\ell^m$.
\item If \m{u\in \princ{\ell}{a}}, then \m{\vl{\ell}{u^{\ell^m}-1}=\vl{\ell}{u-1}+m} for every integer \m{m\ge0}. In
particular, if $u$ has level~$a$, then $u^{\ell^m}$ has level \m{a+m}.
\item An element $u$ of~$\princ{\ell}{a}$ generates $\princ{\ell}{a}$ topologically, that is, $\princ{\ell}{a}$ lies in the
closure of the set of the powers of~$u$, if and only if $u$ does not lie in~$\princ{\ell}{a+1}$; and then $u$ has
level~$a$.
\end{enumerate}
\end{lemma}
\begin{proof}
(a) Reduction modulo~$\ell^n$ is a surjective homomorphism from $\Z_\ell^\times$ onto the group of units of the ring
\m{\Z_\ell/\ell^n\Z_\ell}, which has \m{\ell^{n-1}(\ell-1)} elements, and its kernel is~$\princ{\ell}{n}$. The second
index is the quotient of two of these numbers.

(b) If \m{u=1}, both sides are infinite. Otherwise $u$ has the level \m{\vl{\ell}{u-1}\ge a}, and Lemma~\ref{l:lte}
gives \m{\vl{\ell}{u^{\ell^m}-1}=\vl{\ell}{u-1}+m}.

(c) An element of~$\princ{\ell}{a}$ that is not in~$\princ{\ell}{a+1}$ has level~$a$. If \m{u\in \princ{\ell}{a+1}}, all
powers of~$u$ lie in the closed subgroup~$\princ{\ell}{a+1}$, which does not contain the element
\m{1+\ell^a} of~$\princ{\ell}{a}$. Now let $u$ have level~$a$, and let \m{m\ge0}. By Lemma~\ref{l:lte} an integer~$n$
satisfies \m{u^n\in \princ{\ell}{a+m}} if and only if $\ell^m$ divides~$n$. So the powers
\m{u^0,u^1,\dots,u^{\ell^m-1}} lie in different cosets of $\princ{\ell}{a+m}$ in~$\princ{\ell}{a}$, and by (a) there are no
other cosets. Thus every element of~$\princ{\ell}{a}$ is congruent modulo~$\ell^{a+m}$ to a power of~$u$, for
every~$m$: it lies in the closure of the set of powers of~$u$.
\end{proof}
\leanline{\lean{princGen}\sep (a) \lean{pgIndex}\sep (b) \lean{pgEmult}, from \lean{lteOdd} and \lean{lteSq}\sep (c)
\lean{pgExists}, \lean{princClosed}}

\subhead{Layers.} The groups of the layers are subgroups of the Galois group that the archive defines through the
exponent of Definition~\ref{d:czero}.

\begin{definition}[the group of a layer]\label{d:lay}
Let $L$ be a number field with an algebraic closure~$\overline L$, and let \m{m\ge0} be an integer. In the notation
of §\ref{s:arith}, and with \m{a=\cZero{\ell}{L}},
\[
\layU{\ell}{L}{\overline L}{m}=\Un{L}{\overline L}{\ell}{a+m}
=\bigl\{\sigma\in\GalG{L}{\overline L}:\ \cyc{L}{\overline L}{\ell}{\sigma}\in \princ{\ell}{a+m}\bigr\}.
\]
It is a function of the prime, the field, its closure and~$m$.
\end{definition}
\leanline{\lean{layU}}

\begin{lemma}[roots of unity]\label{l:roots}
Let $L$ be a number field with an algebraic closure~$\overline L$; the hypothesis of Setting~\ref{set:tower} on the
roots of unity is not assumed. Write \m{G=\GalG{L}{\overline L}}, and for \m{n\ge0} let \m{\zeta_n\in\overline L} be a
root of unity of order~$\ell^n$.
\begin{enumerate}[label=\textup{(\alph*)},topsep=2pt,itemsep=2pt]
\item For an integer \m{n\ge0}, the field $L$ contains a primitive root of unity of order~$\ell^n$ if and only if
\m{n\le\cZero{\ell}{L}}.
\item Let \m{F\subseteq F'} be subfields of~$\overline L$ such that $F'$ contains~$\zeta_n$ and is generated
over~$F$ by the roots of unity of order dividing~$\ell^n$. Then the elements of~$G$ that fix every element of~$F'$
are the elements~$\sigma$ of~$G$ that fix every element of~$F$ and satisfy
\m{\cyc{L}{\overline L}{\ell}{\sigma}\in \princ{\ell}{n}}. In particular \m{\galOf{L}{L(\zeta_n)}{\overline L}=\levU{\ell}{G}{n}},
and \m{\levU{\ell}{G}{n}=G} if and only if \m{n\le\cZero{\ell}{L}}.
\item The field $L$ contains the $\ell$-th roots of unity, and the fourth roots of unity if \m{\ell=2}, if and only
if \m{\cZero{\ell}{L}\ge1}, and \m{\cZero{\ell}{L}\ge2} if \m{\ell=2}.
\item If $L$ contains the $\ell$-th roots of unity, and the fourth roots of unity if \m{\ell=2}, then $L$ is
totally complex.
\item Let \m{\ell=2}. Then \m{\cyc{L}{\overline L}{2}{\sigma}\equiv1\pmod 4} for all \m{\sigma\in G} if and only
if $L$ contains the fourth roots of unity.
\end{enumerate}
\end{lemma}
\begin{proof}
(a) A power of a primitive root of unity of order~$\ell^n$ is one of order~$\ell^{n'}$, for every \m{n'\le n}; so
the integers in question are those up to the largest one, which is \m{\cZero{\ell}{L}} by
Definition~\ref{d:czero}.

(b) An automorphism that fixes~$F$ fixes $F'$ if and only if it fixes the roots of unity of order
dividing~$\ell^n$. Since \m{\sigma(\zeta)=\zeta^u} with \m{u=\cyc{L}{\overline L}{\ell}{\sigma}} for these roots,
this happens if and only if \m{u\equiv1\pmod{\ell^n}}. The case \m{F=L}, \m{F'=L(\zeta_n)} is the first formula.
Finally \m{\levU{\ell}{G}{n}=G} means that $\zeta_n$ lies in~$L$, and (a) applies.

(c) This is (a) for \m{n=1}, and for \m{n=2} if \m{\ell=2}.

(d) A real embedding of~$L$ would map a root of unity of order $\ell$, or of order~$4$ if \m{\ell=2}, to a real
root of unity of the same order; but the real roots of unity are $1$ and~$-1$.

(e) By (b), the congruence for all~$\sigma$ says that \m{\levU{2}{G}{2}=G}, that is, \m{2\le\cZero{2}{L}}; by (a) this
means that $L$ contains a root of unity of order~$4$.
\end{proof}
\leanline{(a), the second part of (b), (c), (d), (e) \lean{muQ}, where the condition of (e), or \m{\ell\ne2}, is one
notion\sep (b) \lean{galOfCyc}}

\begin{lemma}\label{l:layers}
In Setting~\ref{set:tower}, write \m{G=\GalG{L}{\overline L}} and \m{a=\cZero{\ell}{L}}, and for \m{n\ge0} let
\m{\zeta_n\in\overline L} be a root of unity of order~$\ell^n$.
\begin{enumerate}[label=\textup{(\alph*)},topsep=2pt,itemsep=2pt]
\item \m{\cyc{L}{\overline L}{\ell}{G}=\princ{\ell}{a}}, where \m{a\ge1}, and \m{a\ge2} if \m{\ell=2}. For every
\m{m\ge0} the index of \m{\layU{\ell}{L}{\overline L}{m}} in~$G$ is~$\ell^m$, the field $L_m$ is
\m{L(\zeta_{a+m})}, it is a layer of index~$m$ over~$L$ for~$\ell$ (Definition~\ref{d:layer}), and
\m{\galOf{L}{L_m}{\overline L}=\layU{\ell}{L}{\overline L}{m}}.
\item Let \m{L'\subseteq\overline L} be an extension of~$L$ that contains~$\zeta_k$ and is generated over~$L$ by the
roots of unity of order dividing~$\ell^k$, for an integer \m{k\ge0}, and put \m{m=\max(k-a,0)}. Then
\m{L'=L_m}, the degree of $L'$ over~$L$ is~$\ell^m$, and
\m{\galOf{L}{L'}{\overline L}=\layU{\ell}{L}{\overline L}{m}}. Consequently a subfield of~$\overline L$ that is a layer
of index~$m$ over~$L$ for~$\ell$ is~$L_m$, and a number field that is a layer of index~$m$ over~$L$ for~$\ell$ is
isomorphic to~$L_m$ over~$L$.
\item There is \m{\sigma\in G} that is a topological generator of~$G$ for the cyclotomic character of~$\ell$
(Definition~\ref{d:isgen}), and $L$ contains a primitive root of unity of order~$\ell^a$. For every such~$\sigma$
the element \m{\cyc{L}{\overline L}{\ell}{\sigma}} of~$\Z_\ell$ has level~$a$, the element $\sigma$ lies
in~$\levU{\ell}{G}{1}$, and for every \m{m\ge0} the element $\sigma^{\ell^m}$ is a topological generator of
\m{\layU{\ell}{L}{\overline L}{m}} for the cyclotomic character of~$\ell$.
\end{enumerate}
\end{lemma}
\begin{proof}
(a) By Lemma~\ref{l:roots}(a) the field $L$ contains~$\zeta_a$ and not~$\zeta_{a+1}$, and by
Lemma~\ref{l:roots}(c) and Setting~\ref{set:tower} the inequalities for~$a$ hold. Every element of~$G$
fixes~$\zeta_a$, so \m{\cyc{L}{\overline L}{\ell}{G}\subseteq \princ{\ell}{a}}. Some \m{\sigma\in G}
moves~$\zeta_{a+1}$. Then \m{\cyc{L}{\overline L}{\ell}{\sigma}\notin \princ{\ell}{a+1}}, so this element has
level~$a$, and by Lemma~\ref{l:princ}(c) the closure of the set of its powers contains~$\princ{\ell}{a}$. The image
\m{\cyc{L}{\overline L}{\ell}{G}} is compact, hence closed, and it contains these powers. So
\m{\cyc{L}{\overline L}{\ell}{G}=\princ{\ell}{a}}.

The cyclotomic character maps the cosets of \m{\layU{\ell}{L}{\overline L}{m}} in~$G$ bijectively onto the cosets of
$\princ{\ell}{a+m}$ in~$\princ{\ell}{a}$, and there are $\ell^m$ of these by Lemma~\ref{l:princ}(a). By
Lemma~\ref{l:roots}(b) this group is \m{\galOf{L}{L(\zeta_{a+m})}{\overline L}}, and by Galois theory its index in~$G$ is
the degree of \m{L(\zeta_{a+m})} over~$L$. So \m{L(\zeta_{a+m})} is a subfield of~$L_\infty$ of degree~$\ell^m$
over~$L$, and it is a layer of index~$m$. It is the only subfield of this degree: such a subfield is the fixed
field of the preimage, under the cyclotomic character, of a closed subgroup of index~$\ell^m$
in~$\princ{\ell}{a}$. The quotient by this subgroup has $\ell^m$ elements, so the subgroup contains $u^{\ell^m}$ for a
topological generator $u$ of~$\princ{\ell}{a}$. Hence it contains $\princ{\ell}{a+m}$ by Lemma~\ref{l:princ}(b)
and~(c), and the two groups have the same index.

(b) We have \m{L'=L(\zeta_k)}. If \m{k\le a}, then \m{L'=L=L_0} and \m{m=0}. If \m{k\ge a}, then by (a)
the field \m{L'=L_{k-a}} has degree~$\ell^{k-a}$ over~$L$. The group is given by~(a). If moreover $L'$ has
degree~$\ell^{m'}$ over~$L$ for an integer~$m'$, then \m{m'=m}; and a number field that is a layer over~$L$ can be
embedded into~$\overline L$ over~$L$.

(c) The element~$\sigma$ of the proof of~(a) has \m{\cyc{L}{\overline L}{\ell}{\sigma}} of level~$a$, so by
Lemma~\ref{l:princ}(c) the closure of the set of the powers of this element contains
\m{\princ{\ell}{a}=\cyc{L}{\overline L}{\ell}{G}}: it is a topological generator. The root of unity~$\zeta_a$ lies
in~$L$. Now let $\sigma$ be any topological generator of~$G$. The element
\m{\cyc{L}{\overline L}{\ell}{\sigma}} of~$\princ{\ell}{a}$ generates $\princ{\ell}{a}$ topologically, so it has
level~$a$ by Lemma~\ref{l:princ}(c); as \m{a\ge1}, the element $\sigma$ lies in~$\levU{\ell}{G}{1}$. By
Lemma~\ref{l:princ}(b), \m{\cyc{L}{\overline L}{\ell}{\sigma^{\ell^m}}=\cyc{L}{\overline L}{\ell}{\sigma}^{\ell^m}}
has level \m{a+m}. Hence \m{\sigma^{\ell^m}\in\layU{\ell}{L}{\overline L}{m}}, and by Lemma~\ref{l:princ}(c), with
\m{a+m} in place of~$a$, the closure of the set of the powers of its image contains~$\princ{\ell}{a+m}$, which is the
image of \m{\layU{\ell}{L}{\overline L}{m}} by~(a).
\end{proof}
\leanline{(a) \lean{cycRange}, \lean{layerGen}, \lean{layerUniq}, with \lean{galEmbTop}\sep (b) \lean{luKey},
\lean{layerUniq}\sep (c) \lean{layerGen}, \lean{genLevel}}

\subhead{Twists.} The actions of the Galois group on the two modules of Definition~\ref{d:dinf} are twisted by powers
of the cyclotomic character.

\begin{definition}[twisted modules]\label{d:tw}
Let $F$ be a number field with an algebraic closure~$\overline F$, let \m{N\ge0} be an integer, $U$ a normal
subgroup of \m{\GalG{F}{\overline F}}, $V$ any subgroup of \m{\GalG{F}{\overline F}}, and $h$ an integer. Then
\[
\dTw{\ell}{N}{U}{h}\qquad\text{and}\qquad\xTw{\ell}{N}{U}{h}
\]
are the modules \m{\Dinf{\ell}{N}{U}} and \m{\Xinf{\ell}{N}{U}} on which \m{\sigma\in\GalG{F}{\overline F}} acts
as \m{\cyc{F}{\overline F}{\ell}{\sigma}^{-h}\sigma} and as \m{\cyc{F}{\overline F}{\ell}{\sigma}^{h}\sigma},
respectively, and
\[
\ShaD{\ell}{N}{U}{V}{h}
\]
is the submodule of the elements of \m{\dTw{\ell}{N}{U}{h}} that are fixed by this action of every element of~$V$.
\end{definition}
\leanline{\lean{dTw}\sep\lean{xTw}\sep\lean{ShaD}}

\begin{lemma}[the two twists are dual]\label{l:twdual}
In the situation of Definition~\ref{d:tw}, with the normal subgroup~$U$ and the integers $N$ and~$h$, let
\m{\sigma\in\GalG{F}{\overline F}}, \m{x\in \Xinf{\ell}{N}{U}} and \m{\varphi\in \Dinf{\ell}{N}{U}}, and put
\m{u=\cyc{F}{\overline F}{\ell}{\sigma}}. Then
\[
\bigl(u^{h}\sigma x\bigr)(\varphi)=x\bigl(u^{h}\sigma^{-1}\varphi\bigr),
\]
and \m{u^h=\cyc{F}{\overline F}{\ell}{\sigma^{-1}}^{-h}}. In the notation of Proposition~\ref{p:dual}: the operator
by which $\sigma$ acts on \m{\xTw{\ell}{N}{U}{h}} is~$g^\vee$, where $g$ is the operator by which $\sigma^{-1}$ acts
on \m{\dTw{\ell}{N}{U}{h}}.
\end{lemma}
\begin{proof}
By Definition~\ref{d:dinf}, \m{(\sigma x)(\varphi)=x(\sigma^{-1}\varphi)}, and $x$ is $\Z_\ell$-linear.
\end{proof}
\leanline{\lean{xTwApply}}

\statement{inksoft}{\phantomsection\label{i:2a}\textbf{Input 2a} (the classes over a layer).\enspace\emph{In
Setting~\ref{set:tower}, write \m{G=\GalG{L}{\overline L}}. Let \m{h\ne0} be an integer, let $F$ be a number field
that is a finite extension of~$L$ contained in~$L_\infty$, and let $\mathcal D$ be a local datum of~$F$ at~$\ell$
whose algebraic closure is~$\overline L$. Then there is an isomorphism of $\Z_\ell$-modules from
$\ShaOf{\mathcal D}{h}$ onto \m{\ShaD{\ell}{\ell}{G}{\galOf{L}{F}{\overline L}}{h}}, given by restricting classes to
\m{\galOf{L}{L_\infty}{\overline L}}.}

\vspace{3pt}\phantomsection\label{i:2b}\textbf{Input 2b} (invariants).\enspace\emph{In Setting~\ref{set:tower}, write
\m{G=\GalG{L}{\overline L}}. Let $V$ be a subgroup of~$G$, and let \m{\eta\in G} be a topological generator of~$V$
for the cyclotomic character of~$\ell$ (Definition~\ref{d:isgen}). Then for every integer~$h$ the module
\m{\ShaD{\ell}{\ell}{G}{V}{h}} is the kernel of \m{g-1}, where $g$ is the operator by which $\eta^{-1}$ acts on
\m{\dTw{\ell}{\ell}{G}{h}}; and it is also the kernel of the same expression formed with~$\eta$.}}
\leanline{2a \lean{desc}\sep 2b \lean{shaDGen}}

Both are proved in the neighbouring directory of the archive, there for any integer~$N$ in the second slot of
Input~2b and any closed normal subgroup in the third. We indicate the standard arguments; they are not
deformalized here.

\emph{On Input 2a.} Write \m{G=\GalG{L}{\overline L}}. The group \m{\GalG{F}{L_\infty}} is isomorphic to~$\Z_\ell$, and a topological generator of it
acts on \m{\qzG{F}{\overline L}{\ell}{h}} by a scalar different from~$1$ when \m{h\ne0}. So the cohomology of
this group with these coefficients vanishes in degrees $1$ and~$2$, and restriction identifies the classes
over~$F$ with the homomorphisms
\m{\galOf{L}{L_\infty}{\overline L}\to\QZl{\ell}} that are invariant under the twisted action of \m{\galOf{L}{F}{\overline L}}. The
same holds place by place, so the local conditions correspond, and $\ShaOf{\mathcal D}{h}$ becomes the module of
invariants in~\m{\dTw{\ell}{\ell}{G}{h}}.

\emph{On Input 2b.} With $G$, $V$ and~$\eta$ as in the statement, an element of~\m{\Dinf{\ell}{\ell}{G}} has finite order and is fixed by an open subgroup
of~$G$, so it is fixed by the twisted action of~$\levU{\ell}{G}{n}$ when $n$ is large. For such~$n$ every \m{\rho\in V} is
the product of a power of~$\eta$ and an element of~$\levU{\ell}{G}{n}$, because \m{\cyc{L}{\overline L}{\ell}{\rho}} is
congruent modulo~$\ell^n$ to a power of~\m{\cyc{L}{\overline L}{\ell}{\eta}}. Hence an element fixed by~$\eta$ is
fixed by~$V$; and it is fixed by~$\eta$ exactly when it is fixed by~$\eta^{-1}$.

\begin{proof}[Proof of Input 2]
Write \m{G=\GalG{L}{\overline L}}. Let \m{h\ne0} and \m{m\ge0} be integers, put \m{\eta=\sigma^{\ell^m}} and
\m{V=\layU{\ell}{L}{\overline L}{m}}. By Lemma~\ref{l:layers}(a) the field $L_m$ lies in~$L_\infty$ and
\m{\galOf{L}{L_m}{\overline L}=V}. Choose a local datum $\mathcal D$ of~$L_m$ at~$\ell$ whose algebraic closure
is~$\overline L$. By Definition~\ref{d:ldata} and Input~2a for \m{F=L_m},
\[
\shaInv{L_m}{\ell}{h}=\len\ShaOf{\mathcal D}{h}=\len\ShaD{\ell}{\ell}{G}{V}{h} .
\]
By Lemma~\ref{l:layers}(c) the element $\eta$ is a topological generator of~$V$ for the cyclotomic character
of~$\ell$. So Input~2b applies: the module on the right is the kernel of \m{g-1}, where $g$ is the operator by which
$\eta^{-1}$ acts on \m{\dTw{\ell}{\ell}{G}{h}}. By Lemma~\ref{l:twdual}, $g^\vee$ is the operator by which $\eta$
acts on \m{\xTw{\ell}{\ell}{G}{h}}, that is, \m{\cyc{L}{\overline L}{\ell}{\eta}^h\eta}. Now
\m{\cyc{L}{\overline L}{\ell}{\eta}=u^{\ell^m}}, and the scalar $u^h$ commutes with~$\sigma$, so
\[
g^\vee=u^{\ell^mh}\sigma^{\ell^m}=(u^h\sigma)^{\ell^m}.
\]
Since \m{(g-1)^\vee=g^\vee-1}, Lemma~\ref{l:cokdual} gives
\[
\len\cok\bigl((u^h\sigma)^{\ell^m}-1\bigr)=\len\ker(g-1)=\shaInv{L_m}{\ell}{h}.
\]
For \m{h\ge1} the operator \m{u^h\sigma} on \m{\Xinf{\ell}{\ell}{G}} is \m{u^h\gamma}, so the left-hand side is
\m{\tw{\gamma}{u}{h}{m}} by Definition~\ref{d:tau}.
\end{proof}
\leanline{\lean{step0}\sep\lean{descends}\sep\lean{eqTwLen}}

The hypothesis \m{h\ne0} enters through Input~2a only. In the archive a layer is not a subfield of a fixed
algebraic closure: it is any number field over~$L$ that is a layer in the sense of Definition~\ref{d:layer}. The
formal proof begins by embedding such a field into~$\overline L$, and Lemma~\ref{l:layers}(b) shows that the image
is~$L_m$.
% ------------------------------------------------------------------------------------- 9
\sect{Two examples}{s:ex}
Both examples have \m{\ell=3} and \m{u=4}, which has level \m{a=1}; the special values stand in the slots of the
notation, so the twisted length is \m{\tw{\gamma}{4}{h}{m}} and the valuation is \m{\vl{3}{\cdot}}. The modules are free, so the lengths are
valuations of determinants (Lemma~\ref{l:smith}); the build script of this note computes them in exact rational
arithmetic from the matrices and compares them with the formulas below.

\begin{example}[two centres]\label{e:one}
Let \m{X=\Z_3^2} and
\[
\gamma=\begin{pmatrix}4^{-5}&1\\0&-2\end{pmatrix},\qquad \gamma\equiv\begin{pmatrix}1&1\\0&1\end{pmatrix}\pmod 3 ,
\]
so $\gamma$ is unipotent modulo~$3$ and \m{R=2}; the reference layer is \m{\mz{3}{2}=\ExMzero}. The eigenvalues are
$4^{-5}$ and~$-2$, and \m{(-2)^2=4}. The operator \m{(u^h\gamma)^{3^m}-1} is triangular, with diagonal entries
\m{4^{3^m(h-5)}-1} and \m{z-1}, where $z=(-2\cdot4^h)^{3^m}$. Since \m{z\equiv1\pmod 3}, the numbers \m{z-1} and
$z^2-1=4^{3^m(2h+1)}-1$ have the same valuation. By Lemma~\ref{l:lte},
\[
\tw{\gamma}{4}{h}{m}=2\,(m+1)+\vl{3}{5-h}+\vlb{3}{-\tfrac12-h}\qquad(h\ne5),\qquad\tw{\gamma}{4}{5}{m}=\infty .
\]
The centres are \m{c_1=5} and \m{c_2=-\frac12}, and Theorem~A is an equality for \m{m\ge\mz{3}{2}}. The first centre is
an integer: at \m{h=5} the operator \m{(u^h\gamma)^{3^m}-1} is singular and the cokernel is infinite. The second
is not an integer, and the integers \m{1,4,13,40,121,\dots} approach it. For \m{Z=\ExZ} we have
\m{\lfloor\log_\ell Z\rfloor=\ExK}, and the exceptional twists are \ExExc. Every other twist up to $\ExZ$ has
\m{\vl{3}{c_1-h}+\vl{3}{c_2-h}\le\ExWorst}. The formula was confirmed for all \m{1\le h\le\ExZ} and the \ExLevels{} layers
\m{0\le m<\ExLevels}.
\end{example}

\par\vspace{2pt}
\noindent\makebox[\linewidth][c]{\input{figure.tex}}
\par\vspace{-2pt}
{\fontsize{9.2}{12.4}\selectfont\color{inksoft}\noindent\textbf{Figure.} Example~\ref{e:one}: the part of the
length that depends on the twist. Each stem belongs to one centre; a label is the twist. The two stems that pass
the dashed line are the exceptional twists for \m{Z=\ExZ}.\par}

\begin{example}[why a reference layer]\label{e:two}
Let \m{X=\Z_3^3} and let $\gamma$ permute the coordinates cyclically. Then \m{\gamma^3=1}, so
\m{(\gamma-1)^3\equiv0\pmod 3} and $\gamma$ is unipotent modulo~$3$; here \m{R=3}. Since
\m{\det(t\gamma-1)=t^3-1}, Lemma~\ref{l:lte} gives \m{\tw{\gamma}{4}{h}{0}=2+\vl{3}{h}}, while \m{(u^h\gamma)^{3^m}} is the
scalar \m{u^{3^mh}} for \m{m\ge1}, so \m{\tw{\gamma}{4}{h}{m}=3\,(1+m+\vl{3}{h})}.

\par\vspace{3pt}
{\fontsize{9.6}{12.6}\selectfont\noindent\hspace*{2.1em}\begin{tabular}{@{}l@{\hspace{16pt}}*{5}{r@{\hspace{14pt}}}@{}}
layer $m$ & 0 & 1 & 2 & 3 & 4\\\midrule
$\tw{\gamma}{4}{1}{m}$ & \ExTwoRowOne\\
$\tw{\gamma}{4}{9}{m}$ & \ExTwoRowNine
\end{tabular}\par}
\vspace{3pt}

From the layer $1$ on, each step costs exactly \m{R=3}: Theorem~B is sharp. The step from the layer $0$ to the
layer~$1$ costs \ExTwoJump{} at \m{h=9}, more than~$R$. Lemma~\ref{l:step} does not apply there, because
\m{u^h\gamma} is not congruent to $1$ modulo~$9$. A reference layer is needed, and
Proposition~\ref{p:growth} shows that it can be taken to depend on $\ell$ and $R$ alone.
\end{example}

\subhead{Modules with torsion.} The two examples are free modules. As a test of Proposition~\ref{p:growth} where
determinants are not available, the script also draws \RandCases{} random modules, for \m{\ell=2,3,5}, each the sum
of a free module of rank at least~$1$ and a non-zero finite module, with a random endomorphism that is unipotent
modulo~$\ell$ and a random twist, and computes the lengths from invariant factors. In all \RandSteps{} steps between
consecutive layers, starting at the layer \m{\lfloor\log_\ell R\rfloor+2}, the length grew by at most~$R$.

% ------------------------------------------------------------------------------------- 10
\sect{Notes on the formal text}{s:notes}
Reading the Lean text for the mathematics, and not only for its statements, brought out the following points.

\begin{enumerate}[label=\textup{\arabic*.},leftmargin=1.6em,itemsep=5pt]
\item \emph{No logarithm, no structure theorem.} A textbook proof would take the $\ell$-adic logarithm of the
eigenvalues, or the structure theorem for modules over the Iwasawa algebra, and would have to treat eigenvalues
that are not principal units, as well as torsion. The formal proof replaces all of this by the ultrametric
inequality of Lemma~\ref{l:units} and the compactness of~$\Z_\ell$. The centres are not computed and are not
canonical: they depend on the presentation, on the lift~$S$, and on a choice in a non-empty intersection of balls.
\item \emph{The reference layer has one layer to spare.} The formal proof of Proposition~\ref{p:growth} proves
the step for all \m{n\ge\lfloor\log_\ell R\rfloor+2} and applies it from~$\mz{\ell}{R}$ on, and the centres can be taken
at any layer. So Theorems A, B and~D hold with \m{\lfloor\log_\ell R\rfloor+2} in place of~$\mz{\ell}{R}$. This sharper
form is not in the archive; the random test of §\ref{s:ex} starts at that layer. The archive's value is kept in
the statements above.
\item \emph{Unipotence is needed for growth only.} Corollary~\ref{c:level} holds for every endomorphism.
\item \emph{The formal statement is weaker than its proof.} The archive's form of Theorem~D\,(i) asserts only
that there exist a field with an ultrametric norm extending that of~$\Q_\ell$, centres in it, at most $R$ in
number, and an exponent in place of~$a$ for which $L$ contains the roots of unity of that order; this is the tower
bound of Definition~\ref{d:abound}, and it is stated for every layer in the sense of Definition~\ref{d:layer}. Its
proof gives centres in~$\Z_\ell$, exactly~$R$ of them, and the exponent~$a$; this is the form stated in
Theorem~D\,(i). The extra generality is carried through Corollary~\ref{c:exc} and is used nowhere else.
\item \emph{Infinite lengths are part of the statement.} Lengths take values in \m{\{0,1,\dots,\infty\}}, and
Theorem~A holds at a singular twist only because both sides are infinite. Finiteness at every twist is a
consequence of Theorem~D\,(ii) and of the finiteness of the group $H^2$, not of the algebra.
\item \emph{Where the hypotheses on the field enter.} The roots of unity in~$L$ make the image of the cyclotomic
character of \m{\GalG{L}{\overline L}} procyclic and give the level of~$u$. They also make $L$ totally complex
(Lemma~\ref{l:roots}(d)), which Input~3 needs when \m{\ell=2}; the archive lists this as a separate hypothesis of
Theorem~D and proves elsewhere that it is implied (Lemma~\ref{l:rank}). That lemma is not
used by the five final theorems: none of them reaches it in the graph of dependencies of the archive.
\item \emph{The exponent is lifted twice.} Lemma~\ref{l:lte} rests in the archive on a statement about discrete
valuation rings, which has its own proof in a directory on local rings. Lemma~\ref{l:princ}(b) rests there on two
formulas for the multiplicity of~$\ell$ in \m{u^n-1}, one for odd~$\ell$ and one for \m{u\equiv1\pmod{\ell^2}},
which a directory on bounds for class numbers deduces from the corresponding statement of the Lean library Mathlib.
Neither refers to the other. Here both are Lemma~\ref{l:lte}.
\item \emph{What is proved about descent, and where.} The formal proof of Input~2 belongs to the same unit of the
archive as the algebra of this note (Appendix~A), and §\ref{s:descent} follows it. The step that is cohomology,
Input~2a, and Input~2b lie outside this unit. Input~2b is stated for every integer~$h$, Input~2a and Input~2 for \m{h\ne0};
Theorem~D uses \m{h\ge1}.
\end{enumerate}

Some proofs above differ from the formal ones. The archive proves Lemma~\ref{l:cong}(c) by an explicit
computation with elements, where the ideals~$I_n$ are used here, and states Lemma~\ref{l:cong}(a) as an
equivalence: one side is congruent to~$0$ if and only if the other is. Lemma~\ref{l:eigen} is proved there through
the characteristic polynomial, and Lemma~\ref{l:lte} for a general discrete valuation ring. In
Lemma~\ref{l:char}(c) the archive shows that the characters of a simple module are the multiples of one of them,
where the isomorphism with \m{\Z_\ell/\ell\Z_\ell} is used here. It proves Lemma~\ref{l:princ}(c) from a statement
of another directory: a closed subgroup of~$\Z_\ell^\times$ that contains an element of level~$a$
contains~$\princ{\ell}{a}$. The uniqueness of~$L_m$ as a subfield of~$L_\infty$, in the proof of
Lemma~\ref{l:layers}(a), is not in the archive, where a layer is defined as in Definition~\ref{d:layer}. The short
proofs of Lemma~\ref{l:roots}, of Lemma~\ref{l:twdual} and of Corollary~\ref{c:exc} are written for this note from
the definitions; they were not compared with the formal proofs step by step. In all other places the proofs follow
the formal text step by step.

For §\ref{s:descent} the formal proofs were not read in the Lean text itself. They were read in the structured
proofs of the reading project, in which every step quotes the lines of the archive that it renders. A program
compared the quotations with the archive, for all \UnitSteps{} steps of the \UnitProofs{} proofs of the unit, and
found that no line of these proofs is left without a step. The statements of §\ref{s:descent} were compared with
the Lean text, as everywhere in this note.

\subhead{Changes in this edition.} Earlier editions wrote the twisted length as a function of the twist and the
layer alone, and stated the hypotheses of Theorems A to~D in sentences before the statements. That hid the
endomorphism and the scalar for which the centres are chosen, and with them the assumption on the endomorphism.
In this edition every function that the archive defines is written with all the arguments of its formal
definition, in every place, and every theorem cites the Setting that holds its hypotheses.

Beyond the twisted length, the reference layer and the valuation, this changes the arithmetic part. The two
modules of §\ref{s:arith}, written without arguments before, are \m{\Dinf{\ell}{N}{U}} and
\m{\Xinf{\ell}{N}{U}}, with \m{N=\ell} and the whole Galois group in the slots for the modules of Theorem~D. The
exponent that was~$a$ is \m{\cZero{\ell}{L}}. The groups of principal units and their preimages in the Galois group, which carried one index before, are
$\princ{\ell}{n}$ and \m{\Un{L}{\overline L}{\ell}{n}}, the group of a layer is \m{\layU{\ell}{L}{\overline L}{m}}, and
the Galois group is written with the field and its closure. The cyclotomic character carries the prime, the field
and the closure, and so do the Tate twists. The twisted module and its invariants carry all their arguments. The length of the classes is
\m{\shaInv{F}{\ell}{h}}, and the length of the second cohomology group is \m{\VLen{\mathcal D}{h}}, a function of
a local datum, which earlier editions did not name.

The definitions that stood in running text (the valuation, generators, cokernels and congruences in
§\ref{s:results}; the objects of §\ref{s:arith} and §\ref{s:descent}) are now definitions of their own. The three
hypotheses on Galois cohomology, which earlier editions named in a paragraph after Theorem~D, are
Setting~\ref{set:coh}, and Theorem~D and Input~3 assume them in their statements. Theorem~D\,(ii) is stated for
every local datum, and the theorem ends with the form of its conclusion that the archive states
(Definition~\ref{d:abound}). The statements of the last part of §\ref{s:arith}, which earlier editions described in
three paragraphs, are stated with their hypotheses. Lemma~\ref{l:small}(b), Lemma~\ref{l:sep} and
Lemma~\ref{l:princ}(b) are stated in the generality of the archive; Lemma~\ref{l:len}(c), Lemma~\ref{l:step},
Proposition~\ref{p:dual}, Lemma~\ref{l:roots}, Lemma~\ref{l:twdual}, Corollary~\ref{c:unip} and
Proposition~\ref{p:adesc} state, as parts or as statements of their own, what earlier editions said inside proofs or between the statements; and
Corollary~\ref{c:exc} is the form of Corollary~C that the archive proves.

Theorems A and~B, Corollary~C, the bounds of Theorem~D and the inputs say what they said before, and the arguments
of the proofs are the same; some of them moved with the statements they prove.

% ------------------------------------------------------------------------------------- A
\head{Appendix A. Concordance with the Lean archive}
{\fontsize{9.6}{12.8}\selectfont\color{inksoft}\noindent Each name is a link to the lines of the declaration at
commit {\leanfont\Pin} of the public mirror of the archive. The line ranges and the quotations below are read from
the archive by the build script, which also checks them against the SHA-256 of each file. The reading project
divides the part of the archive that the final theorems use into units. This note is the note of one unit: the
\UnitInDir{} declarations of the directory that the final theorems use, and \UnitOutside{} declarations of the neighbouring
directory on towers of number fields, those that prove Input~2. The table contains all \UnitMembers{} of them, and
the build of the note fails if one is missing.\par}
\vspace{4pt}
{\fontsize{8.6}{11.6}\selectfont\renewcommand{\leanfont}{\ttfamily\fontsize{7.7}{10}\selectfont\color{leancol}}
\setlength{\LTleft}{0pt}\setlength{\LTright}{\fill}\setlength{\tabcolsep}{0pt}
\begin{longtable}{@{}p{30mm}p{56mm}p{43mm}r@{}}
\textbf{In this note} & \textbf{Lean declaration} & \textbf{File} & \textbf{Lines}\\\midrule\endhead
\ConcordanceRows
\end{longtable}}

\Needspace{8\baselineskip}
\subhead{Four statements as they are written in the archive.} Theorem~D:
\quoted{aTower}
Its conclusion, the tower bound of Definition~\ref{d:abound}, with the second bound; here {\leanfont e} is the
twist~$h$, {\leanfont s} the centres, {\leanfont c} the exponent~$a'$, which the proof takes to be the level~$a$,
{\leanfont jZero l R} the reference layer~$\mz{\ell}{R}$, {\leanfont shaInv L₁ l e} the length
$\shaInv{L'}{\ell}{h}$ and {\leanfont VLen D e} the length $\VLen{\mathcal D}{h}$:
\quoted{ABound}
\quoted{SecondBound}
Theorem~A, where {\leanfont twLen γ u e m} is $\tw{\gamma}{u}{h}{m}$, with all four arguments in the same order, and
{\leanfont vl l x} is the valuation~$\vl{\ell}{x}$. The endomorphism {\leanfont γ} and the scalar {\leanfont u} are bound, with
their hypotheses, before the centres {\leanfont s}:
\quoted{centres}
Input~2, where {\leanfont γ} is the element~$\sigma$ of \m{G=\GalG{L}{\overline L}}, {\leanfont xTw l l ⊤ e γ} is
the operator $u^h\sigma$ on \m{\Xinf{\ell}{\ell}{G}}, with the prime, \m{N=\ell} and the whole group in its first
three slots, {\leanfont cokerLen} is the length of a cokernel, and {\leanfont IsLayer l m L₀ L₁} says that
{\leanfont L₁} is a layer of index~$m$ over~{\leanfont L₀} for~$\ell$ (Definition~\ref{d:layer}):
\quoted{step0}

% ------------------------------------------------------------------------------------- notation
\head{Notation}
{\fontsize{9.6}{12.6}\selectfont\color{inksoft}\noindent The symbols of the text. A function that the archive
defines is listed with all its arguments; the brackets give its formal name with its arguments, or the letter used
in the archive.\par}
\vspace{4pt}
\begingroup\fontsize{8.6}{10.9}\selectfont\raggedright
\renewcommand{\leanfont}{\ttfamily\fontsize{7.6}{10}\selectfont\color{leancol}}
\setlength{\columnsep}{9mm}\setlength{\parskip}{0pt}
\begin{multicols}{2}
\nt{$\ell$, $\Z_\ell$, $\Q_\ell$, $\overline{\Q}_\ell$}{a prime; the $\ell$-adic integers and numbers; an algebraic closure}
\nt{$\vl{\ell}{x}$, $|z|$}{$\ell$-adic valuation of~$x$, with \m{\vl{\ell}{\ell}=1} [{\leanfont vl l x}]; absolute value $\ell^{-\vl{\ell}{z}}$}
\nt{$X$, $R$, $\pi$}{a $\Z_\ell$-module; a number of generators; a surjection \m{\Z_\ell^R\to X} [{\leanfont GenBy l R X}]}
\nt{$\len M$}{length of a module $M$}
\nt{$\End(X)$, $\cok f$, $\ker f$}{endomorphisms; cokernel \m{X/f(X)} [its length: {\leanfont cokLen f}]; kernel}
\nt{$I_n$, $f\equiv g\pmod{\ell^n}$}{endomorphisms with image in $\ell^nX$; \m{f-g\in I_n} [{\leanfont InPow n (f - g)}]}
\nt{$\gamma$}{an endomorphism of~$X$, an argument of the twisted length}
\nt{$u$, $a$}{an element of $\Z_\ell$, an argument of the twisted length; its level \m{\vl{\ell}{u-1}} [$c$; {\leanfont IsLevel l u c}]}
\nt{$h$, $m$}{twist [$e$]; layer}
\nt{$\mz{\ell}{R}$}{reference layer \m{\lfloor\log_\ell R\rfloor+3} of the prime~$\ell$ and the number~$R$ [{\leanfont jZero l R}]}
\nt{$\tw{\gamma}{u}{h}{m}$}{twisted length of the pair $(\gamma,u)$ at twist~$h$ and layer~$m$ (Definition~\ref{d:tau}) [{\leanfont twLen γ u e m}]}
\nt{$c$, $c_i$}{a centre; the centres [$s_i$]}
\nt{$Z$, $\mathcal E_Z$, $\mathcal E$, $B$}{range of twists; exceptional twists up to~$Z$; a finite set of twists; a bound for \m{\vl{\ell}{c_i-h}}}
\nt{$K$, $a'$, $R'$}{a field with an ultrametric absolute value extending that of~$\Q_\ell$; an exponent and a number of centres (Definition~\ref{d:abound})}
\nt{$f$, $g$, $w$, $\nu$}{endomorphisms; \m{\nu=1+w+\dots+w^{\ell-1}}}
\nt{$A$, $P$, $Q$, $d_i$}{a square matrix; invertible matrices; elementary divisors (Lemma~\ref{l:smith})}
\nt{$S$, $t$, $t'$, $z$, $z_i$}{a matrix lifting an endomorphism; scalars; an element of $\overline{\Q}_\ell$; eigenvalues}
\nt{$n$, $n'$, $i$, $j$, $k$, $r$, $m'$}{integers: exponents, indices, layers}
\nt{$x$, $y$, $b_n$, $r_n$}{elements of a module, of $\Z_\ell$ or of a field; points and radii (Lemma~\ref{l:helly})}
\nt{$T$}{an indeterminate}
\nt{$D$, $D[\ell]$, $\varphi$, $g^\vee$}{a module, its elements killed by~$\ell$, an element; the dual endomorphism}
\nt{$M^\vee$, $\chi$}{the characters \m{M\to\QZl{\ell}} of a module~$M$ (§\ref{s:descent}); a character}
\nt{$W$, $W'$, $W^\perp$}{submodules; the characters that vanish on~$W$}
\nt{$F$, $\overline F$, $F'$, $L$, $\overline L$, $L'$, $L''$}{number fields, algebraic closures of them, extensions}
\nt{$\GalG{F}{\overline F}$, $G$}{the Galois group, with the field and its closure [{\leanfont G F Ω}]; inside a statement, a name for \m{\GalG{L}{\overline L}}}
\nt{$\cyc{F}{\overline F}{\ell}{\sigma}$}{the $\ell$-adic cyclotomic character of \m{\GalG{F}{\overline F}} at~$\sigma$ [{\leanfont cycZ F Ω l σ}]}
\nt{$\sigma$, $\rho$, $\rho'$, $\eta$, $\zeta$, $\zeta_n$}{elements of a Galois group; a root of unity, one of order~$\ell^n$}
\nt{$\princ{\ell}{n}$}{the units of~$\Z_\ell$ congruent to~$1$ modulo~$\ell^n$ [{\leanfont princ l n}]}
\nt{$U$, $V$, $\levU{\ell}{U}{n}$, $\cycU{\ell}{U}$}{subgroups of a Galois group; the elements of~$U$ whose cyclotomic character lies in~$\princ{\ell}{n}$ [{\leanfont levU l U n}], or is~$1$ [{\leanfont cycU l U}]}
\nt{$\cZero{\ell}{L}$}{the largest $n$ with a primitive root of unity of order~$\ell^n$ in~$L$ [{\leanfont cZero l L}]}
\nt{$\layU{\ell}{L}{\overline L}{m}$}{the group of the layer~$m$ (Definition~\ref{d:lay}) [{\leanfont layU l L Ω m}]}
\nt{$L_m$, $L_\infty$}{the layers of the cyclotomic tower of Setting~\ref{set:tower}; their union}
\nt{$\mathcal O$, $\mathfrak p$, $F_{\mathfrak p}$}{a valuation ring of~$\overline F$; a prime of~$F$; the completion}
\nt{$N$, $\Lambda$, $\XDual{N}{U}{\Lambda}{M}$}{an integer, a ring; the characters of~$U$ with values in the $\Lambda$-module~$M$ and local conditions (Definition~\ref{d:dinf}) [{\leanfont XDual N U R M}]}
\nt{$\Dinf{\ell}{N}{U}$, $\Xinf{\ell}{N}{U}$}{the case of $\cycU{\ell}{U}$ and $\QZl{\ell}$ [{\leanfont Dinf l N U}]; its dual [{\leanfont Xinf l N U}], the Iwasawa module for \m{N=\ell} and the whole group}
\nt{$\dTw{\ell}{N}{U}{h}$, $\xTw{\ell}{N}{U}{h}$}{the two modules with the twisted actions [{\leanfont dTw l N U e}, {\leanfont xTw l N U e}]}
\nt{$\ShaD{\ell}{N}{U}{V}{h}$}{the elements of the first of them fixed by~$V$ [{\leanfont ShaD l N U V e}]}
\nt{$\qzG{F}{\overline F}{\ell}{h}$, $\qzTw{\mathcal D}{h}$, $\Z_\ell(h)$}{Tate twists: $\QZl{\ell}$ with the action of \m{\GalG{F}{\overline F}} through the $(-h)$-th power of the cyclotomic character [{\leanfont qzG F Ω l e}]; the same for the group of a local datum [{\leanfont qzTw D e}]; $\Z_\ell$ with the $h$-th power}
\nt{$\Sha^1_\ell(F,h)$, $\shaInv{F}{\ell}{h}$}{classes trivial above~$\ell$ and unramified elsewhere (Definition~\ref{d:sha}); the length of this group [{\leanfont shaInv F l e}]}
\nt{$\mathcal D$, $\Gamma_{\mathcal D}$, $H^i(\mathcal D,\cdot)$}{a local datum of~$F$ at~$\ell$ [{\leanfont D : LData F Ω l}]; the Galois group unramified outside~$\ell$ and infinity, formed in the closure of the datum; its continuous cohomology}
\nt{$\ShaOf{\mathcal D}{h}$, $\VLen{\mathcal D}{h}$}{the classes in the description by a local datum [{\leanfont Sha D e}]; the length of \m{H^2(\mathcal D,\Z_\ell(h+1))} [{\leanfont VLen D e}]}
\nt{$\mu$, $R_\ell$}{Iwasawa's invariant (§\ref{s:arith}); number of generators for~$\ell$}
\nt{$q$, $\Delta$, $p$, $p'$}{a modulus, a subgroup of \m{(\Z/q\Z)^\times}; two primes (Proposition~\ref{p:fixed})}
\nt{$C$, $d_L$, $[L:\Q]$, $\#$}{a constant; discriminant and degree of~$L$; number of elements (Definition~\ref{d:K})}
\nt{$\theta$, $\Theta$, $\Xi$, $J$, $Y$, $\omega$}{the real numbers of Lemma~\ref{l:window}}
\nt{$\Phi(T)$, $s$}{characteristic polynomial; an $\ell$-adic variable (Remark~\ref{r:classical})}
\nt{$E$}{in §\ref{s:intro}, the number field of the guide}
\end{multicols}
\endgroup
% ------------------------------------------------------------------------------------- references
\head{References}
\begingroup\fontsize{9.4}{12.6}\selectfont\raggedright\renewcommand{\section}[2]{}
\begin{thebibliography}{NSW}\setlength{\itemsep}{2pt}
\bibitem[Al]{Al} L.\,Alpöge, announcement of proofs of the Riemann hypothesis and of the generalized Riemann
  hypothesis for Dirichlet characters, 20 November 2026, with a Lean archive;
  \href{\PinUrl}{public mirror of the archive at commit {\leanfont\Pin}}.
\bibitem[At]{At} \href{\AtlasUrl}{Riemann proof atlas}: the declarations of the archive on a map of arXiv papers;
  data at commit {\leanfont\AtlasCommit}.
\bibitem[FW]{FW} B.\,Ferrero and L.\,C.\,Washington, The Iwasawa invariant $\mu_p$ vanishes for abelian number
  fields, \emph{Ann. of Math.} (2) \textbf{109} (1979), 377--395.
\bibitem[Gu]{Gu} \emph{The proof of the Riemann hypothesis announced by L.\,Alpöge: a guide to its architecture},
  riemann-rev, \texttt{guide/proof-guide.pdf}.
\bibitem[Iw]{Iw} K.\,Iwasawa, On $\Z_\ell$-extensions of algebraic number fields, \emph{Ann. of Math.} (2)
  \textbf{98} (1973), 246--326.
\bibitem[MW]{MW} B.\,Mazur and A.\,Wiles, Class fields of abelian extensions of~$\Q$, \emph{Invent. Math.}
  \textbf{76} (1984), 179--330.
\bibitem[NSW]{NSW} J.\,Neukirch, A.\,Schmidt and K.\,Wingberg, \emph{Cohomology of Number Fields}, 2nd ed.,
  Grundlehren der mathematischen Wissenschaften 323, Springer, 2008.
\bibitem[Wa]{Wa} L.\,C.\,Washington, \emph{Introduction to Cyclotomic Fields}, 2nd ed., Graduate Texts in
  Mathematics 83, Springer, 1997.
\end{thebibliography}
\endgroup

\end{document}
